% concatenated from main.tex, Headers_aut.tex, 1-Intro.tex, 2-Preliminaires.tex, 3-6_PartI_Isom.tex, 7-9_PartII_Autom.tex, 10-11_Part_III.tex
\documentclass[a4paper,english,11pt]{article}
\pdfoutput=1

\usepackage{amsmath,amsfonts,amssymb,amsthm}
\usepackage{mathabx} 
\usepackage{blindtext}
\usepackage{titlesec}
\usepackage{titletoc}

\usepackage[utf8]{inputenc} 
% ==================================
% ==================================
%  Font parametrization
% =================================
% ==================================
% To use Computer Modern (Latex default font)
% comment the lines
% “\usepackage{sansmathfonts}” and
% “\usepackage[sfdefault]{atkinson}” and
% “\usepackage[default]{opensans}” below.
\usepackage{sansmathfonts}
% ----------------------------------------------
% Atkinson hyperlegible font (see https://ctan.org/tex-archive/fonts/atkinson/)
% To use this font, comment the line
%“\usepackage[default]{opensans}” below and 
% uncomment the following one
\usepackage[sfdefault]{atkinson}
\usepackage[T1]{fontenc}
% ----------------------------------------------
% OpenSans
% To use this font, comment the line
% “\usepackage[sfdefault]{atkinson}” above and 
% uncomment the following one
%\usepackage[default]{opensans}
% ==================================
% ==================================
\usepackage[colorlinks=true,
    citecolor=DeepSkyBlue4,
    linkcolor=DeepSkyBlue4,
    urlcolor=DeepSkyBlue4,
    linktoc=page,
    hyperindex=true,
    pdfcreator={}]{hyperref} 
\hypersetup{
        pdfauthor={Amandine Escalier and Camille Horbez},
        pdftitle={Isomorphisms and automorphisms of graph products of groups},
        pdfkeywords={Graph products of groups, Isomorphism problem, Automorphism group},
        %MSC_classes={20F28,05C25},
}
\usepackage[english]{babel} 
\usepackage[x11names]{xcolor}
\definecolor{marronfonce}{cmyk}{0,0.06,0.20,0.6}
\definecolor{marron}{cmyk}{0,0.03,0.08,0.26} 
\colorlet{marronclair}{marronfonce!25}
% -----------------------------------------------
\setcounter{tocdepth}{1}
% -- Page style
\usepackage{geometry}
\geometry{vmargin=2.25cm}
%
\usepackage{setspace} %Interline space
\linespread{1.06}
%
\usepackage{enumitem}
\setlist{nosep}
% --
\usepackage{graphics}
% -----------------------------------------------
\newtheorem{de}{Definition}[section]
\newtheorem{theo}[de]{Theorem} 
\newtheorem{prop}[de]{Proposition}
\newtheorem{lemma}[de]{Lemma}
\newtheorem{cor}[de]{Corollary}
\newtheorem{claim}[de]{Claim}
\newtheorem{Claim}[de]{Claim}
%
\theoremstyle{definition}
\newtheorem{Fact}[de]{Fact}
\theoremstyle{remark}
\newtheorem{rk}[de]{Remark}
\newtheorem{Rq}[de]{Remark}
\newtheorem{ex}[de]{Example}
\newtheorem{setting}[de]{Setting}
%
\normalsize
%	
\newcommand{\st}{\mathrm{st}}
\newcommand{\lk}{\mathrm{lk}}
\newcommand{\Out}{\mathrm{Out}}
\newcommand{\Aut}{\mathrm{Aut}}
\newcommand{\Inn}{\mathrm{Inn}}
\newcommand{\PConj}{\mathrm{PConj}}
\newcommand{\calf}{\mathcal{F}}
\newcommand{\calF}{\mathcal{F}}
\newcommand{\bZ}{\mathbb{Z}}
\newcommand{\twistor}{\xi}
% ====================================================
% Code for \llangle \\rangle without using the Mnsymbol package
% ====================================================
\makeatletter
\DeclareFontFamily{OMX}{MnSymbolE}{}
\DeclareSymbolFont{MnLargeSymbols}{OMX}{MnSymbolE}{m}{n}
\SetSymbolFont{MnLargeSymbols}{bold}{OMX}{MnSymbolE}{b}{n}
\DeclareFontShape{OMX}{MnSymbolE}{m}{n}{
    <-6>  MnSymbolE5
   <6-7>  MnSymbolE6
   <7-8>  MnSymbolE7
   <8-9>  MnSymbolE8
   <9-10> MnSymbolE9
  <10-12> MnSymbolE10
  <12->   MnSymbolE12
}{}
\DeclareFontShape{OMX}{MnSymbolE}{b}{n}{
    <-6>  MnSymbolE-Bold5
   <6-7>  MnSymbolE-Bold6
   <7-8>  MnSymbolE-Bold7
   <8-9>  MnSymbolE-Bold8
   <9-10> MnSymbolE-Bold9
  <10-12> MnSymbolE-Bold10
  <12->   MnSymbolE-Bold12
}{}
%
\let\llangle\@undefined
\let\rrangle\@undefined
\DeclareMathDelimiter{\llangle}{\mathopen}%
                     {MnLargeSymbols}{'164}{MnLargeSymbols}{'164}
\DeclareMathDelimiter{\rrangle}{\mathclose}%
                     {MnLargeSymbols}{'171}{MnLargeSymbols}{'171}
\makeatother
%% ==========================================
%% =============================================
% ----------------
\newcommand{\deffont}[1]{\textbf{#1}}
% -------- 
\usepackage{wasysym}
\renewcommand{\qedsymbol}{\pentagon}
% ----------------- COLORBOX
\usepackage{tcolorbox}
\newtcolorbox{SettingBox}{colframe=black!75,colback=white}

	


\title{Isomorphisms and automorphisms of graph products of groups}
\author{Amandine Escalier and Camille Horbez}
\date{\today}



\begin{document}
\maketitle
\begin{abstract}
We solve the isomorphism problem for graph products of groups. We give a generating set for the automorphism group of a graph product of groups, generalizing the one given by Laurence and Servatius for right-angled Artin groups. 
\end{abstract}

%\vfill

~\\
\setcounter{tocdepth}{1}
{\small \startcontents  \printcontents{}{0}[1]{}} \vfill
\noindent \textcolor{black!70}{\small This work is openly licensed via \href{https://creativecommons.org/licenses/by/4.0/}{Creative Commons CC-BY 4.0}}
\smallskip

\noindent \textcolor{black!70}{\small This work was completed without use of artificial intelligence. The authors do not consent for all or part of this work to be processed by artificial intelligence models.}
\newpage

% ======================================
% SECTION 1 INTRODUCTION 
% ======================================
\section{Introduction}
\subsection{Statement of results}

Graph products of groups, introduced by Green in \cite{Gre}, provide a general construction that encompasses both direct products and free products. Starting with a finite simple graph $\Gamma$, and a family of {non-trivial} groups $(G_v)_{v\in V\Gamma}$ indexed by the vertex set of $\Gamma$, the associated graph product is the group generated by the groups $G_v$, where the only extra relations impose that $G_v$ and $G_w$ commute whenever $v$ and $w$ are adjacent. If the graph $\Gamma$ is complete, then one recovers the direct product of the groups $G_v$, and if $\Gamma$ is edgeless, then one obtains their free product. Graph products encompass right-angled Artin groups (when all vertex groups are isomorphic to $\mathbb{Z}$) and right-angled Coxeter groups (when they are isomorphic to $\mathbb{Z}/2\mathbb{Z}$).

In the present paper, we resolve the following two problems for graph products.

\medskip

\begin{itemize}
\item \textbf{(Isomorphism problem)}. If the same group splits as a graph product in two ways, over $\Gamma$ with vertex groups $(G_v)_{v\in V\Gamma}$, and over $\Lambda$ with vertex groups $(H_w)_{w\in V\Lambda}$, under what conditions does this force the existence of an isomorphism $\theta:\Gamma\to\Lambda$ such that $H_{\theta(v)}$ is isomorphic to $G_v$ for every $v\in V\Gamma$?
\item \textbf{(Generators for $\boldsymbol{\Aut(G)}$)}. Give a natural set of generators for the automorphism group of a graph product.
\end{itemize}

\medskip

The isomorphism problem, first raised in \cite[Chapter~4]{Gre}, was only solved in full generality under very restrictive assumptions on the vertex groups, e.g.\ for right-angled Artin groups \cite{Dro}, or when the vertex groups are finite \cite{Gre,Rad} or abelian \cite{GP}. In another direction, solutions to the isomorphism problem with no assumptions on the vertex groups, but with restrictions on the graphs, were provided in \cite{Gen,EH}. {See also \cite[Section~5.5]{Ber} and references therein for a discussion of this problem.}

An explicit set of generators of $\Aut(G)$ was given in the case when $G$ is a right-angled Artin group by Laurence \cite{Lau}, confirming a conjecture of Servatius \cite{Ser}. This was later extended to all graph products of {finitely generated} abelian groups in \cite{CG}. Automorphisms of right-angled Coxeter groups were first studied by Tits \cite{TitsRACG} and a generating set was given by Mühlherr \cite{Muhlherr}. More recent developments for general graph products include \cite{GM,Gen} and will be reviewed below. The question of describing $\Aut(G)$ in full generality was formulated by Genevois--Martin in \cite[]{GM}.

\paragraph*{The isomorphism problem.}

Two caveats should be taken into account when solving the isomorphism problem for graph products. The first is that if a vertex group $G_v$ is itself a graph product over a graph $\Gamma_v$, then one gets a new graph product structure for $G$ over a graph $\Lambda$ obtained by replacing the vertex $v$ by $\Gamma_v$, and joining every vertex of $\Gamma_v$ to every vertex $w\in V\Gamma$ that is adjacent to $v$. 
For this reason, one wants to either consider maximal decompositions as graph products, where the vertex groups cannot split further as graph products, or minimal decompositions. 
To make this formal, we import the following notions from \cite{EH}. A full subgraph $\Lambda\subseteq\Gamma$ is \deffont{collapsible} if for any two vertices $x,y\in V\Lambda$, we have $\lk_\Gamma(x)\cap (\Gamma\setminus\Lambda)=\lk_\Gamma(y)\cap (\Gamma\setminus\Lambda)$, where $\lk_\Gamma(x)$ denotes the full subgraph spanned by the vertices of $\Gamma$ that are adjacent to $x$. For instance, the graph $\Gamma_v$ described above is a collapsible subgraph of $\Lambda$. A finite simple graph $\Gamma$ is 
\begin{itemize}
    \item \deffont{strongly reduced} if it does not contain any proper collapsible subgraph on at least two vertices;
    \item \deffont{clique-reduced} if it does not contain any collapsible subgraph on at least two vertices which is a clique.
\end{itemize} 
The second caveat is that a group can split as a direct product in different ways even if the factors do not further split as direct products. 
A concrete example of this phenomenon is given in \cite[Chapter XI, pages 81–83]{Kurosh}, so extra care is required when treating direct products. As in \cite{Gen}, we say that a group $H$ is \deffont{graphically irreducible} if whenever $H$ splits as a graph product over a graph $\Lambda$, the graph $\Lambda$ is complete.

The following theorem solves the isomorphism problem for graph products, taking these two difficulties into account.

\begin{theo}[Isomorphisms between graph products, see Theorems~\ref{theo:iso-strongly-reduced} and~\ref{theo:isomorphism-complete}]
\label{theo:isomorphism}
Let $\Gamma,\Lambda$ be finite simple graphs with at least three vertices. Let $G$ be a graph product over $\Gamma$ with non-trivial vertex groups $(G_v)_{v\in V\Gamma}$, and let $H$ be a graph product over $\Lambda$ with non-trivial vertex groups $(H_w)_{w\in V\Lambda}$. Assume that one of the following two holds:
\begin{itemize}
\item $\Gamma$ and $\Lambda$ are strongly reduced;
    \item $\Gamma$ and $\Lambda$ are clique-reduced, and the groups $G_v,H_w$ are graphically irreducible.
\end{itemize}
Then the following are equivalent.
\begin{enumerate}
    \item The groups $G$ and $H$ are isomorphic.
    \item There exists a graph isomorphism $\theta:\Gamma\to\Lambda$ such that for every $v\in V\Gamma$, the groups $G_v$ and $H_{\theta(v)}$ are isomorphic.
\end{enumerate}
\end{theo}

The two sets of hypotheses amount to looking at minimal (non-trivial) decompositions as a graph product for the first, and maximal decompositions for the second. In the second set of assumptions, the combination of the clique-reduced assumption on the graph, and the graphical irreducibility of vertex groups is here to tackle the problem coming from multiple decompositions of groups as direct products. It is possible to consider intermediate sets of assumptions, where one only forbids a certain family $\mathbb{T}$ of collapsible subgraphs in $\Gamma,\Lambda$ (containing all cliques), and imposes in return that vertex groups can only split as graph products over a graph in $\mathbb{T}$. We refer to Theorem~\ref{theo:isomorphism-complete} on page~\pageref{theo:isomorphism-complete} for the complete statement.

We also mention that the theorem remains true if one only assumes that every connected component of $\Gamma,\Lambda$ is strongly reduced with at least three vertices (and likewise with join components instead of connected components). We refer to Theorem~\ref{theo:iso-strongly-reduced} on page~\pageref{theo:iso-strongly-reduced} for this extension.

\begin{rk}\label{Rq:at-least-three-vertices}
The assumption that $\Gamma,\Lambda$ have at least three vertices cannot be dropped. Indeed, graphs having two vertices are automatically strongly reduced. But the isomorphism $(A\times B)\times C\simeq A\times (B\times C)$ gives two decompositions of a direct product as a graph product over an edge. And writing $(A\ast B)\ast C\simeq A\ast (B\ast C)$ gives two decompositions of a free product as a graph product over an edgeless graph with two vertices.
\end{rk}

\paragraph{Generators for $\Aut(G)$.}

We now introduce our basic automorphisms that will serve as generators for $\Aut(G)$. These were already considered by Servatius \cite{Ser} and Laurence \cite{Lau} in the case of a right-angled Artin group $A_\Gamma$; they proved that they generate $\Aut(A_\Gamma)$. In \cite{GM}, Genevois and Martin considered local automorphisms and partial conjugations in the general context of graph products, and proved that they generate a well-identified subgroup $\PConj(G)\subseteq \Aut(G)$. This is the subgroup consisting of \deffont{conjugating automorphisms}, i.e.\ automorphisms $\varphi$ that send every vertex group $G_v$ to a conjugate of a vertex group $G_w$. 

In what follows we denote by $Z(G)$ the center of a group $G$. Recall that the \deffont{link} $\lk_\Gamma(v)$ of a vertex $v\in V\Gamma$ is the full subgraph spanned by all vertices adjacent to $v$; the \deffont{star} $\st_\Gamma(v)$ is the full subgraph spanned by $v$ and $\lk_\Gamma(v)$. Given a full subgraph $\Lambda\subseteq\Gamma$, we denote by $G_\Lambda$ the subgroup of $G_\Gamma$ generated by all vertex groups $G_v$ with $v\in V\Lambda$; it turns out to be isomorphic to the graph product over $\Lambda$ with vertex groups $(G_v)_{v\in V\Lambda}$.

\begin{de}[Basic automorphisms]\label{de:basic}
Let $G$ be a graph product over a finite simple graph $\Gamma$, with vertex groups $(G_v)_{v\in V\Gamma}$. An automorphism $\varphi\in\Aut(G)$ is
\begin{itemize}
    \item a \deffont{local automorphism} if there exists a graph automorphism $\sigma\in\Aut(\Gamma)$ such that for every $v\in V\Gamma$, one has $\varphi(G_v)=G_{\sigma(v)}$; 
    \item a \deffont{partial conjugation} if there exist a vertex $v\in V\Gamma$, a connected component $C$ of $\Gamma\setminus\st_\Gamma(v)$, and $g\in G_v$, such that $\varphi(h)=ghg^{-1}$ for every $h\in G_C$, and $\varphi(h)=h$ for every $h\in G_{\Gamma\setminus C}$;  
    \item a \deffont{transvection} if it has one of the following two forms:
    \begin{itemize}
        \item $\varphi$ is a \deffont{twist} if there exist distinct vertices $v,v'\in V\Gamma$ such that $\st_\Gamma(v)\subseteq\st_\Gamma(v')$, and a homomorphism $\twistor:G_v\to Z(G_{v'})$, such that $\varphi(g)=g\twistor(g)$ for every $g\in G_v$, and $\varphi(h)=h$ for every $h\in G_{\Gamma\setminus\{v\}}$;
        \item $\varphi$ is a \deffont{fold} if there exist distinct vertices $v,v'\in V\Gamma$ such that $\lk_\Gamma(v)\subseteq\lk_\Gamma(v')$, with $G_v$ isomorphic to $\mathbb{Z}$ (generated by $z_v$), and $g'\in G_{v'}$, such that $\varphi(z_v)=z_vg'$ and $\varphi(h)=h$ for every $h\in G_{\Gamma\setminus\{v\}}$. 
    \end{itemize}
\end{itemize}
\end{de}

Note that every inner automorphism is a product of local automorphisms and partial conjugations.

These automorphisms depend not only on $G$, but on the chosen graph product structure over $\Gamma$. We will talk about \deffont{$\Gamma$-local automorphisms}, \deffont{$\Gamma$-partial conjugations} and \deffont{$\Gamma$-transvections} (including \deffont{$\Gamma$-twists} and \deffont{$\Gamma$-folds}) when we need to emphasize the graph product structure.

\begin{rk}
The terminology \emph{twists} and \emph{folds} comes from the usual terminology for automorphisms of right-angled Artin groups. In this case, if $v,v'$ are adjacent, then $G_{\{v,v'\}}\simeq\mathbb{Z}^2$, and after identifying $\mathbb{Z}^2$ with the fundamental group of a $2$-torus, a twist is nothing but the automorphism realized geometrically by a Dehn twist of the torus. If $v,v'$ are not adjacent, then $G_{\{v,v'\}}\simeq F_2$, and after identifying $F_2$ with a bouquet of two circles, a fold is nothing but the automorphism realized by a homotopy equivalence that folds one circle over the other.  
\end{rk}

In order to state our second main theorem, we need the following definition. A finite simple graph is \deffont{almost strongly reduced} if every proper collapsible subgraph on at least two vertices is edgeless.


\begin{theo}[Generators for $\Aut(G)$, see Theorems~\ref{theo:conclusion} and~\ref{theo:aut-complete}]%
\label{theo:main-aut}
Let $G$ be a graph product over a finite simple graph $\Gamma$ which is not an edge, with non-trivial vertex groups $(G_v)_{v\in V\Gamma}$. Assume that one of the following two holds:
\begin{itemize}
\item $\Gamma$ is almost strongly reduced and for every $v\in V\Gamma$, the group $G_v$ does not split non-trivially as a free product;
\item $\Gamma$ is clique-reduced and for every $v\in V\Gamma$, the group $G_v$ is graphically irreducible.
\end{itemize}
Then $\Aut(G)$ is generated by local automorphisms, partial conjugations and transvections.
\end{theo}

As for the isomorphism problem, a more general version with intermediate sets of assumptions is provided in Theorem~\ref{theo:aut-complete}.

\begin{rk}\label{rk:aut}
The first set of assumptions cannot be replaced by the sole requirement that $\Gamma$ is strongly reduced. Indeed, $\Gamma$ could then contain two vertices $v,w$ with $\lk_\Gamma(v)\subseteq\lk_\Gamma(w)$, respectively carrying vertex groups $G_v=\mathbb{Z}\ast H_v$ (for some group $H_v$) and $G_w$. Then, denoting by $z_v$ a generator of the free factor of $G_v$ isomorphic to $\mathbb{Z}$, and taking a non-trivial element $g\in G_w$, the map sending $z_v$ to $z_vg$, sending every $h\in H_v$ to itself, and sending every $g'\in G_{\Gamma\setminus\{v\}}$ to itself, determines an automorphism of $G$ which is not a $\Gamma$-transvection. And in fact, it does not belong to the subgroup generated by our basic automorphisms for the graph product structure over $\Gamma$. Indeed, one can check that any automorphism in this subgroup sends every commutator contained in a vertex group, in a conjugate of a vertex group. Likewise, any conjugation of $\mathbb{Z}$ or of $H_v$ by a non-trivial element of $G_w$ yields an automorphism of $G$ which is not a $\Gamma$-partial conjugation for the graph product structure over $\Gamma$, and for the same reason as above it does not belong to the subgroup generated by basic $\Gamma$-automorphisms.
\end{rk}

\subsection{Warm-up and strategy of proof}

As a way of presenting our strategy for Theorems~\ref{theo:isomorphism} and~\ref{theo:main-aut}, we will give a proof of the following proposition, which treats a special case of Theorem~\ref{theo:main-aut} when there are no transvections. We recall from the {paragraphs} above Definition~\ref{de:basic} that $\PConj(G)$ is the subgroup of $\Aut(G)$ consisting of conjugating automorphisms, sending every vertex group to a conjugate of a vertex group.

\begin{prop}\label{prop:trivial-center}
Let $\Gamma$ be a finite simple graph, let $(G_v)_{v\in V\Gamma}$ be a family of {non-trivial} groups, and let $G$ be the corresponding graph product. Assume that for every $v\in V\Gamma$, the group $G_v$ does not split as a graph product over a graph with at least two vertices, and $G_v$ has trivial center.

Then $\Aut(G)=\mathrm{PConj}(G)$.
\end{prop}

We refer the reader to Section~\ref{sec:background-graph-product} for the terminology regarding parabolic subgroups of graph products used in the following proof.

\begin{proof}
Let $\varphi\in\Aut(G)$ and $v\in V\Gamma$; we aim to prove that $\varphi(G_v)$ is conjugate to a vertex group. Consider a chain of subgroups \[G_{v}=A_n\subseteq F_{n-1}\subseteq P_{n-1}\subseteq A_{n-1}\subseteq\dots\subseteq F_1\subseteq P_1\subseteq A_1\subseteq F_0=G,\]
where for every $\ell\in\{1,\dots,n\}$,
    \begin{itemize}
        \item $A_\ell$ is a parabolic subgroup of $F_{\ell-1}$ whose type is a connected component of the type of $F_{\ell-1}$, not equal to $G_v$ unless $\ell=n$ (in particular $A_\ell$ is a freely indecomposable free factor of $F_{\ell-1}$);
        \item If $\ell\neq n$ then $P_\ell$ is a maximal parabolic subgroup of $A_\ell$ that splits non-trivially as a direct product and contains $G_v$: this exists because
        \begin{itemize}
        \item the type of $A_\ell$ is not reduced to $v$ as $\ell\neq n$, so $v$ belongs to an edge $e=vw$ with $G_e\subseteq A_\ell$;
        \item the group $G_e$ is a direct product $G_e=G_v\times G_w$;
        \item since any increasing chain of parabolic subgroups terminates, there exists a parabolic subgroups splitting as a direct product and containing $G_v$ which is maximal for inclusion. 
        \end{itemize}
        \item If $\ell\neq n$ then $F_\ell$ is a parabolic subgroup of $P_\ell$ whose type is a join component of the type of $P_\ell$ (in particular $F_\ell$ is a directly indecomposable direct factor of $P_\ell$).
        \end{itemize}
        
We will prove by induction on $\ell\in \{1,\dots,n-1\}$ that $\varphi(A_\ell),\varphi(P_\ell)$ and $\varphi(F_\ell)$ are parabolic subgroups of $H$ and $\varphi(A_\ell)$ is not conjugate to a vertex group. This is true when $\ell=0$ (for $\varphi(F_0)$). Assuming that this is true for $\ell-1$, let us prove it for $\ell$. 
\begin{itemize}
\item We start with $\varphi(A_\ell)$. Consider the free product decomposition \[\varphi(F_{\ell-1})=Q_1\ast\dots\ast Q_k\]
where the subgroups $Q_i$ are the parabolic subgroups associated to the connected components of the type of $\varphi(F_{\ell-1})$. Since vertex groups do not split as free products and are not isomorphic to $\mathbb{Z}$, this is a Grushko decomposition of $\varphi(F_{\ell-1})$, and $\varphi(A_\ell)$ is conjugate to one of the subgroups $Q_i$. Since $\ell\in\{1,\dots,n-1\}$, the parabolic subgroup $A_\ell$ is not conjugate to a vertex group, so it splits non-trivially as a graph product. Therefore, so does $\varphi(A_\ell)$, and it follows from our assumption on the vertex groups that $\varphi(A_\ell)$ is not conjugate to a vertex group.
\item For $\varphi(P_\ell)$, since maximal direct product subgroups of a graph product over a finite connected graph with at least two vertices are precisely the maximal direct product parabolic subgroups (see Lemma~\ref{lemma:maximal-direct-product} below, due to Minasyan--Osin and Genevois), $\varphi(P_\ell)$ is a parabolic subgroup.
\item We finally consider $\varphi(F_\ell)$. Consider the decomposition $\varphi(P_\ell)=K_1\times\dots\times K_m$, where each $K_i$ is a parabolic subgroup whose type does not split as a join. Our assumption on vertex groups ensures that the factors $K_i$ have trivial center and do not split as direct products -- and so do the parabolic direct factors of the analogous decomposition of $P_\ell$. But when a group splits as a direct product $G=G_1\times\dots\times G_k$ and the factors $G_i$ have trivial center and do not split as direct products, the decomposition is unique up to permutation of the factors (see e.g.\ \cite[Lemma~3.5]{Fio}). So $\varphi(F_\ell)$ is equal to one of the factors $K_i$.
\end{itemize}
When $\ell=n$, the same argument as in the first bullet point ensures that $\varphi(A_n)$ is a parabolic subgroup of $H$. Since $A_n=G_v$ is not decomposable as a graph product, necessarily $\varphi(A_n)$ is conjugate to a vertex group, as desired.
\end{proof}

In the general context of Theorems~\ref{theo:isomorphism} and~\ref{theo:main-aut}, it is no longer true that an isomorphism between two free products, or between two direct products, sends factors to factors up to conjugation -- this is where the assumptions on the vertex groups made in Proposition~\ref{prop:trivial-center} led to a significant simplification. But the idea is that, up to precomposing $\varphi$ by an automorphism of the free product or of the direct product, we can assume that they do. The automorphisms we will need in this process are precisely those given in Definition~\ref{de:basic}. There are many technical subtleties to be taken care of -- notably, we need to make sure that the automorphisms we use to precompose actually extend to automorphisms of the whole graph product. This is why we need a careful analysis of free products (Sections~\ref{sec:free-products} and~\ref{sec:free-products-2}) and direct products (Section~\ref{sec:direct-products}). To carry them out successfully, several arguments
arising from the original work of Laurence \cite{Lau} are particularly useful, notably in Section~\ref{sec:free-products-2} regarding free products.

\subsection{Structure of the article}
After a preliminary section recording some standard tools on graph products and free products (Section~\ref{sec:prelims}), the article is split into three parts. Part~\ref{part:isomorphism} solves the isomorphism problem for graph products (Theorem~\ref{theo:isomorphism}) assuming that the graphs are strongly reduced. Part~\ref{part:automorphisms} gives a generating set of their automorphism group (Theorem~\ref{theo:main-aut}), assuming that the graphs are almost strongly reduced. Part~\ref{part:3} generalizes the previous theorems to graphs that are not strongly reduced (or almost strongly reduced), and completes in particular the proofs of all theorems from the introduction.

In Part~\ref{part:isomorphism}, after a short section giving the general strategy (Section~\ref{sec:blueprint}), the next two sections provide the necessary tools to analyze free products (Section~\ref{sec:free-products}) and direct products (Section~\ref{sec:direct-products}). Our proof of Theorem~\ref{theo:isomorphism}, in the case where the graphs are assumed strongly reduced, is completed in Section~\ref{sec:endgame-isomorphism}.

Part~\ref{part:automorphisms} starts in Section~\ref{sec:permutationless} by exploiting our solution to the isomorphism problem to show that every automorphism can be precomposed by a local automorphism in order to ensure that for every vertex $v$, the parabolic support of $\varphi(G_v)$ contains a conjugate of $G_v$. We then develop on the necessary tools to study free products (Section~\ref{sec:free-products-2}), before completing our proof of Theorem~\ref{theo:main-aut}, in the almost strongly reduced case, in Section~\ref{sec:endgame-automorphisms}.

In Part~\ref{part:3}, we extend our solution to the isomorphism problem (Section~\ref{sec:3.1}), and our description of a generating set for $\Aut(G)$ to graphs that are not strongly reduced (Section~\ref{sec:3.2}).

\paragraph{Acknowledgments.} We thank Anthony Genevois for pointing us to
\cite{MO}, which enabled us to correct a mistake in an earlier version of the
paper.



%
% ======================================
% SECTION 2 
% ======================================
\section{Preliminaries}\label{sec:prelims}
\subsection{Background on graph products}\label{sec:background-graph-product}

\paragraph*{General background.}

Let $\Gamma$ be a finite simple graph, i.e.\ $\Gamma$ has no loop-edge and no multiple edges. We will denote by $V\Gamma$ and $E\Gamma$ the vertex and edge set of $\Gamma$, respectively. Let $(G_v)_{v\in V\Gamma}$ be a family of non-trivial vertex groups,
and let $G$ be the graph product over $\Gamma$ with vertex groups $(G_v)_{v\in V\Gamma}$, i.e.\ 
\begin{equation*}
    G:= \left(\bigast_{v\in V\Gamma}G_v\right) \bigg/ \left\llangle [g_v,g_w] \mid  g_v\in G_v,\, g_w\in G_w,\, (v,w)\in E\Gamma \right\rrangle.
\end{equation*}
We insist that by convention, all graph products in this paper have non-trivial vertex groups. A subgraph $\Delta\subseteq \Gamma$ is \deffont{full} if two vertices of $\Delta$ are adjacent in $\Delta$ if and only if they are adjacent in $\Gamma$. Given a full subgraph $\Delta\subseteq\Gamma$, the subgroup $G_\Delta\subseteq G_\Gamma$ generated by the vertex groups $G_v$ with $v\in V\Delta$ is isomorphic to the graph product over $\Delta$ with vertex groups $(G_v)_{v\in V\Delta}$. We call $G_\Delta$ a \deffont{standard parabolic subgroup} of $G_\Gamma$. A \deffont{parabolic subgroup} of $G_\Gamma$ is a subgroup that is conjugate to a standard parabolic subgroup. If $P$ is a parabolic subgroup of $G$, then the subgraph $\Delta$ of $\Gamma$ such that $P$ is conjugate to $G_\Delta$ is unique, see e.g.\ \cite[Corollary~3.8]{AM}; it is called the \deffont{type} of $P$. We insist that the notion of being a parabolic subgroup does not only depend on $G$, but really on the choice of a presentation of $G$ as a graph product. 
Unless otherwise specified, in the sequel of the paper we will always assume a presentation of $G$ as a graph product chosen and fixed. When we need to insist on the graph product structure we consider, we will talk about \deffont{$\Gamma$-parabolic subgroups}.

Given a full subgraph $\Delta\subseteq \Gamma$, there is a \deffont{parabolic
  retraction} $r_{\Delta}:G\to G_{\Delta}$, defined as the identity on
$G_\Delta$ and the trivial hmomorphism on $G_{\Gamma\setminus\Delta}$.
Given a vertex $v\in V\Gamma$, we will use the notation $r_v$ as a shortening for $r_{\{v\}}$.

Given a full subgraph $\Delta\subseteq\Gamma$, we denote by $\Delta^\perp$ the full subgraph whose vertices are the vertices of $\Gamma\setminus\Delta$ that are joined to all vertices of $\Delta$. By \cite[Proposition~3.13]{AM}, the normalizer of $G_\Delta$ is then $G_\Delta\times G_{\Delta^{\perp}}$. Given a parabolic subgroup $P=gG_\Delta g^{-1}$, we let $P^\perp=gG_{\Delta^{\perp}}g^{-1}$ -- this is well defined independently of the choice of $g$. Then the normalizer of $P$ is $P\times P^\perp$.

We note that if $Q\subseteq P$ are two parabolic subgroups, then $P^\perp\subseteq Q^\perp$. Indeed, if $P$ and $Q$ are standard, this follows from the observation that if $\Upsilon\subseteq\Delta$ are two full subgraphs, then $\Delta^\perp\subseteq\Upsilon^\perp$. In general, this follows from \cite[Lemma~3.7]{AM}, which ensures that up to a global conjugation, we can always assume that $P$ and $Q$ are standard.

Given a vertex $v\in V\Gamma$, the \deffont{link} of $v$, denoted $\lk_\Gamma(v)$, is the full subgraph whose vertices are the vertices of $\Gamma$ that are adjacent to $v$ (in other words $\lk_\Gamma(v)=\{v\}^\perp$). The \deffont{star} of $v$, denoted $\st_\Gamma(v)$, is the full subgraph whose vertices are $v$ and those of $\lk_\Gamma(v)$.

Following \cite[Section~2]{CV}, we consider the following pre-ordering of the vertices of~$\Gamma$: for $v,w\in V\Gamma$, we have $v\le w$ if $\lk_\Gamma(v)\subseteq\st_\Gamma(w)$. We write $v\sim w$, and say that $v$ and $w$ are \deffont{equivalent}, if $v\le w$ and $w\le v$. {We write $v<w$ to mean that $v\le w$ and $v,w$ are not equivalent.} We let $G_v^+$ be the subgroup of $G$ generated by all vertex groups $G_w$ with $w\ge v$. More generally, given a parabolic subgroup $P=gG_vg^{-1}$ that is conjugate to a vertex group, we let $P^+=gG_v^+g^{-1}$. This definition does not depend on the choice of $g$: changing $g$ to $gg_1g_2$ with $g_1\in G_v$ and $g_2\in G_v^\perp$ does not affect the value of $gG_v^+g^{-1}$, because $g_1\in G_v^+$, and $g_2$ commutes with $G_v^+$ by definition.

\label{de:saturated} 
We say that a standard parabolic subgroup is \deffont{saturated} if whenever its type contains a vertex $v$, then it contains the whole equivalence class of $v$. More generally, a parabolic subgroup is \deffont{saturated} if it is conjugate to a saturated standard parabolic subgroup.

Any intersection of parabolic subgroups is a parabolic subgroup \cite[Corollaries~3.6 and~3.18]{AM}. In particular, any subgroup $A\subseteq G$ is contained in a unique smallest parabolic subgroup of $G$ (equal to the intersection of all parabolic subgroups of $G$ that contain $A$), called the \deffont{parabolic support} of $A$. We also define the parabolic support of an element $a\in G$ as the smallest parabolic subgroup of $G$ that contains $a$. 

\begin{Rq}\label{Fact:Normalizer-support}
    If $N\subseteq G$ normalizes $A$, then it normalizes the parabolic support $R$ of $A$. In particular, we have $N_G(A)=N_R(A)\times R^\perp$.
\end{Rq}

\begin{lemma}[{see e.g.\ \cite[Lemma 2.7]{EH}}]\label{lemma:parabolic-in-product}
    Let $G$ be a graph product over a finite simple graph $\Gamma$. Let $P=P_1\times P_2$ be a parabolic subgroup that splits as a direct product of two parabolic subgroups $P_1,P_2$. Let $Q\subseteq P$ be a parabolic subgroup.

    Then $Q=(Q\cap P_1)\times (Q\cap P_2)$.
\end{lemma}

\paragraph*{Parabolic free factors and direct factors.} Let $G$ be a graph product over a finite simple graph $\Gamma$, and let $P\subseteq G$ be a parabolic subgroup. Write $P=gG_\Delta g^{-1}$ for some $g\in G$ and some full subgraph $\Delta\subseteq \Gamma$. \label{thick}

By \cite[Lemma~4.7]{Gre}, if $\Delta$ is connected and has at least two vertices, then $G_\Delta$ does not split non-trivially as a free product. In general, let $\Delta_1,\dots,\Delta_k$ be the connected components of $\Delta$. Then $G_\Delta$ splits as a free product $G_{\Delta}=G_{\Delta_1}\ast\dots\ast G_{\Delta_k}$. This yields a free product decomposition $P=(gG_{\Delta_1}g^{-1})\ast\dots\ast (gG_{\Delta_k}g^{-1})$, that we call a \deffont{full parabolic free product decomposition} of $P$. Notice that any two such decompositions differ only by permutation of the factors, and conjugating all factors by the same element in $G_{\Delta}$. A \deffont{minimal parabolic free factor} of $P$ is a subgroup of $P$ that is conjugate to $gG_{\Delta_i}g^{-1}$ for some $i\in\{1,\dots,k\}$; it is \deffont{thick} if $|V\Delta_i|\ge 2$, and \deffont{thin} if $|V\Delta_i|=1$. We will use the phrases 
\deffont{thick free factor} and \deffont{thin free factor} as shortenings of \emph{thick minimal parabolic free factor} and \emph{thin minimal parabolic free factor}.

By e.g.\ \cite[Proposition~2.8]{Gen}, if $\Delta$ has at least two vertices and does not split non-trivially as a join, then $G_{\Delta}$ does not split non-trivially as a direct product. 
In general, let $\Delta'_1,\dots,\Delta'_\ell$ be the \deffont{join components} of $\Delta$: this means that $\Delta$ decomposes as a join $\Delta=\Delta'_1\circ\dots\circ\Delta'_\ell$ of non-empty subgraphs, in such a way that no subgraph $\Delta'_i$ splits non-trivially as a join. Equivalently, the join components of $\Delta$ are exactly the connected components of the \deffont{opposite graph} $\Delta^{\mathrm{opp}}$, having the same vertex set as $\Delta$, and where two vertices $v,w$ are adjacent if and only if they are non-adjacent in $\Delta$. This induces a decomposition $G_\Delta=G_{\Delta'_1}\times\dots\times G_{\Delta'_\ell}$ as a direct product.  Likewise we have $P=(gG_{\Delta'_1}g^{-1}) \times \dots \times (gG_{\Delta'_\ell}g^{-1})$, and this decomposition is independent of the choice of $g$. We call it the \deffont{full parabolic direct product decomposition} of $P$; it is well defined up to permutation of the factors. We say that $P$ is a \deffont{direct product parabolic subgroup} if this decomposition is non-trivial. A \deffont{minimal parabolic direct factor} of $P$ is a subgroup that is equal to one of the subgroups $gG_{\Delta'_i}g^{-1}$. Up to a permutation of the factors, we can assume that there exists $i_0\in\{0,\dots,\ell\}$ such that $\Delta'_1,\dots,\Delta'_{i_0}$ are the join factors of $\Delta$ reduced to one vertex. We then call $g(G_{\Delta'_1}\times\dots\times G_{\Delta'_{i_0}})g^{-1}$ the \label{de:clique}\deffont{clique factor} of $P$. We will also say that the subgraph spanned by $\Delta'_1,\dots,\Delta'_{i_0}$ is the \deffont{clique factor} of $\Delta$. Up to permutation, we can assume that there exists $i_1\ge i_0$ such that $\Delta'_{i_0+1},\dots,\Delta'_{i_1}$ are exactly the totally disconnected join components of $\Delta$ with at least two vertices and consisting of pairwise equivalent vertices. We then call $g(G_{\Delta'_1}\times\dots\times G_{\Delta'_{i_1}})g^{-1}$ the \deffont{extended clique factor}\label{de:extended-clique-factor} of $P$.

\begin{rk}\label{rk:center}
By \cite[Theorem~3.34]{Gre}, the center of every direct factor $G_{\Delta'_i}$ with $i>i_0$ is trivial. In particular 
\begin{equation*}
    Z(P)=g\left(Z\left(G_{\Delta'_1}\right)\times\dots\times Z\left(G_{\Delta'_{i_0}}\right)\right)g^{-1}.
\end{equation*}
\end{rk}

We now record a few lemmas that follow from the works of Minasyan--Osin \cite{MO} and Genevois \cite{Gen}.

\begin{lemma}\label{lemma:cu}
 Let $G$ be a graph product over a finite simple graph $\Gamma$, let $G=G_1\times\dots\times G_\ell$ be a full parabolic direct prduct decomposition of $G$, and let $A\subseteq G$ be a subgroup whose parabolic support is $G$.

 Then for every $i\in\{1,\dots,\ell\}$, there exists an element $g_i\in A$ whose parabolic support contains $G_i$.
\end{lemma}

\begin{proof}
Let $i\in\{1,\dots,\ell\}$, and let $r_i:G\to G_i$ be the projection. Then $r_i(A)$ is a subgroup of $G_i$ whose parabolic support is $G_i$. It thus follows from \cite[Theorem~6.16]{MO} that there exists $g_i\in A$ such that the parabolic support of $r_i(g_i)$ is $G_i$. Therefore the parabolic support of $g_i$ contains $G_i$. 
\end{proof}

\begin{lemma}\label{lemma:maximal-direct-product}
Let $G$ be a graph product over a finite simple graph $\Gamma$. Let $H\subseteq G$ be a subgroup that is not contained in a conjugate of a vertex group. The following assertions are equivalent.
\begin{enumerate}
\item\label{item:product-1} $H$ is a maximal subgroup of $G$ that splits non-trivially as a direct product.
\item\label{item:product-2} $H$ is a maximal direct product parabolic subgroup of $G$. 
\end{enumerate}
\end{lemma}

\begin{proof}
The implication \ref{item:product-1}$\Rightarrow$\ref{item:product-2} is the content of \cite[Proposition~2.8]{Gen}. For \ref{item:product-2}$\Rightarrow$\ref{item:product-1}, let $H$ be as in \ref{item:product-2}, and let $\hat H$ be a subgroup of $G$ that contains $H$ and splits non-trivially as a direct product. By \cite[Corollary~6.15]{MO}, the group $\hat H$ is contained in a direct product parabolic subgroup $P$. By maximality we have $H=P$ and in particular $H=\hat H$.
\end{proof}


\paragraph*{Normal forms in graph products.}\label{Def:Normal-forms}

Let $G$ be a graph product over a finite simple graph $\Gamma$. A \emph{$G$-word} is a tuple $\underline{\mathsf{w}}=(g_1,\dots,g_k)$, where every $g_i$ is an element of some vertex group. 
The $G$-word $\underline{\mathsf{w}}$ \emph{represents} $g\in G$ if $g=g_1\dots g_k$. A word $\underline{\mathsf{w}}$ is \emph{graphically reduced} if, letting $v_i\in V\Gamma$ be such that $g_i\in G_{v_i}$, the following two conditions hold: 
\begin{itemize}
    \item for every $i\in\{1,\dots,k\}$, one has $g_i\neq e$;
    \item for every $i<j$, if $v_i=v_j$, then there exists $i<k<j$ such that $v_k$ is neither equal nor adjacent to $v_i$. 
\end{itemize}
It is \deffont{graphically cyclically reduced} if in addition, given any $i<j$ in $\{1,\dots,k\}$, if $v_i=v_j$, then either there exists $l<i$ such that $v_l$ is neither equal nor adjacent to $v_i$, or there exists $l>j$ such that $v_l$ is neither equal nor adjacent to $v_j$. An element $g\in G$ is \deffont{graphically cyclically reduced} if it is represented by a graphically cyclically reduced $G$-word.

Green proved in \cite[Theorem~3.9]{Gre} that every element of $G$ is represented by a graphically reduced $G$-word. In addition, two graphically reduced words $\underline{\mathsf{w}}=(g_1,\dots,g_k)$ and $\underline{\mathsf{w}}'=(g'_1,\dots,g'_\ell)$ represent the same element if and only if $k=\ell$, and $\underline{\mathsf{w}}$ and $\underline{\mathsf{w}}'$ are obtained from one another by successive applications of the following operation: if $g_i$ and $g_{i+1}$ belong to adjacent vertex groups, swap them. In particular, if two graphically reduced words $\underline{w}$ and $\underline{w}'$ represent the same element $g\in G$, for every $v\in V\Gamma$, if there exists $i\in\{1,\dots,k\}$ such that $g_i\in G_v$, then there exists $j\in\{1,\dots,\ell\}$ such that $g'_{j}\in G_v$.

Green also proved in \cite[Lemma~3.16]{Gre} that every element of $G$ is conjugate to a graphically cyclically reduced element.

If an element $g$ is graphically cyclically reduced, represented by a graphically cyclically reduced $G$-word $\underline{\mathsf{w}}$, and if $\Lambda\subseteq\Gamma$ is the smallest subgraph such that every syllable of $\underline{\mathsf{w}}$ belongs to $G_\Lambda$, then $G_\Lambda$ is the parabolic support of $g$: this can be proved using, for instance, \cite[Theorem~5.5]{Gen2}.

\paragraph*{Centralizers in graph products.} We will need the following theorem of Barkauskas \cite{Bar}, which describes centralizers in a graph product (see also \cite[Proposition~2.6]{Gen} for a rephrasing in a terminology closer to ours).

\begin{theo}[Barkauskas \cite{Bar}] \label{theo:barkauskas}
{Let $G$ be a graph product over a finite simple graph~$\Gamma$, with vertex groups $(G_v)_{v\in V\Gamma}$. Let $g\in G$, let $P$ be the parabolic support of $g$, and let $P=P_1\times\dots\times P_k$ be a full parabolic direct product decomposition of $P$, chosen so that $P_1\times\dots\times P_{i_0}$ is the clique factor of $P$. Write $g=g_1\dots g_k$, with $g_i\in P_i$. Then 
\begin{equation*}
    C_G(g)=C_{P_1}(g_1)\times\dots\times C_{P_{i_0}}(g_{i_0})\times Z_{i_0+1}\times\dots\times Z_k\times P^\perp,
\end{equation*}
where for every $i\in\{i_0+1,\dots,k\}$, $Z_i$ is the maximal (infinite) cyclic subgroup of $P_i$ that contains $g_i$.}  
\end{theo}

\paragraph{Collapsing subgraphs.} Recall that a full subgraph $\Delta\subseteq\Gamma$ is \deffont{collapsible} if for any two vertices $x,y\in V\Delta$, we have $\lk_\Gamma(x)\cap (\Gamma\setminus\Delta)=\lk_\Gamma(y)\cap (\Gamma\setminus\Delta)$. Equivalently, for every vertex $x\in V\Delta$, one has $\lk_\Gamma(x)\subseteq \Delta\circ\Delta^{\perp}$. Recall that a graph is \deffont{strongly reduced} if it does not contain any proper collapsible subgraph having at least two vertices. It is \deffont{almost strongly reduced} if every proper collapsible subgraph on at
least two vertices is edgeless. Note that a graph is almost strongly reduced if and only if it does not contain any connected collapsible subgraph on at least two vertices.

\begin{ex} In Figure~\ref{fig:collapsing} the rightmost graph is strongly reduced and hence almost strongly reduced while the graph in the middle is almost strongly reduced but not strongly reduced.  
\end{ex}



\begin{lemma}\label{Lmm:vertex-group-iso-strong-redution}
Let $G,H$ be graph products over finite simple graphs $\Gamma,\Lambda$, not reduced to one vertex, such that $\Lambda$ is strongly reduced. Let $\varphi:G\to H$ be an isomorphism. Let $v\in V\Gamma$, let $Q$ be a parabolic subgroup of $H$, and assume that $\varphi(G_v)=Q$.

Then $Q$ is conjugate to a vertex group of $H$.
\end{lemma}

\begin{proof}
Assume otherwise. Since $\Gamma$ is not reduced to one vertex, we have $\varphi(G_v)\neq H$, so the type of $Q$ is a proper subgraph of $\Lambda$, which we denote by $\Upsilon$. Since $\Lambda$ is strongly reduced, the subgraph $\Upsilon$ is not collapsible, so there exists a vertex $w\in V\Upsilon$ such that $\lk_\Lambda(w)\not\subseteq\Upsilon\circ \Upsilon^\perp$, in particular $N_H(H_w)\not\subseteq N_H(Q)$. On the other hand, for every non-trivial subgroup $A\subseteq G_v$, the parabolic support of $A$ is $G_v$, and therefore $N_G(A)\subseteq N_G(G_v)$ by Remark~\ref{Fact:Normalizer-support}. This is a contradiction.
\end{proof}

Given a set $\calf$ of pairwise disjoint collapsible subgraphs of $\Gamma$, we now define a new graph $\overline{\Gamma}^{(\calf)}$ by \deffont{collapsing} $\calf$, as follows. First, let $\hat\calf$ be obtained from $\calf$ by adding all singletons $\{v\}$, with $v\in V\Gamma$ not contained in any subgraph in $\calf$. In this way, the subgraphs in $\hat\calf$ are collapsible, and $\hat\calf$ induces a partition of $V\Gamma$. Collapsibility implies that if $\Delta,\Upsilon\in\hat{\calf}$, either all vertices of $\Delta$ are adjacent to all vertices of $\Upsilon$, or no vertex of $\Delta$ is adjacent to a vertex of $\Upsilon$. We then let $\overline{\Gamma}^{(\calf)}$ be the graph whose vertex set is $\hat\calf$, where $\Delta,\Upsilon\in\hat\calf$ are joined by an edge if all vertices of $\Delta$ are adjacent to all vertices of $\Upsilon$. It comes with a map $\pi_\calf:\Gamma\to\overline{\Gamma}^{(\calf)}$, sending every vertex $v\in V\Gamma$ to the unique $\Delta\in\hat\calf$ with $v\in V\Delta$. Note that this map sends every vertex to a vertex, and every edge $e=vw$ to either a vertex (if $v,w$ belong to the same subgraph in $\hat\calf$) or to an edge (otherwise). We call $\pi_\calf$ a \deffont{graph collapse}. Note that $G$ splits as a graph product over $\overline{\Gamma}^{(\calf)}$, with vertex groups $(G_\Delta)_{\Delta\in\hat\calf}$.
A \deffont{strong reduction} of $\Gamma$ is a strongly reduced graph $\overline{\Gamma}$ obtained by collapsing a family $\calF$ of pairwise disjoint proper collapsible subgraphs of $\Gamma$, equipped with a {graph collapse} $\pi:\Gamma\to\overline{\Gamma}$ (this means more precisely that there exists an isomorphism between $\overline{\Gamma}$ and $\overline{\Gamma}^{(\calf)}$ that identifies $\pi$ with $\pi_\calf$). 
\begin{ex}
    The leftmost graph on Figure~\ref{fig:collapsing}, denoted by $\Gamma$, has three proper collapsible subgraphs that are not reduced to one vertex: one collapsible edge –~whose vertices are in \textcolor{DarkOrange3}{orange}~– and two collapsible disconnected subgraphs on two vertices –~whose vertices are in \textcolor{marronfonce}{brown}. Collapsing only the \textcolor{DarkOrange3}{orange} subgraph gives the graph in the middle. Collapsing moreover the two \textcolor{marronfonce}{brown} subgraphs gives the strong reduction of $\Gamma$ on the right.
\end{ex}
\begin{figure}[htbp]
    \centering
    \includegraphics[width=0.97\textwidth]{Tikz/Almost_and_strong_reduction.pdf}
    \caption{Examples of graph collapse images.}
    \label{fig:collapsing}
\end{figure}

\begin{rk}\label{Rq:strong-reduction-as-at-least-3-vertices}\label{Rq:Unicite-des-reductions}
If $\Gamma$ is disconnected, then any strong reduction $\overline\Gamma$ of $\Gamma$ is an edgeless graph with two vertices. But the family $\mathcal{F}$ such that $\overline{\Gamma}$ is isomorphic to $\Gamma^{(\mathcal{F})}$ is not unique
as soon as $\Gamma$ has at least three connected components.

If $\Gamma$ splits non-trivially as a join, then any strong reduction $\overline\Gamma$ of $\Gamma$ is a graph reduced to one edge. As above, the family $\mathcal{F}$
is not unique as soon as $\Gamma$ has at least three join components. See e.g.\ Remark~\ref{Rq:at-least-three-vertices}.

If $\Gamma$ is connected and does not split non-trivially as a join, then the strong reduction $\overline{\Gamma}$ of $\Gamma$ is unique (and it has at least three vertices, otherwise $\Gamma$ is disconnected or splits as a join). This can be proved using that if two proper subgraphs $\Delta_1,\Delta_2\subseteq\Gamma$ are collapsible and have non-empty intersection, then the full subgraph induced by $\Delta_1\cup\Delta_2$ is again a proper collapsible subgraph. Therefore, the maximal (for inclusion) {proper} collapsible subgraphs of $\Gamma$ are pairwise disjoint. Letting $\calf$ be the set of all maximal {proper} collapsible subgraphs of $\Gamma$, we then have $\overline{\Gamma}=\overline{\Gamma}^{(\calf)}$, with the associated graph collapse.

Note that if we had taken for $\calf$ the set of all maximal \emph{connected} collapsible subgraphs of $\Gamma$, then the graph $\overline{\Gamma}^{(\calf)}$ would have been almost strongly reduced, but not necessarily strongly reduced.
\end{rk}

\paragraph*{Basic automorphisms and graph collapses.} Remark that the definition of basic automorphisms given in Definition~\ref{de:basic} is set forth \emph{once a presentation of $G$ has been chosen}. When working at the same time with a presentation of a graph product and a collapse of it, we will thus need to keep track of the defining graph. We will call \deffont{$\Gamma$-local automorphisms, $\Gamma$-partial conjugations, $\Gamma$-twists} and \deffont{$\Gamma$-folds} the automorphisms arising as basic automorphisms of $G$ seen as a graph product over $\Gamma$, as defined in Definition~\ref{de:basic}.

\begin{lemma}[Basic automorphisms and graph collapses]\label{Lmm:Aut-and-graph-collapse}
    Let $\Gamma$ be a finite simple graph and let $G$ be a graph product over $\Gamma$ with vertex groups $(G_v)_{v\in V\Gamma}$. Let $\pi:\Gamma \rightarrow \overline{\Gamma}$ be a graph collapse of $\Gamma$. Let $\varphi \in \Aut(G)$.
    \begin{enumerate}
        \item If $\varphi$ is a $\overline{\Gamma}$-twist, then it is a product of $\Gamma$-twists.
        \item If $\varphi$ is a $\overline{\Gamma}$-fold, then it is a product of $\Gamma$-folds.
        \item If $\varphi$ is a $\overline{\Gamma}$-partial conjugation, then it is a product of $\Gamma$-partial conjugations.
    \end{enumerate}
\end{lemma}

\begin{proof} 
The lemma is satisfied if $|{V\overline\Gamma}|\le 1$. We now assume $|{V\overline\Gamma}|\geq 2$. In the following we denote by $\bar{v}$ and $\bar{w}$ two distinct vertices of $\overline{\Gamma}$ and let $\Delta:=\pi^{-1}(\bar{v})$ and $\Upsilon:=\pi^{-1}(\bar{w})$.

\medskip

\textbf{1 – $\overline{\Gamma}$-twists.} Assume that $\st_{\overline{\Gamma}}(\bar{v})\subseteq \st_{\overline{\Gamma}}(\bar{w})$. Let $\twistor:G_{\bar{v}}\rightarrow Z(G_{\bar{w}})$ and $\varphi\in \Aut(G)$ be such that $\varphi(g)= g\twistor(g)$ for all $g\in G_{\bar{v}}$ and $\varphi(h)=h$ if $h\in G_{\overline{\Gamma}\backslash \{\bar{v}\}}$. We aim to prove that $\varphi$ is a product of $\Gamma$-twists.

First let $w_1,\dots,w_{i_0}$ denote the vertices of the clique factor of $\Upsilon$. 
Then, by Remark~\ref{rk:center}, we have 
\begin{equation*}
    Z\left(G_{\bar{w}} \right)
    =Z\left( G_{w_1}\right)\times \dots \times Z\left( G_{w_{i_0}}\right).
\end{equation*}
So for all $j\in \{1,\dots,i_0\}$ there exists $\twistor_j:G_{\bar{v}}\rightarrow Z\left( G_{w_j}\right)$ satisfying $ \twistor(g)=\twistor_1(g)\cdots \twistor_{i_0}(g)$ for all $g\in G_{\bar v}$.

Now let $v\in V \Delta$ and $j\in \{1,\dots,i_0\}$ and let us prove that $\st_\Gamma(v)\subseteq \st_\Gamma(w_j)$. Since $\bar{v}\in \st_{\overline{\Gamma}}(\bar{w})$ and by definition of $\Delta$ and $\Upsilon$, any vertex of $\Delta$ is adjacent to any vertex of $\Upsilon$. Therefore,
\begin{itemize}
    \item $v\in \st_\Gamma(w_j)$;
    \item and using that $w_j$ is in the clique factor of $\Upsilon$, we obtain that if $u\in \lk_\Gamma(v)\cap (\Delta\cup \Upsilon)$ then $u\in \st_\Gamma(w_j)$.
\end{itemize}
It thus only remains to show that $\lk_{\Gamma \backslash(\Delta\cup \Upsilon)}(v)\subseteq \st_\Gamma(w_j)$. So let $u\in \lk_{\Gamma \backslash(\Delta\cup \Upsilon)}(v)$. Then $\pi(u)$ belongs to $\lk_{\overline{\Gamma}}(\bar{v})\setminus\{\bar w\}$, which in turn is contained in $\lk_{\overline{\Gamma}}(\bar{w})$. This implies that $u\in \lk_\Gamma(w)$ for all $w\in V\Upsilon$. 

Therefore, $\st_\Gamma(v)\subseteq \st_\Gamma(w_j)$. 
In particular, the automorphism $\varphi_{j,v}$ defined by $\varphi_{j,v}(g)=g\twistor_j(g)$ for all $g\in G_v$ and $\varphi_{j,v}(h)=h$ for all $h\in G_{\Gamma\backslash \{v\}}$ is a $\Gamma$-twist. The wanted assertion now follows noting that $\varphi$ can be written as a product of such $\varphi_{j,v}$.
\medskip

\textbf{2 – $\overline{\Gamma}$-folds.} Let $\bar{v},\bar{w}\in V\Gamma$ be such that $\lk_{\overline{\Gamma}}(\bar{v})\subseteq \lk_{\overline{\Gamma}}(\bar{w})$ and $G_{\bar{v}}\simeq \bZ$. Note that in this case $\Delta$ is reduced to a vertex, denoted by $v$. 

Let $z$ be a generator of $G_{\bar{v}}=G_v$ and let $\varphi\in \Aut(G)$ be such that $\varphi(z)=zg'$ for some element $g'$ belonging to $G_{\bar{w}}=G_\Upsilon$, and $\varphi(h)=h$ for every $h\in G_{\overline{\Gamma}\setminus\{\bar{v}\}}$. By definition of a collapsible subgraph, for all $w\in V\Upsilon$ we have $\lk_\Gamma(v)\subseteq \lk_\Gamma(w)$. Therefore, writing $g'$ as a product of elements of the $G_w$ for $w\in V\Upsilon$, we get that $\varphi$ is a product of $\Gamma$-folds.
\medskip

\textbf{3 – $\overline{\Gamma}$-partial conjugations.} Let $g\in G_{\bar{v}}$, let $\overline{C}$ be a connected component of $\overline{\Gamma}\backslash\st_{\overline{\Gamma}}(\bar v)$ and define $\varphi\in \Aut(G)$ such that $\varphi(h)=ghg^{-1}$ for all $h\in G_{\overline{C}}$ and $\varphi(h)=h$ for all $h\in G_{\overline{\Gamma}\backslash\overline{C}}$. 

Let $C:=\pi^{-1}(\overline{C})$ and let us prove that $C$ is a union of connected components of $\Gamma\backslash\st_\Gamma(v)$ for all $v\in V\Delta$. To this end, letting $u\in V\Gamma\backslash VC$ be such that $u$ is adjacent to some vertex $u'$ in $C$, we aim to prove that $u\in\lk_\Gamma(v)$ for all $v\in V\Delta$. Let $\bar u=\pi(u)$.
\begin{itemize}
    \item We first prove that $\bar u\in\st_{\overline{\Gamma}}(\bar v)$. So assume for the sake of contradiction that this fails, and let $\overline{C}'$ be the connected component of $\overline{\Gamma}\backslash\st_{\overline{\Gamma}}(\bar v)$ containing $\bar u$. By assumption $u\notin VC$ so  $\bar{u}\notin V\overline{C}$ and thus $\overline{C}'\neq \overline{C}$. Moreover, since $u\notin VC$ and is {adjacent} to $u'\in VC$, then $\bar{u}$ is adjacent to $\pi(u')\in V\overline{C}$. This implies that $\overline{C}'=\overline{C}$, which is the desired contradiction. 
    \item We now prove that $\bar u\neq\bar{v}$. Since $u$ is adjacent to $u'$, we have $d_{\overline\Gamma}(\bar{u},\pi(u'))\le 1$. On the other hand, since $\pi(u')\in \overline{C}$ and $\overline{C}$ is a connected component of $\overline{\Gamma}\setminus\st_{\overline{\Gamma}}(\bar{v})$, we have $d_{\overline{\Gamma}}(\pi(u'),\bar{v})\ge 2$. So $\bar{u}\neq \bar{v}$.
    \item It follows from the above two points that $\bar{u}\in \lk_{\overline{\Gamma}}(\bar{v})$, so by definition of $\pi$ we have $u\in\lk_\Gamma(v)$ for all $v\in V\Delta$.
\end{itemize}

This shows that $C$ is a union of connected components of $\Gamma\backslash \st_{\Gamma}(v)$ for all $v\in V\Delta$.

Then, since $g$ belongs to $G_{\bar{v}}=G_{\Delta}$, we can decompose $g$ as a product of elements of elements of vertex groups of $G_{\Delta}$ and thus write $\varphi$ as a composition of $\Gamma$-partial conjugations.
\end{proof}

\begin{lemma}[Extending automorphisms defined on collapsible subgraphs]\label{Lmm:Extending-aut-defined-on-collapsible}
Let $\Gamma$ be a finite simple graph, let $G$ be a graph product over $\Gamma$ with vertex groups $(G_v)_{v\in V\Gamma}$. Let $\Delta\subseteq\Gamma$ be a collapsible subgraph, and view $G_{\Delta}$ as a graph product over $\Delta$ with vertex groups $(G_v)_{v\in V\Delta}$. 

Then every automorphism $\alpha_\Delta\in\Aut(G_\Delta)$ extends uniquely to an automorphism $\alpha\in\Aut(G)$ which restricts to the identity on $G_{\Gamma\setminus\Delta}$. If in addition $\alpha_\Delta$ is a $\Delta$-local automorphism (resp.\ a $\Delta$-partial conjugation, a $\Delta$-twist, a $\Delta$-fold), then $\alpha$ is a $\Gamma$-local automorphism (resp.\ a $\Gamma$-partial conjugation, a $\Gamma$-twist, a $\Gamma$-fold). 
\end{lemma}

\begin{proof}
  Let $\pi_\Delta:\Gamma \rightarrow \overline{\Gamma}^{(\Delta)}$ be the graph collapse. Then $\alpha_\Delta$ is an automorphism of the vertex group $G_{\pi_\Delta(\Delta)}$, therefore it extends (uniquely) to a $ \overline{\Gamma}^{(\Delta)}$-local automorphism of $G$, denoted by $\alpha$, which is the identity on every vertex group $G_v$ with $v\in V\Gamma\backslash V\Delta$. The additional assertion readily follows for local automorphisms, we now prove it for partial conjugations, twists and folds. 

  The important observation, which follows from the collapsibility of $\Delta$, is that if $v\in V\Delta$, then
  \begin{equation}\label{lkst}
    \lk_\Gamma(v)=\lk_{\Delta}(v)\circ\Delta^{\perp} \quad \text{and} \quad \st_{\Gamma}(v)=\st_{\Delta}(v)\circ\Delta^{\perp}.
  \end{equation}

  For partial conjugations, if $\st_{\Delta}(v)$ disconnects $\Delta$, and if $C\subseteq \Delta$ is a connected component of $\Delta\setminus\st_{\Delta}(v)$, it follows from Equation~\eqref{lkst} that $C$ is a connected component of $\Gamma\setminus\st_{\Gamma}(v)$. Therefore every $\Delta$-partial conjugation extends to a $\Gamma$-partial conjugation which is the identity on $G_{\Gamma\setminus\Delta}$.

  For twists and folds, if $v,w\in V\Delta$ and if $\lk_{\Delta}(v)\subseteq\st_\Delta(w)$, then Equation~\eqref{lkst} ensures that $\lk_{\Gamma}(v)\subseteq\st_\Gamma(w)$. It follows that every $\Delta$-twist (resp.\ $\Delta$-fold) extends to a $\Delta$-twist (resp.\ $\Delta$-fold) which restricts to the identity on $G_{\Gamma\setminus\Delta}$. 
\end{proof}

\subsection{Background on free products}

\paragraph*{General background.} As explained in the introduction, the proofs of our main theorems will require to consider free products. Here we gather some terminology and standard facts about them. 

A group $G$ is \deffont{freely indecomposable} if it does not split as a free product of two non-trivial groups, and is not isomorphic to $\mathbb{Z}$ (equivalently, every action of $G$ on a simplicial tree with trivial edge stabilizers has a global fixed point). 

A \deffont{Grushko decomposition} of a group $G$ is a splitting of $G$ as a free product of the form $G=G_1\ast\dots\ast G_k\ast F$, where $G_1,\dots,G_k$ are freely indecomposable, and $F$ is a finitely generated free group. By Grushko's theorem \cite{Gru}, every finitely generated group has a Grushko decomposition. In addition, if $G$ has two Grushko decompositions $G=G_1\ast\dots\ast G_k\ast F$ and $G=G'_1\ast\dots\ast G'_{k'}\ast F'$, then the following hold:
\begin{itemize}
\item $k=k'$ and there exists a bijection $\sigma:\{1,\dots,k\}\to\{1,\dots,k'\}$ such that for every $i\in\{1,\dots,k\}$, the subgroups $G_i$ and $G'_{\sigma(i)}$ are conjugate in $G$;
\item the free groups $F$ and $F'$ have the same rank.
\end{itemize}
\begin{rk}\label{rk:grushko}
  If we have two decompositions $G=A_1\ast\dots\ast A_m$ and $G=B_1\ast\dots\ast B_\ell$ where each subgroup $A_i$ and $B_j$ is either freely indecomposable or isomorphic to $\mathbb{Z}$, then $\ell=m$ and up to permuting the factors, we have $A_i\simeq B_i$ for every $i\in\{1,\dots,m\}$.
\end{rk}

A \deffont{free factor} of a group $G$ is a subgroup $A\subseteq G$ such that there exists $B\subseteq G$ with $G=\langle A,B\rangle$ and $\langle A,B\rangle\simeq A\ast B$. We will use the following facts satisfied by any free factor $A\subseteq G$.
\begin{itemize}
\item $A$ is \deffont{malnormal}, i.e.\ for every $g\in G$, either $g\in A$ or $gAg^{-1}\cap A=\{1\}$. \label{Def:Malnormal}
\item If $H\subseteq G$ is a subgroup, then $A\cap H$ is a free factor of $H$.
\end{itemize}

\paragraph*{Automorphisms of free products.} Let $G$ be a group that splits as a free product $G=G_1\ast\dots\ast G_k\ast F_N$, where $F_N=\langle x_1,\dots,x_N\rangle$ is a rank $N$ free group (we no longer assume that the subgroups $G_i$ are freely indecomposable). We denote by $\mathcal{F}$ the set of all conjugacy classes in $G$ of the subgroups $G_i$. We let $\Aut(G,\calf)$ be the subgroup of $\Aut(G)$ that preserves $\calf$ setwise. We let $\Aut(G,\calf^{(t)})$ be the subgroup consisting of all automorphisms that restrict to a conjugation on each subgroup $G_i$ (where the conjugating element belongs to $G$ and depends on $i$). Note that if the above free product decomposition of $G$ is a Grushko decomposition, then $\Aut(G,\calf)=\Aut(G)$.

In \cite{FR}, Fouxe-Rabinovitch gave generating sets for $\Aut(G,\calf^{(t)})$ and $\Aut(G,\calf)$, extending the one given by Nielsen for $\Aut(F_N)$ in \cite{Nie}.

\begin{theo}[{Fouxe-Rabinovitch \cite{FR}}]\label{theo:FR}
Let $G=G_1\ast\dots\ast G_k\ast F_N$ be a free product, where $F_N=\langle x_1,\dots,x_N\rangle$ is a rank $N$ free group, and let $\calf$ be the set of all conjugacy classes of the subgroups $G_i$. 

\noindent $\drsh$ The group $\Aut(G,\calf^{(t)})$ is generated by the following automorphisms:
\begin{enumerate}
    \item \emph{transvections}, i.e.\ automorphisms that 
    \begin{itemize}
    \item send $x_i$ to $x_ig$ for some $i\in\{1,\dots,N\}$ and some element $g\in G$ which either belongs to some $G_j$, or is equal to a power of some $x_l$ with $l\neq i$, 
    \item and restrict to the identity on every $G_j$ and send every $x_j$ with $j\neq i$ to itself;
    \end{itemize}
    \item \emph{partial conjugations}, i.e.\ automorphisms that 
    \begin{itemize}
    \item restrict to the conjugation by an element $g$ to one subgroup $G_i$, for some $g$ that either belongs to some $G_j$ or is equal to a power of some $x_l$,
    \item and are the identity on all  subgroups $G_j$ with $j\neq i$ and on $F_N$;
    \end{itemize}
    \item \emph{inversions}, i.e.\ automorphisms that send one $x_i$ to $x_i^{-1}$, are the identity on every $G_j$, and fix every $x_j$ with $j\neq i$;
    \item \emph{permutations}, i.e.\ automorphisms that are the identity on every $G_i$, and send every $x_i$ to $x_{\sigma(i)}$ for some permutation $\sigma$ of $\{1,\dots,N\}$.
\end{enumerate}
$\drsh$ The group $\Aut(G,\calf)$ is generated by the above four types of automorphisms together with:
\begin{enumerate}
    \item[5.] automorphisms that restrict to an automorphism of some $G_i$, are identity on all $G_j$ with $j\neq i$, and fix all $x_l$;
    \item[6.] automorphisms which, given some isomorphism $\varphi_{ij}:G_i\to G_j$ with $i\neq j$, restrict to $\varphi_{ij}$ on $G_i$ and to $\varphi_{ij}^{-1}$ on $G_j$, are identity on all $G_l$ with $l\notin\{i,j\}$ and fix all $x_p$.
\end{enumerate}
\end{theo}

\begin{rk}\label{rk:FR}
The above says that when viewing $G$ as a graph product on an edgeless graph with $k+N$ vertices, with $k$ vertices carrying the groups $G_1,\dots,G_k$ and the last $N$ vertices each carrying a vertex group isomorphic to $\mathbb{Z}$, then the basic automorphisms listed in Definition~\ref{de:basic} generate $\Aut(G,\calf)$. Note that there are no twists. 
\end{rk}

\begin{lemma}\label{lemma:fouxe-rabinovitch}
Let $G=G_1\ast\dots\ast G_k\ast F_N$, where $F_N=\langle x_1,\dots,x_N\rangle$ is a rank $N$ free group. Let $\calf$ be the set of all conjugacy classes in $G$ of the subgroups $G_i$.

The orbits of $(x_1,\dots,x_N)$ under $\Aut(G,\calf)$ and $\Aut(G,\calf^{(t)})$ coincide.
\end{lemma}

\begin{proof}
We will prove that every automorphism $\varphi\in\Aut(G,\calf)$ can be written as $\varphi=\varphi_0\psi$, with $\varphi_0\in\Aut(G,\calf^{(t)})$ and $\psi\in \Aut(G,\calf)$ such that for every $\ell\in\{1,\dots,N\}$, one has $\psi(x_\ell)=x_\ell$. The lemma follows from this fact.

First, we explain how to write $\varphi=\varphi_1\psi_1$, with $\psi_1(x_\ell)=x_\ell$ for every $\ell\in\{1,\dots,N\}$, and $\varphi_1(G_i)$ conjugate to $G_i$ for every $i\in\{1,\dots,k\}$. Let  $\sigma:\Aut(G,\calf)\to\mathfrak{S}_k$ be the homomorphism defined by letting $\sigma(\varphi)(i)=j$ whenever $\varphi(G_i)$ is conjugate to $G_j$. For any $i,j\in\{1,\dots,k\}$ such that $G_i$ is isomorphic to $G_j$, we fix an isomorphism $\varphi_{i,j}:G_i\to G_j$, in such a way that $\varphi_{i,k}=\varphi_{j,k}\circ\varphi_{i,j}$ whenever $G_i,G_j,G_k$ are all isomorphic. This yields a section $\kappa$ of $\sigma$: for $\tau\in\mathrm{im}(\sigma)$, we let $\kappa(\tau)$ be the automorphism such that, for all $i\in \{1,\dots,k\}$, $\kappa(\tau)$ is $\varphi_{i,\tau(i)}$ in restriction to $G_i$, and it fixes $x_1,\dots,x_N$. Letting $\psi_1=\kappa\circ\sigma(\varphi)$ and $\varphi_1=\varphi\circ\psi_1^{-1}$ then gives the desired decomposition.

Let $\Aut^0(G,\calf)$ be the kernel of $\sigma$. Now, for every $i\in\{1,\dots,k\}$, there is a homomorphism $\mathrm{res}_i:\Aut^0(G,\calf)\to\Out(G_i)$, defined as follows. For $\alpha\in\Aut^0(G,\calf)$, we have $\alpha(G_i)=gG_ig^{-1}$ for some $g\in G$. We then let $\mathrm{res}_i(\alpha)$ be the outer class of $(\mathrm{ad}_{g^{-1}}\circ\alpha)_{|G_i}$. Since $G_i$ is its own normalizer in $G$, this outer class does not depend on the choice of $g$ as above. We now define a set-theoretic section $s_i$ of $\mathrm{res}_i$, i.e.\ a map $s_i:\Out(G_i)\to\Aut^0(G,\calf)$ such that $\mathrm{res}_i\circ s_i=\mathrm{id}_{\Out(G_i)}$. First, choose a set-theoretic section $\alpha_i:\Out(G_i)\to\Aut(G_i)$. Now let $s_i:\Out(G_i)\to\Aut^0(G,\calf)$ be defined by letting $s_i(\Phi)$ be the automorphism that restricts to $\alpha_i(\Phi)$ on $G_i$, and is the identity on all $G_j$ with $j\neq i$ and on $F_N$. 

Letting $\psi=(s_1\circ\mathrm{res}_1(\varphi_1))\circ\dots\circ (s_k\circ\mathrm{res}_k(\varphi_1))\circ\psi_1$, and $\varphi_0=\varphi\circ\psi^{-1}$, yields the conclusion.
\end{proof}

% ======================================
% Part 1 - Isomorphisms
% ======================================
\part{The isomorphism problem}\label{part:isomorphism}
The goal of this first part is to solve the isomorphism problem for graph products (Theorem~\ref{theo:isomorphism}), under the assumption that the graphs $\Gamma,\Lambda$ are strongly reduced.

\section{Blueprint for the isomorphism problem}\label{sec:blueprint}

In this section we describe our general strategy for solving the isomorphism problem. Assume that we have an isomorphism $\varphi:G\to H$ between two graph products over finite simple graphs $\Gamma,\Lambda$. Under natural sets of conditions on the graph products, we aim to derive a graph isomorphism $\theta:\Gamma\to\Lambda$. If we knew that $\varphi(G_v)$ was conjugate to a vertex group $H_w$ for every $v\in V\Gamma$, then the assignment $v\mapsto w$ would yield such an isomorphism (see \cite[Theorem~3.11]{GM}). But to solve the isomorphism problem, it will be enough to build an isomorphism from $G$ to $H$ that satisfies a weaker condition, encapsulated by the following definitions.
Recall that for $v\in V\Gamma$ we denote by $r_v:G\rightarrow G_v$ the parabolic retraction on $G_v$, namely the group homomorphism defined as being the identity on $G_v$ and the trivial homomorphism on $G_u$, for any $u\in V\Gamma\backslash \{v\}$.

\begin{de}[$\varphi$-paired]\label{de:paired}
Let $\Gamma,\Lambda$ be finite simple graphs. Let $G,H$ be graph products over $\Gamma,\Lambda$, respectively, and let $\varphi:G\to H$ be an isomorphism. Given $v\in V\Gamma$ and $w\in V\Lambda$, we say that $(v,w)$ is \deffont{$\varphi$-paired} if the following conditions hold: 
\begin{itemize}
    \item $(r_w\circ\varphi)_{|G_v}$ is an isomorphism from $G_v$ to $H_w$;
    \item $(r_v\circ \varphi^{-1})_{|H_w}$ is an isomorphism from $H_w$ to $G_v$.
\end{itemize}
We say that $(v,w)$ is \deffont{exclusively $\varphi$-paired} if $w$ is the unique vertex in $V\Lambda$ such that $(v,w)$ is $\varphi$-paired, and $v$ is the unique vertex in $V\Gamma$ such that $(w,v)$ is $\varphi^{-1}$-paired. 
\end{de}

\begin{ex}\label{Rk:vertex-exclusively-phi-paired}
We give three examples.
\begin{enumerate}
\item If $\varphi(G_v)=hH_wh^{-1}$ for some $h\in H$, then $(v,w)$ is exclusively $\varphi$-paired. 
\item \label{item:Ex_phi_pared2}Let $G$ be the graph product over an edge $e=vw$ with vertex groups $G_v$ and $G_w$, so that $G=G_v\times G_w$. Assume that we are given a homomorphism $\twistor:G_w\to Z(G_v)$, and let $\varphi:G\to G$ be the isomorphism (a twist) which is the identity on $G_v$, and sends every $g\in G_w$ to $g\twistor(g)$. Then $(v,v)$ and $(w,w)$ are both exclusively $\varphi$-paired, even though $\varphi(G_w)$ is not conjugate to any vertex group as soon as $\twistor$ is non-trivial.
\item \label{item:Ex_phi_pared3}Let $G$ be the graph product over an edgeless graph with two vertices $v,w$, with vertex groups $G_v$ and $G_w$, so that $G=G_v*G_w$. Assume that  $G_w\simeq\mathbb{Z}$, and let $g\in G_v\backslash \{e\}$.
Let $\varphi:G\to G$ be the isomorphism (a fold) which is the identity on $G_v$, and sends a generator $z\in G_w$ to $zg$. Then $(v,v)$ and $(w,w)$ are both exclusively $\varphi$-paired, even though $\varphi(G_w)$ is not conjugate to any vertex group.
\end{enumerate}
\end{ex}

\begin{de}[Graphical and twistless isomorphisms]\label{de:clean}
Let $\Gamma,\Lambda$ be finite simple graphs, and let $G,H$ be graph products over $\Gamma,\Lambda$, respectively. An isomorphism $\varphi:G\to H$ is  
\begin{itemize}
\item \deffont{graphical} if for every $v\in V\Gamma$, there exists $w\in V\Lambda$ such that $(v,w)$ is exclusively $\varphi$-paired;
\item \deffont{twistless} if for every $v\in V\Gamma$, the type of the parabolic support of $\varphi(G_v)$ does not split non-trivially as a join.
\end{itemize}
\end{de}

\begin{ex} Following on the previous example, we make the following observations.
\begin{enumerate}
\item If $\varphi$ sends every vertex group to a conjugate of a vertex group, then $\varphi$ is both graphical and twistless.
\item In point~\ref{item:Ex_phi_pared2} of Example~\ref{Rk:vertex-exclusively-phi-paired}, the automorphism $\varphi$ is graphical, but it is not twistless if $\twistor$ is non-trivial: in this case the parabolic support of $\varphi(G_w)$ is $G_v\times G_w$.
\item In point~\ref{item:Ex_phi_pared3} of Example~\ref{Rk:vertex-exclusively-phi-paired}, the automorphism $\varphi$ is both graphical and twistless. 
\end{enumerate}
\end{ex}

Remark that these definitions are set forth once a presentation of $G$ has been chosen. If we need to keep track of the defining graphs we will talk about $(\Gamma,\Lambda)$-graphical and $(\Gamma,\Lambda)$-twistless isomorphisms. When $\Gamma=\Lambda$, as will be the case in Part~\ref{part:automorphisms}, we will talk about $\Gamma$-graphical and $\Gamma$-twistless automorphisms.

\begin{lemma}[Blueprint for the isomorphism problem]\label{lemma:blueprint}
Let $\Gamma,\Lambda$ be finite simple graphs, and let $G,H$ be graph products over $\Gamma,\Lambda$, respectively. Let $\varphi:G\to H$ and $\psi:H\to G$ be isomorphisms, and assume that $\varphi$ and $\psi$ are graphical and twistless. Let $\theta:V\Gamma\to V\Lambda$ be such that for every $v\in V\Gamma$, the pair $(v,\theta(v))$ is exclusively $\varphi$-paired.

Then $\theta:\Gamma\to\Lambda$ is an isomorphism such that for every $v\in V\Gamma$, one has $G_v\simeq H_{\theta(v)}$.
\end{lemma}

\begin{proof}
We will show that $\theta$ is a graph isomorphism; the fact that $G_v\simeq H_{\theta(v)}$ for every $v\in V\Gamma$ is a consequence of the fact that $(v,\theta(v))$ is $\varphi$-paired.

We let $\kappa:V\Lambda\to V\Gamma$ be a map such that for every $w\in V\Lambda$, the pair $(w,\kappa(w))$ is exclusively $\psi$-paired. 

The fact that $\theta$ is a graph isomorphism is a consequence of the three points below. 

\medskip

\textbf{1 – Bijection.}  We first prove that $\theta$ is a bijection. 
\begin{itemize}
\item For injectivity, let $v_1,v_2\in V\Gamma$ be such that $\theta(v_1)=\theta(v_2)$, and denote by $w$ this common vertex. Then $(v_1,w)$ and $(v_2,w)$ are both exclusively $\varphi$-paired. This means in particular that $v_1$ is the unique vertex of $\Gamma$ such that $(w,v_1)$ is $\varphi^{-1}$-paired, and so is $v_2$. So $v_1=v_2$.
\item For surjectivity, the injectivity of $\theta$ shows that $|V\Gamma|\le |V\Lambda|$, and the injectivity of $\kappa$ shows the reverse inequality. So $|V\Gamma|=|V\Lambda|$, and as the graphs are finite $\theta$ is automatically surjective. 
\end{itemize}
\medskip

\textbf{2 – Adjacency.} We now prove that $\theta$ preserves adjacency. Let $v_1,v_2\in V\Gamma$ be adjacent. Let $Q_1,Q_2\subseteq H$ be the respective parabolic supports of $\varphi(G_{v_1})$ and $\varphi(G_{v_2})$. The definition of $\theta$ ensures that for every $i\in\{1,2\}$, the vertex $\theta(v_i)$ belongs to the type of $Q_i$. Since $G_{v_1}$ and $G_{v_2}$ commute, so do $\varphi(G_{v_1})$ and $\varphi(G_{v_2})$. Therefore $\varphi(G_{v_2})$ normalizes $Q_1$ (see Remark~\ref{Fact:Normalizer-support}), so $\varphi(G_{v_2})\subseteq Q_1\times Q_1^\perp$. As $Q_2$ is the smallest parabolic subgroup containing $\varphi(G_{v_2})$, it follows that $Q_2\subseteq Q_1\times Q_1^\perp$. Since $\varphi$ is twistless, $Q_2$ does not split non-trivially as a direct product, so either $Q_2\subseteq Q_1$ or $Q_2\subseteq Q_1^\perp$ (Lemma~\ref{lemma:parabolic-in-product}).
\begin{itemize}
\item If $Q_2\subseteq Q_1^\perp$, then $\theta(v_2)$ belongs to the type of $Q_1^\perp$, so it is adjacent to $\theta(v_1)$ and we are done. 
\item We now assume that $Q_2\subseteq Q_1$, and aim for a contradiction.
\\ -- We first note that the type of $Q_1$ cannot be reduced to one vertex $w_1$. Indeed, otherwise the type of $Q_2$ is also reduced to $w_1$. So $\varphi(G_{v_1})$ and $\varphi(G_{v_2})$ are both contained in a conjugate of $H_{w_1}$. By definition of being $\varphi$-paired, this implies that $w_1=\theta(v_1)=\theta(v_2)$, contradicting the injectivity of $\theta$.   
\\ -- Since $\varphi$ is twistless, $Q_1$ does not split non-trivially as a
direct product. Lemma~\ref{lemma:cu} therefore enables us to choose an element
$h_1\in\varphi(G_{v_1})$ whose parabolic support is equal to $Q_1$.
Using again that $Q_1$ does not split non-trivially as a direct product, the description of centralizers in a graph product (Theorem~\ref{theo:barkauskas}) ensures that $C_{Q_1}(h_1)$ is an infinite cyclic subgroup. Recalling that $\varphi(G_{v_2})$ is contained in $Q_2$, whence in $Q_1$, and commutes with $h_1$, we obtain that $\varphi(G_{v_2})$ and $h_1$ are both contained in $C_{Q_1}(h_1)$, contradicting that $\left\langle{\varphi^{-1}(h_1),G_{v_2}}\right\rangle$ is not infinite cyclic.
\end{itemize}

\medskip

\textbf{3 – Non-adjacency.} We finally prove that $\theta$ preserves non-adjacency. The above shows that $\theta$ preserves adjacency, therefore $|E\Gamma|\le |E\Lambda|$. Since $\kappa$ also preserves adjacency, the reverse equality also holds, so $|E\Gamma|=|E\Lambda|$, as desired. 
\end{proof}

In Section~\ref{sec:endgame-isomorphism}, we will prove that when the defining graphs are strongly reduced with at least three vertices,  every isomorphism $\varphi:G\to H$ can be precomposed by an automorphism $\tau\in\Aut(G)$ in order to ensure that $\varphi\circ\tau$ is graphical and twistless. In combination with Lemma~\ref{lemma:blueprint}, this will solve the isomorphism problem for graph products over strongly reduced graphs.

\section{Dealing with free products}\label{sec:free-products}

Throughout the section, we will work in the following setting.

\begin{SettingBox}
\begin{setting}\label{setting:free-product}
Let $\Gamma,\Lambda$ be finite simple graphs with at least three vertices. Let $G$ be a graph product over $\Gamma$ with vertex groups $(G_v)_{v\in V\Gamma}$, and let $H$ be a graph product over $\Lambda$ with vertex groups $(H_w)_{w\in V\Lambda}$. Let $\varphi:G\to H$ be an isomorphism. Let $P\subseteq G$ and $Q\subseteq H$ be parabolic subgroups such that $\varphi(P)=Q$. 
\end{setting}
\end{SettingBox}

We refer the reader to page~\pageref{thick} for all notions related to parabolic free factors. Our goal in this section is to prove that in Setting~\ref{setting:free-product}, and  under natural conditions on the graphs $\Gamma,\Lambda$ such as being strongly reduced, every thick free factor of $P$ is sent by $\varphi$ to a thick free factor of~$Q$, and if $v\in V\Gamma$ is a vertex for which a conjugate of $G_v$ is a thin free factor of $P$, then $v$ is exclusively $\varphi$-paired with some vertex in the type of $Q$. This is the content of Lemma~\ref{lemma:free-product-1} below, our main result in this section. The first step (showing that $\varphi$ sends thick free factors of $P$ to thick free factors of $Q$) is treated in Section~\ref{sec:thick}. Thin free factors will be dealt with in a second step, in Section~\ref{sec:thin}. Throughout the section, we will work in the case where $\Gamma,\Lambda$ are either strongly reduced, or almost strongly reduced (the latter case will be used only in Part~\ref{part:automorphisms} regarding automorphisms of graph products).

\subsection{Thick free factors are preserved}\label{sec:thick}

Our goal in this section is to prove that under Setting~\ref{setting:free-product}, and under a strong reduction assumption on the graphs $\Gamma,\Lambda$, the isomorphism $\varphi$ sends thick free factors of $P$ to thick free factors of $Q$. The key ingredient for this is the following purely algebraic characterization of thick free factors. 

\begin{lemma}[Algebraic characterization of thick free factors] 
\label{Lmm:Algebraic-characterization-thick-free}
Let $G$ be a graph product over a finite simple graph $\Gamma$, and let $P\subseteq G$ be a parabolic subgroup. Assume that $\Gamma$ is almost strongly reduced. Let $A\subseteq P$ be a subgroup, and assume that $A\neq G$. Then the following are equivalent:
\begin{enumerate}
\item\label{thick1} $A$ is a thick free factor of $P$;
\item\label{thick2} $A$ is a free factor of $P$, it is freely indecomposable, and there exists a non-trivial subgroup $B\subseteq A$ such that $N_G(B)\not\subseteq N_G(A)$.
\end{enumerate}
\end{lemma}

\begin{proof}
\ref{thick1}$\Rightarrow$\ref{thick2}. Assume that $A$ is a thick free factor of $P$. Then $A$ is a free factor of $P$, and it is freely indecomposable by \cite[Lemma~4.7]{Gre}. Let $\Delta$ be the type of $A$, so that $A=gG_{\Delta}g^{-1}$ for some $g\in G$. Since $\Gamma$ is almost strongly reduced, and $\Delta$ is a connected subgraph with $|V\Delta|\ge 2$, and not equal to $\Gamma$ because $A\neq G$, the subgraph graph $\Delta$ is not collapsible. Therefore, we can find a vertex $u\in V\Delta$ such that $\lk_{\Gamma}(u)\not\subseteq {\Delta\circ\Delta^\perp}$. Letting $B$ be a conjugate of $G_u$ contained in $A$, we then have $B\subseteq A$ and $N_G(B)\not\subseteq N_G(A)$.

\medskip

\noindent \ref{thick2}$\Rightarrow$\ref{thick1}. Let $A$ be a freely indecomposable free factor of $P$ such that there exists a non-trivial subgroup $B\subseteq A$ with $N_G(B)\not\subseteq N_G(A)$. Let $P=P_1\ast\dots\ast P_k$ be a full parabolic free product decomposition of $P$. Since $A$ is freely indecomposable, there exist $i\in\{1,\dots,k\}$ and $g\in P$ such that $A\subseteq g P_ig^{-1}$. Since $A$ is a free factor of $P$, it is also a free factor of $gP_ig^{-1}$.
\begin{itemize}
\item We first assume that $P_i$ is a thick free factor of $P$. Then $P_i$ is freely indecomposable by \cite[Lemma~4.7]{Gre}, so $A=gP_ig^{-1}$ and $A$ is a thick free factor of $P$, as desired.
\item We now assume that $P_i$ is a thin free factor of $P$ (i.e.\ its type is reduced to one vertex), and aim for a contradiction. Since $B$ is non-trivial and $B\subseteq A\subseteq gP_ig^{-1}$, the parabolic support of $B$ (and of $A$) is $gP_ig^{-1}$. By Remark~\ref{Fact:Normalizer-support}, we have 
\[N_G(B)=N_{gP_ig^{-1}}(B)\times (gP_ig^{-1})^\perp \quad\text{and}\quad N_G(A)=N_{gP_ig^{-1}}(A)\times (gP_ig^{-1})^\perp.\] 
Since $A$ is a free factor of $gP_ig^{-1}$, it is malnormal, so we have $N_{gP_ig^{-1}}(B)\subseteq A$ and $N_{gP_ig^{-1}}(A)=A$. Altogether we get $N_{G}(B)\subseteq N_G(A)$, a contradiction. \qedhere
\end{itemize}
\end{proof}

Lemma~\ref{Lmm:Algebraic-characterization-thick-free} gives a characterization of thick free factors in purely algebraic terms, i.e.\ via a condition that is preserved by group isomorphisms. We thus reach the following consequence.

\begin{cor}[Thick free factors are preserved]\label{cor:recognizing-free-factors}
Under Setting~\ref{setting:free-product}, assume that $\Gamma,\Lambda$ are almost strongly reduced.  

If $A\subseteq P$ is a thick free factor, then $\varphi(A)$ is a thick free factor of $Q$. 
\end{cor}

\begin{proof}
The case where $A\neq G$ is a consequence of Lemma~\ref{Lmm:Algebraic-characterization-thick-free}. If $A=G$, then $\varphi(A)=Q=H$. The fact that $A$ is a thick free factor implies that $\Gamma$ is connected, so $G$ is freely indecomposable \cite[Lemma~4.7]{Gre}. Therefore so is $H$, so $\Lambda$ is connected. Since $\Lambda$ is not reduced to one vertex, it follows that $\varphi(A)=Q$ is a thick free factor of itself.
\end{proof}

In a similar spirit, the following lemma highlights another situation where, under Setting~\ref{setting:free-product}, certain (possibly non-thick) parabolic free factors of $P$ are automatically sent by $\varphi$ to parabolic free factors. Its content is only needed for Part~\ref{part:automorphisms}.
 
\begin{lemma}\label{lemma:free}
Under Setting~\ref{setting:free-product}, let $A\subseteq P$ be a minimal parabolic free factor. 
    
    If for each $v,w$ in the respective types of $P,Q$, the vertex group $G_v,H_w$ is either a free group or freely indecomposable, and if $A$ is not a free group, then $\varphi(A)$ is a minimal parabolic free factor of $Q$.
\end{lemma}

\begin{proof}
Our assumption on the vertex groups $G_v,H_w$ ensures that 
any full parabolic free product decomposition of $P,Q$ is a decomposition as a finite free product of groups that are either free groups, or freely indecomposable. In particular, the non-free minimal parabolic free factors of $P$ are exactly the maximal non-trivial freely indecomposable subgroups of $P$ (and likewise for $Q$). It follows that $\varphi$ sends every minimal parabolic free factor which is not a free group to a minimal parabolic free factor. 
\end{proof}


\subsection{Dealing with thin free factors}\label{sec:thin}

Thanks to the previous section, we know (under suitable assumptions on $\Gamma,\Lambda$) that $\varphi$ sends thick free factors to thick free factors. We now analyze the images of thin free factors under $\varphi$. A crucial tool in this analysis is to take advantage of subgroups sitting outside of the free products $P,Q$ and normalizing one of the thin free factors of $P,Q$.


\begin{lemma}\label{lemma:allumette}
Under Setting~\ref{setting:free-product}, assume that $P$ and $Q$ are standard. Let $v\in V\Gamma$ be such that $G_v$ is a thin free factor of $P$, and let $w\in V\Lambda$ be a vertex in the type of the parabolic support of $\varphi(G_v)$.

Then $\varphi^{-1}(H_w)$ is contained in a conjugate of $G_v^+$.
\end{lemma}

\begin{proof}
It is enough to show that for every vertex $v'$ in the type of {the parabolic support of} $\varphi^{-1}(H_w)$, we have $v'\ge v$. So let $v'$ be such a vertex. Note that $v'$ belongs to the type of $P$; as $G_v$ is a thin free factor of $P$, it follows that $v$ and $v'$ are non-adjacent, so our goal is to prove that $\lk_\Gamma(v)\subseteq\lk_\Gamma(v')$.

Let $u\in\lk_\Gamma(v)$. Note that since $G_v$ is a free factor of $P$, the vertex $u$ does not belong to the type of $P$. Let $R$ be the parabolic support of $\varphi(G_v)$, write $R=hH_{\Upsilon}h^{-1}$ for some $h\in H$ and some full subgraph $\Upsilon\subseteq\Lambda$. Notice that $R\subseteq Q$. Then $\varphi(G_u)$ normalizes $\varphi(G_v)$, so $\varphi(G_u)$ normalizes $R$ (Remark~\ref{Fact:Normalizer-support}), so $\varphi(G_u)$ is contained in $h (H_\Upsilon\times H_{\Upsilon}^\perp)h^{-1}$. Thus
\begin{equation*}
    G_u\subseteq\varphi^{-1}\left(hH_\Upsilon h^{-1} \right) \times \varphi^{-1}\left(hH_{\Upsilon}^\perp h^{-1}\right).
\end{equation*}
Let $r_u:G\to G_u$ be the parabolic retraction and remark that $u$ does not belong to the type of $P$ and $\varphi^{-1}(hH_\Upsilon h^{-1})\subseteq P$. Therefore $\varphi^{-1}(hH_\Upsilon h^{-1})$ has trivial image under $r_u$. 
So we deduce that $G_u\subseteq r_u(\varphi^{-1}(hH_{\Upsilon}^\perp h^{-1}))$; in particular $u$ belongs to the type of the parabolic support $T$ of $\varphi^{-1}(hH_\Upsilon^\perp h^{-1})$. Let $S$ be the parabolic support of $\varphi^{-1}(hH_wh^{-1})$. Since by assumption $w$ belongs to $V\Upsilon$, the subgroup $\varphi^{-1}(hH_{\Upsilon}^\perp h^{-1})$ centralizes $\varphi^{-1}(hH_w h^{-1})$. It follows that $\varphi^{-1}(hH_{\Upsilon}^\perp h^{-1})$ normalizes $S$ 
(Remark~\ref{Fact:Normalizer-support}), and thus $\varphi^{-1}(hH_{\Upsilon}^\perp h^{-1})\subseteq S\times S^{\perp}$. Since $T$ was defined as the smallest parabolic subgroup containing $\varphi^{-1}(hH_{\Upsilon}^\perp h^{-1})$, we deduce that $T\subseteq S\times S^\perp$. Since $u$ belongs to the type of $T$, there exists $g\in G$ such that $gG_ug^{-1}\subseteq S\times S^\perp$. Since $S\subseteq P$ and $u$ does not belong to the type of $P$, necessarily $gG_ug^{-1}\subseteq S^\perp$ (Lemma~\ref{lemma:parabolic-in-product}), so $u$ belongs to the type of $S^\perp$. 
Since $v'$ lies in the type of $S$, it follows that $u\in\lk_\Gamma(v')$, as desired. 
\end{proof}

\begin{lemma}\label{lemma:free-product-1}
Under Setting~\ref{setting:free-product}, let $P=P_1\ast\dots\ast P_k$ and $Q=Q_1\ast\dots\ast Q_m$ be full parabolic free product decompositions. Assume that one of the following holds: 
\begin{enumerate}
\item $\Gamma$ and $\Lambda$ are strongly reduced;
\item $\Gamma$ and $\Lambda$ are almost strongly reduced, and the following hold: 
\begin{enumerate}
\item for any vertex $v,w$ in the respective types of $P,Q$, the vertex group $G_v,H_w$ is either free or freely indecomposable,
\item and the types of $P$ and $Q$ do not contain any edgeless subgraph on at least two vertices that is collapsible as a subgraph of $\Gamma,\Lambda$, and whose vertex groups are all free groups.
\end{enumerate}
\end{enumerate}
Then $k=m$, and up to reordering the factors $P_i$ and $Q_j$, there exists $i_0>0$ such that 
\begin{enumerate}
\item for every $i\in\{1,\dots,i_0\}$, the factors $P_i,Q_i$ are thick, and $\varphi(P_i)$ is conjugate to $Q_i$; 
\item for every $i>i_0$, the factors $P_i,Q_i$ are thin, and denoting by $v_i,w_i$ the respective types of $P_i,Q_i$, the pair $(v_i,w_i)$ is exclusively $\varphi$-paired, and one of the following assertions holds:
\begin{enumerate}
\item\label{2a} $\varphi(P_i)$ is conjugate to $Q_i$, or 
\item\label{2b} $\varphi(P_i^+\cap P)$ is conjugate to $Q_i^+\cap Q$, with $P_i^+\cap P \subseteq\langle P_1,\dots,P_i\rangle$ and $Q_i^+\cap Q\subseteq \langle Q_1,\dots,Q_i\rangle$.
\end{enumerate}
\end{enumerate}
\end{lemma}

\begin{rk}
The additional information that one of \ref{2a} or~\ref{2b} holds will not be needed to solve the isomorphism problem, but only in Part~\ref{part:automorphisms} to provide a generating set for $\Aut(G)$. A reader only interested in the isomorphism problem may therefore skip Step~3 from the present proof.

Likewise, a reader only interested in the isomorphism problem may safely assume that the graphs $\Gamma$ and $\Lambda$ are strongly reduced. The second set of assumptions, where the graphs are assumed to be almost strongly reduced, will only be used in Part~\ref{part:automorphisms}.
\end{rk}

\begin{proof}
Since pre- or post-composition of $\varphi$ by an inner automorphism does not affect the conclusion of the lemma, we may assume without loss of generality that $P$ and $Q$ are standard parabolic subgroups (and we will likewise assume that the subgroups $P_i$ and $Q_j$ are standard).

\medskip

\underline{\textbf{Definition of $i_0$ and $i_1$ and general setup.}} Corollary~\ref{cor:recognizing-free-factors} gives us $i_0$ such that, up to permuting the factors,
\begin{itemize}
\item $P_1,\dots,P_{i_0}$ are thick free factors of $P$, and $Q_1,\dots,Q_{i_0}$ are thick free factors of $Q$, with $\varphi(P_i)$ conjugate to $Q_i$ for all $i\le i_0$; 
\item the parabolic subgroups $P_i$ with $i>i_0$ and $Q_j$ with $j>i_0$ are thin. 
\end{itemize}
We now define $i_1\ge i_0$ in the following way. If $\Gamma$ and $\Lambda$ are strongly reduced, we simply let $i_1=i_0$. Otherwise, by assumption, all vertex groups $G_v,G_w$ are either free or freely indecomposable, so up to permuting the factors again, Lemma~\ref{lemma:free} provides us with $i_1\ge i_0$ such that $P_{i_0+1},\dots,P_{i_1}$ and $Q_{i_0+1},\dots,Q_{i_1}$ are non-free, while $P_{i_1+1},\dots,P_k$ and $Q_{i_1+1},\dots,Q_m$ are free, and $\varphi(P_i)$ is conjugate to $Q_i$ for every $i\le i_1$.

Notice that 
\begin{equation}\label{inequivalent}\tag{$\star$}
\text{for $i,i'>i_1$,} \begin{cases} \text{the types of $P_i$ and $P_{i'}$ are inequivalent vertices of $\Gamma$;} \\ \text{the types of $Q_i$ and $Q_{i'}$ are inequivalent vertices of $\Lambda$.}\end{cases}
\end{equation}
Indeed, this follows from our assumption that $\Gamma,\Lambda$ are strongly reduced (with at least three vertices), or that the types of $P$ and $Q$ do not contain any collapsible edgeless subgraph on at least two vertices whose vertex groups are all free groups. We can therefore reorder the factors so that, denoting by $v_i$ the vertex giving the type of $P_i$, the following holds: 
\begin{equation}\label{ordering}\tag{$\star\star$}
\text{for $i_3>i_2>i_1$, we have $v_{i_2}\not\le v_{i_3}$.} 
\end{equation}

\medskip

\underline{\textbf{Step 1.}} We will construct a bijection $\sigma:\{1,\dots,k\}\to\{1,\dots,m\}$ such that for every $i\in\{1,\dots,k\}$, the following three properties hold:  
\begin{enumerate}  
    \item\label{point-2} if $i>i_1$, then $\varphi^{-1}(Q_{\sigma(i)})$ is contained in a conjugate of $P_i^+$,\\ 
    and $P_i^+\cap P \subseteq\langle P_1,\dots,P_i\rangle$; 
    \item\label{point-1} if $i>i_1$, then $\varphi(P_i)$ is contained in a conjugate of $Q_{\sigma(i)}^+$,\\
    and $Q_{\sigma(i)}^+\cap Q\subseteq \langle Q_{\sigma(1)},\dots,Q_{\sigma(i)}\rangle$;
    \item\label{point-3} $\varphi(\llangle P_1,\dots,P_i\rrangle)=\llangle %\llangle
    Q_{\sigma(1)},\dots,Q_{\sigma(i)}\rrangle$; more precisely, letting $M_i=\llangle P_1,\dots,P_{i-1}\rrangle$ and $N_i=\llangle Q_{\sigma(1)},\dots,Q_{\sigma(i-1)}\rrangle$, the subgroup $\varphi(\langle M_i,P_i\rangle)$ {is conjugate to} $\langle N_i,Q_{\sigma(i)}\rangle$.
\end{enumerate}
Here the normal closures are understood in $G$ and $H$ (but the third item is also valid, with the same proof, if they are understood in $P$ and $Q$).

\smallskip

By induction on $i\in\{1,\dots,k\}$, we will construct an injective map $\sigma:\{1,\dots,k\}\to\{1,\dots,m\}$ satisfying items~\ref{point-2}, \ref{point-1}, \ref{point-3}. 
For $i\le i_1$, we simply let $\sigma(i)=i$; the first two assertions are void, and the third one follows from our definition of $i_1$. We now let $i>i_1$, and construct $\sigma$ inductively. So assume that $\sigma(1),\dots,\sigma(i-1)$ have been defined, and let us define $\sigma(i)$. For notational simplicity, up to reordering the factors $Q_\ell$, we will assume that $\sigma(j)=j$ for every $j<i$. 

Write $P_i=G_v$. Let $w$\label{def:sommet-w-Qi} belong to the type of the parabolic support of $\varphi(G_v)$ but not to the type of any $Q_j$ with $j<i$. This exists because otherwise $\varphi(G_v)$ would be contained 
in a conjugate of $\langle Q_1,\dots,Q_{i-1}\rangle$ contradicting point~\ref{point-3} of our inductive hypothesis.  
We reorder the free factors $\{Q_j\}_{j\geq i}$ such that $Q_i=H_w$, and set $\sigma(i)=i$. We will now prove that $Q_{\sigma(i)}=H_w$ satisfies the three items from Step~1.

\smallskip

\textbf{$\drsh$ Item~\ref{point-2}.} Lemma~\ref{lemma:allumette} ensures that $\varphi^{-1}(H_w)$ is {contained in a conjugate of} $G^+_v=P^+_i$. In addition $P_i^+\cap P\subseteq \langle P_1,\dots,P_i\rangle$ {by condition \eqref{ordering}}. So item~\ref{point-2} is established.

\smallskip

\textbf{$\drsh$ Item~\ref{point-1}.} We first observe that $v$ belongs to {the type of the parabolic} support of $\varphi^{-1}(H_w)$. Indeed, item \ref{point-3} of our inductive hypothesis ensures that this {type} contains a vertex $z$ which is not in the type of any $P_j$ with $j<i$. Since $\varphi^{-1}(H_w)$ is contained both in $P$ and in a conjugate of $P_i^+$ (item~\ref{point-2}), it follows from \cite[Proposition~3.4]{AM} that it is contained in a conjugate of $P_i^+\cap P$. Using item~\ref{point-2} again, we deduce that $\varphi^{-1}(H_w)$ is contained in a conjugate of $\langle P_1,\dots,P_i\rangle$. So necessarily $z=v$.

We can therefore apply Lemma~\ref{lemma:allumette} to the isomorphism $\varphi^{-1}$ instead of $\varphi$, and we obtain that $\varphi(G_v)$ is contained in a conjugate of $H_w^+$.

To conclude the proof of item~\ref{point-1}, there remains to prove that $Q_i^+\cap Q\subseteq\langle Q_1,\dots,Q_i\rangle$. So assume towards a contradiction that this is not the case. Then there exists $j>i$ such that the thin free factor $Q_j$ is contained in $Q_i^+$. Let $w'\in V\Lambda$ be such that $Q_j=H_{w'}$. Since $Q_j\subseteq Q_i^+$, we have $w'\ge w$. Since the vertices $w$ and $w'$ are inequivalent (Fact~\eqref{inequivalent}), in fact $w'>w$. Let $v'\in V\Gamma$ be a vertex in the type of the parabolic support $R$ of $\varphi^{-1}(H_{w'})$, chosen so that $G_{v'}=P_{i'}$ for some $i'\ge i$ -- this exists by item~\ref{point-3} of our inductive assumption. Lemma~\ref{lemma:allumette} ensures that $\varphi(G_{v'})$ is contained in a conjugate of $H_{w'}^+$. {But} $w$ belongs to the type of the parabolic support of $\varphi(G_v)$ and $w<w'$, so we have $v'\neq v$. Condition~\eqref{ordering} ensures that $v\not\leq v'$, so we can find a vertex $u\in V\Gamma$ which is adjacent to $v$ but not to $v'$. Note that since $v$ is isolated in the type of $P$, the vertex $u$ does not belong to the type of $P$. Then for all $x$ in the type of $P$ such that $x\geq v$ we have $u\in \lk_{\Gamma}(x)$ and thus $G_u$ commutes with $G_x$. In particular $G_u$ commutes with $G^+_v\cap P$. But, by item~{\ref{point-2}} we have $\varphi^{-1}(H_w)\subseteq g'G_v^+{g'}^{-1}\cap P$ for some $g'\in G$. By \cite[Proposition~3.4]{AM} there exists $g\in P$ such that $g'G_v^+{g'}^{-1}\cap P\subseteq g(G^+_v\cap P) g^{-1}$. Therefore $H_w$ commutes with $\varphi(gG_ug^{-1})$ and thus $\varphi(gG_ug^{-1})\subseteq H_w\times H_w^\perp$.

Let $k\in gG_ug^{-1}\setminus\{e\}$. Then there exists $t\in\varphi^{-1}(H_w)$ such that $\varphi(kt)\in H_w^{\perp}$. Since $w<w'$, we have $H_w^\perp\subset H_{w'}^\perp$, so $\varphi(kt)$ commutes with $H_{w'}$. Therefore $kt$ commutes with $\varphi^{-1}(H_{w'})$, and thus $kt$ normalizes the parabolic support $R$ of $\varphi^{-1}(H_{w'})$ (Remark~\ref{Fact:Normalizer-support}). In particular $kt\in R\times R^\perp$. But recall that $u\notin \lk_\Gamma(v')$ and that $v'$ belongs to the type of $R$, so $u$ does not belong to the type of $R^\perp$. Moreover $R\subseteq P$ and $u$ does not belong to the type of $P$, so $u$ does not belong to the type of $R$ either. Altogether, $u$ does not belong to the type of $R\times R^{\perp}$. On the other hand $k\in gG_ug^{-1}$ and $t\in P$, so by retracting to $G_u$ we obtain that $u$ belongs to the type of the support of $kt$. This is the desired contradiction, which proves the wanted assertion. 

\smallskip

\textbf{$\drsh$ Item~\ref{point-3}.}  
By our inductive hypothesis, we have $\varphi(M_i)=N_i$. Item~\ref{point-1} ensures that $\varphi(P_i)$ is contained in a conjugate of $Q^+_i$. Applying \cite[Proposition 3.4]{AM} shows that there exists $h\in H$ such that
\begin{equation*}
    \varphi\left(P_i\right)\subseteq h\left(Q_i^+\cap Q\right)h^{-1}\subseteq h\left\langle{N_i,Q_i}\right\rangle h^{-1}.
\end{equation*}
It follows that $\varphi(\langle M_i,P_i\rangle)\subseteq h \langle N_i,Q_i\rangle h^{-1}$, and therefore $\varphi(M_{i+1})\subseteq N_{i+1}$. The same argument applied to $\varphi^{-1}$ (using item~\ref{point-2} in place of item~\ref{point-1}) shows that $\varphi^{-1}(\langle N_i,Q_i\rangle)$ is contained in $g\langle M_i,P_i\rangle g^{-1}$ for some $g\in G$. Altogether
\begin{equation*}
    \left\langle {M_i,P_i}\right\rangle
    \subseteq \varphi^{-1}(h)\varphi^{-1}\left(\left\langle N_i,Q_i\right\rangle\right)\varphi^{-1}(h)^{-1}\subseteq \varphi^{-1}(h)g \left\langle M_i,P_i\right\rangle (\varphi^{-1}(h)g)^{-1}.
\end{equation*}
We now claim that these inclusions are in fact equalities, and for this it is enough to show that if $\langle M_i,P_i\rangle\subseteq s\langle M_i,P_i\rangle s^{-1}$ for some $s\in G$, then this inclusion is an equality. In this case, using the normality of $M_i$, every $p\in P_i$ can be written as $m_psp's^{-1}$ with $m_p\in M_i$ and $p'\in P_i$, and by retracting to $P_i$ we get $p=r_{P_i}(s)p'r_{P_i}(s)^{-1}$, so the map $\kappa:p\mapsto p'$ is uniquely well defined and bijective (and $m_p$ is also uniquely well defined). So conversely, for every element of the form $msp's^{-1}$ with $m\in M_i$ and $p'\in P_i$, we can find $p\in P_i$ such that $p'=\kappa(p)$, and rewrite $msp's^{-1}=mm_p^{-1}(m_ps\kappa(p)s^{-1})=mm_p^{-1}p$, showing that $msp's^{-1}\in \langle M_i,P_i\rangle$. This proves the surjectivity of the above inclusion.

This shows that $\varphi(\langle M_i,P_i\rangle)=h\langle N_i,Q_i\rangle h^{-1}$. In particular $\varphi(M_{i+1})=N_{i+1}$.

\smallskip

\textbf{$\drsh$ Conclusion of Step~1.} Our inductive argument has enabled us to build an injective map $\sigma:\{1,\dots,k\}\to\{1,\dots,m\}$ satisfying the three requirements from Step~1. In particular $k\le m$. Arguing with $\varphi^{-1}$ shows that the reverse inequality also holds, so in fact $k=m$ and $\sigma$ is a bijection. This concludes Step~1.
\medskip

\underline{\textbf{Notational interlude.}} Up to permuting the factors $Q_i$, from now on we will assume that $\sigma(i)=i$ for every $i\in\{1,\dots,k\}$.
\medskip 

\underline{\textbf{Step 2.}} For $i>i_0$, let $v_i,w_i$ be the respective types of $P_i,Q_i$. We show that $(v_i,w_i)$ is exclusively $\varphi$-paired. 
\smallskip

\textbf{$\drsh$ 1 – Pairing.} We first show that $(v_i,w_i)$ is $\varphi$-paired. By Step 1, we have $\varphi(\langle M_i,P_i\rangle)=h\langle N_i,Q_i\rangle h^{-1}$ for some $h\in H$, while $\varphi(M_i)=N_i$. Note that $M_i$ is the kernel of the parabolic retraction $\langle M_i,P_i\rangle\to P_i$, which thus descends at the quotient to an isomorphism $\langle M_i,P_i\rangle/M_i\to P_i$. Likewise, the parabolic retraction $\langle N_i,Q_i\rangle\to Q_i$ induces an isomorphism $\langle N_i,Q_i\rangle/N_i\to Q_i$. Through these identifications, $\varphi$ descends at the quotient to an isomorphism $\overline\varphi:P_i\to Q_i$. We deduce that $(r_{Q_i}\circ\varphi)_{|P_i}:P_i\to Q_i$ is an isomorphism. The same argument applied to $\varphi^{-1}$ shows that $(r_{P_i}\circ\varphi^{-1})_{|Q_i}:Q_i\to P_i$ is an isomorphism.

\smallskip 

\textbf{$\drsh$ 2 – Exclusivity.} We claim that $(v_i,w_i)$ is exclusively $\varphi$-paired. {For this, we prove that $w_i$ is the unique $w\in V\Lambda$ such that $(v_i,w)$ is $\varphi$-paired. The fact that $v_i$ is the unique $v\in V\Gamma$ such that $(w_i,v)$ is $\varphi^{-1}$-paired is symmetric, arguing with $\varphi^{-1}$ instead of $\varphi$.}
\begin{itemize}
\item For $j<i$, and $w\in V\Lambda$ such that $H_w\subseteq Q_j$, we have $\varphi^{-1}(H_w)\subseteq \llangle P_1,\dots,P_j\rrangle$ (item~\ref{point-2} of Step 1). So $(r_{P_i}\circ\varphi^{-1})_{|H_w}$ is trivial and hence $(v_i,w)$ is not $\varphi$-paired;
\item for $j>i$ and $w\in V\Lambda$ such that $H_w=Q_j$, we have $\varphi(P_i)\subseteq \llangle Q_1,\dots,Q_i\rrangle$ (item~\ref{point-1} of Step 1). So $(r_{Q_j}\circ\varphi)_{|P_i}$ is trivial,  
and thus $(v_i,w)$ is not $\varphi$-paired;
\item likewise, if $w$ is a vertex not in the type of $Q$, then $r_w\circ\varphi$ is trivial in restriction to $P_i$, so $(v_i,w)$ is not $\varphi$-paired.
\end{itemize}
\medskip

\underline{\textbf{Step 3.}} Let $i>i_0$ and let us show that one of \ref{2a} or~\ref{2b} holds.
\smallskip

%We 
First assume that $P_i^+\cap P=P_i$. Then by the first item of Step~1, $\varphi^{-1}(Q_i)$ is contained in a conjugate of $P_i$. In addition, Step~2 ensures that $(r_{P_i}\circ\varphi^{-1})_{|Q_i}:Q_i\to P_i$ is an isomorphism.
Hence $\varphi^{-1}(Q_i)$ is conjugate to $P_i$, so $\varphi(P_i)$ is conjugate to~$Q_i$. 

Likewise, if $Q_i^+\cap Q=Q_i$, then the same reasoning with the roles of $\varphi$ and $\varphi^{-1}$ reversed shows that $\varphi(P_i)$ is conjugate to $Q_i$.

We now assume that $\varphi(P_i)$ is not contained in a conjugate of $Q_i$, and that $\varphi^{-1}(Q_i)$ is not contained in a conjugate of $P_i$; in particular, by the above, $P_i\subsetneq P_i^+\cap P$ and $Q_i\subsetneq Q_i^+\cap Q$. 
\begin{itemize}
\item We first show that $\varphi(P_i^\perp)$ is contained in a conjugate of $Q_i^\perp$.\\
Let $R\subseteq Q$ be the parabolic support of $\varphi(P_i)$. Recall from Step~2 that, denoting by $v_i,w_i$ the respective types of $P_i,Q_i$, the pair $(v_i,w_i)$ is $\varphi$-paired. Hence $w_i$ belongs to the type of $R$, so there exists $h\in Q$ such that $hQ_ih^{-1}\subseteq R$. We also observe that this inclusion is strict, because we assumed that $\varphi(P_i)$ is not contained in a conjugate of $Q_i$. Since the type of $R$ contains $w_i$, which is an isolated vertex in the type of $Q$, it follows that $R$ splits non-trivially as a free product. Since $\varphi(P_i^\perp)$ commutes with $\varphi(P_i)$, it normalizes $R$ (Remark~\ref{Fact:Normalizer-support}), so 
\begin{equation*}
    \varphi\left(P_i^\perp\right)\subseteq R\times R^{\perp}.
\end{equation*} 
Now let $u\in V\Gamma$ be such that $G_{u}\subseteq P_i^+\cap \langle P_1,\dots,P_{i-1}\rangle$. Then $G_u$ commutes with $P_i^\perp$. So 
$\varphi(P_i^\perp)$ commutes with $\varphi(G_{u})$, 
and hence $\varphi(P_i^{\perp})$ centralizes $\langle \varphi(P_i),\varphi(G_u)\rangle$. On the other hand, by the description of centralizers in free products, the centralizer of $\langle \varphi(P_i),\varphi(G_u)\rangle$ in $R$ is trivial because
\begin{itemize}
    \item $\varphi(P_i)$ is not contained in any freely indecomposable free factor of $R$ (otherwise its support would be a proper parabolic subgroup of $R$),
    \item and $\langle \varphi(P_i),\varphi(G_u)\rangle$ is not cyclic.
\end{itemize}
It follows that the projection of $\varphi(P_i^\perp)$ to $R$ is trivial. Hence $\varphi(P_i^\perp)\subseteq R^\perp\subseteq hQ_i^\perp h^{-1}$.
 \item The same reasoning applied to $\varphi^{-1}$ shows that there exists $g\in G$ such that $\varphi^{-1}(Q_i^\perp)\subseteq g P_i^\perp g^{-1}$. 
\end{itemize} 
Altogether, we have thus found $a,b\in G$ such that $P_i^\perp\subseteq a\varphi^{-1}(Q_i^\perp) a^{-1}\subseteq b P_i^\perp b^{-1}$, and it follows from \cite[Lemma~3.9]{AM} that these inclusions are equalities. So 
\begin{equation*}
    \varphi\left(P_i^\perp\right)=\varphi(a)Q_i^\perp\varphi(a)^{-1}.
\end{equation*}
Finally, since $P_i^+\cap P$ commutes with $P_i^\perp$, the above implies that $\varphi(P_i^+\cap P)$ commutes with $\varphi(a)Q_i^{\perp}\varphi(a)^{-1}$, so $\varphi(P_i^+\cap P)\subseteq \varphi(a)Q_i^+\varphi(a)^{-1}$. Since additionally $\varphi(P_i^+\cap P)\subseteq Q$, it follows from \cite[Proposition~3.4]{AM} that $\varphi(P_i^+\cap P)$ is contained in a conjugate of $Q_i^+\cap Q$. Reverting the roles of $\varphi$ and $\varphi^{-1}$ also shows that $\varphi^{-1}(Q_i^+\cap Q)$ is contained in a conjugate of $P_i^+\cap P$, so it follows that $\varphi(P_i^+\cap P)$ is conjugate to $Q_i^+\cap Q$. 
\end{proof}


\section{Dealing with direct products}\label{sec:direct-products}

The goal of this section is to prove Lemma~\ref{lemma:direct-product} below, which says that (under conditions on the underlying graphs, such as being strongly reduced) if an isomorphism $\varphi$ between two graph products $G$ and $H$ sends a parabolic subgroup $P\subseteq G$ to a parabolic subgroup $Q\subseteq H$, then up to precomposing $\varphi$ by a product of twists, it sends every minimal parabolic direct factor of $P$ to a minimal parabolic direct factor of $Q$. This is standard when the clique factors of $P$ and $Q$ are trivial (see e.g.\ \cite[Lemma~3.5]{Fio}) because then all factors are directly indecomposable and have trivial center; there is no need to precompose by twists in this case. In general, we will first use \cite[Appendix~A]{EH} to show that $\varphi$ sends the clique factor of $P$ to the clique factor of $Q$. We will then exploit normalization within the ambient graph product to precompose $\varphi$ by a product of twists and ensure that it sends every vertex group of the clique factor of $P$, to a vertex group of the clique factor of $Q$, before finally untwisting $\varphi$ on the whole of $P$. 

In order to carry out this strategy, we start with a few useful observations.

\begin{lemma}\label{Lmm:detwister-un-facteur-direct}
Let $P_1,R_1,P_2,R_2$ be groups. Let $\varphi:P_1\times R_1\to P_2\times R_2$ be an isomorphism such that $\varphi(R_1)=R_2$. Then there exists a homomorphism $\twistor:P_1\to Z(R_1)$ such that, letting $\tau\in\Aut(P_1\times R_1)$ be the automorphism satisfying $\tau(g)=g$ for every $g\in R_1$, and $\tau(g)=g\twistor(g)$ for every $g\in P_1$, one has $\varphi\circ\tau(P_1)=P_2$.
\end{lemma}

\begin{proof} Since $\varphi(R_1)=R_2$, it induces an isomorphism $\bar\varphi:(P_1\times R_1)/R_1\to (P_2\times R_2)/R_2$, which we view as an isomorphism $\beta:P_1\to P_2$. It means that there exists a homomorphism $\twistor':P_1\to R_1$ such that for every $g\in P_1$, one has $\varphi(g)=\beta(g)\varphi(\twistor'(g))$. But
\begin{itemize}
    \item $g$ commutes with $R_1$, so $\varphi(g)$ commutes with $\varphi(R_1)=R_2$;
    \item and $\beta(g)\in P_2$ which commutes with $R_2$;
\end{itemize}
so $\varphi(\twistor'(g))$ commutes with $R_2$, and hence $\twistor'(g)$ commutes with $\varphi^{-1}(R_2)=R_1$. Therefore $\twistor'(g)\in Z(R_1)$, and the homomorphism $\twistor$ defined for all $g\in P_1$ by $\twistor(g)=(\twistor'(g))^{-1}$ satisfies the conclusion of the lemma.
\end{proof}

Often, when changing an automorphism $\varphi$ by precomposing it by a twist to tighten it on a certain vertex group, it will be important to ensure that this operation does not perturb the image of other parabolic subgroups under $\varphi$. This is the content of the following lemma.

\begin{lemma}\label{lemma:redress}
Let $\Gamma,\Lambda$ be finite simple graphs. Let $G$ be a graph product over $\Gamma$ with vertex groups $(G_v)_{v\in V\Gamma}$, and let $H$ be a graph product over $\Lambda$ with vertex groups $(H_w)_{w\in V\Lambda}$. Let $v\in V\Gamma$, and $A\subseteq G$ and $B,K\subseteq H$ be parabolic subgroups such that $A\subseteq G_{\lk(v)}\cap G_v^+$ and $B\subseteq K^\perp$. Let $\varphi:G\to H$ be an isomorphism, and assume that $\varphi(\langle A,G_v\rangle)\subseteq \langle B,K\rangle$ and $\varphi(A)=B$. Let $\twistor:G_v\to Z(A)$ be a homomorphism, and consider the automorphism $\tau$ of $G$ defined by letting $\tau(g)=g\twistor(g)$ for $g\in G_v$, and $\tau(g)=g$ for $g\in G_{\Gamma\setminus\{v\}}$. Assume that $\varphi\circ\tau(G_v)\subseteq K$.

Then for every standard parabolic subgroup $P\subseteq G$, and every parabolic subgroup $Q\subseteq H$, if $\varphi(P)=Q$, then $\varphi\circ\tau(P)=Q$. 
\end{lemma}

\begin{rk}\label{Rk:redress}
It follows that if $P$ is a (possibly non-standard) parabolic subgroup of $G$ such that $\varphi(P)$ is a parabolic subgroup $Q$, then $\varphi\circ\tau(P)$ is conjugate to $Q$.
\end{rk}

\begin{proof}
The conclusion is clear if $G_v \not\subseteq P$, as $\tau_{|P}=\mathrm{id}_P$ in this case. We thus assume that $G_v\subseteq P$ and aim at showing that $\tau(P)=P$, from which the lemma follows.

Since $G_v\subseteq P$, we have $\varphi(G_v)\subseteq Q$. Let $B'\subseteq B$ be the parabolic support of $\varphi\circ\twistor(G_v)$. We claim that $B'\subseteq Q$. Indeed, consider a full parabolic direct product decomposition $B'=B'_1\times\dots\times B'_k$. We will prove that for every $i\in\{1,\dots,k\}$, we have $B'_i\subseteq Q$, from which the above claim follows.

So let $i\in\{1,\dots,k\}$. Lemma~\ref{lemma:cu} enables us to choose an element $h\in\varphi\circ\twistor(G_v)$ whose parabolic support contains $B'_i$. Let $g\in G_v$ be such that $h=\varphi\circ\twistor(g)$. Then 
\begin{equation*}
    \varphi\circ\tau(g)=\varphi(g) \varphi\circ\twistor(g)=\varphi(g)h,
\end{equation*}
and this element belongs to $K$ by assumption on $\varphi\circ\tau(G_v)$. 
As $\varphi(g)=(\varphi\circ\tau(g))h^{-1}$, it follows that $\varphi(g)$
belongs to $\langle K,B\rangle$, which is isomorphic to $K\times B$ as
$B\subseteq K^\perp$. The parabolic support of $\varphi(g)$ is a parablic
subgroup of $K\times B$ that contains $B'_i$.
Since $G_v\subseteq P$ and $\varphi(P)=Q$, we have $\varphi(g)\in Q$, and therefore $B'_i\subseteq Q$, as desired. This proves our claim that $B'\subseteq Q$.

It follows from the above claim that $\varphi^{-1}(B')\subseteq P$. So $\twistor(G_v)\subseteq P$, and hence $\tau(P)=P$, as desired.
\end{proof}

Our next lemma deals with the persistence of the twistless property (Definition~\ref{de:clean}) when collapsing an almost strongly reduced graph to a strongly reduced graph. Its content is not needed for our solution to the isomorphism problem for graph products, but only for the work on automorphisms of graph products in Part~\ref{part:automorphisms}.

\begin{lemma}\label{lemma:twistless-collapse}
Let $\Gamma,\Lambda$ be finite simple graphs which are almost strongly reduced, and let $\overline{\Gamma},\overline{\Lambda}$ be strong reductions of $\Gamma,\Lambda$. Let $G$ and $H$ be graph products over $\Gamma,\Lambda$, and let $\varphi:G\to H$ be an isomorphism.

If $\varphi$ is $(\Gamma,\Lambda)$-twistless, then $\varphi$ is $(\overline{\Gamma},\overline{\Lambda})$-twistless.
\end{lemma}

\begin{proof}
Let $\pi_\Gamma:\Gamma\to\overline{\Gamma}$ and $\pi_{\Lambda}:\Lambda\to\overline{\Lambda}$ be the collapse maps. Let $\bar v\in V\overline{\Gamma}$, and let $\Delta_{\bar v}:=\pi_\Gamma^{-1}(\bar v)$. As $\Gamma$ is almost strongly reduced, $\Delta_{\bar v}$ is edgeless. Let $\Upsilon_R$ be the type of the parabolic support $R$ of $\varphi(G_{\Delta_{\bar v}})$ for the graph product structure of $H$ over $\Lambda$.

\medskip

\textbf{Observation 1:} $\Upsilon_R$ does not split non-trivially as a join, and neither does $\pi_\Lambda(\Upsilon_R)$. 

The statement about $\pi_\Lambda(\Upsilon_R)$ follows from the one about $\Upsilon_R$, so we focus on $\Upsilon_R$.
\begin{itemize}
\item If $\Delta_{\bar v}$ is reduced to one vertex, this follows from the fact that $\varphi$ is $(\Gamma,\Lambda)$-twistless. 
\item Otherwise, enumerate $V\Delta_{\bar v}=\{v_1,\dots,v_n\}$, and assume towards a contradiction that $\Upsilon_R$ splits non-trivially as a join, giving a minimal parabolic direct product decomposition $R=R_1\times\dots\times R_k$. Since $\varphi$ is $(\Gamma,\Lambda)$-twistless, for every $i\in\{1,\dots,n\}$, there exists $\sigma(i)\in\{1,\dots,n\}$ such that the parabolic support of $\varphi(G_{v_i})$ is contained in $R_i$. As the subgroups $G_{v_1},\dots,G_{v_k}$ span a free product (and thus do not commute), one has $\sigma(1)=\dots=\sigma(k)$. Hence $\varphi(G_{\Delta_{\bar v}})\subseteq R_{\sigma(1)}$, contradicting the fact that $R$ is the smallest parabolic subgroup of $H$ that contains $\varphi(G_{\Delta_{\bar v}})$.   
\end{itemize} 

\medskip

\textbf{Observation 2:} $\pi_\Lambda(\Upsilon_R)$ is the type of the parabolic support of $\varphi(G_{\bar v})=\varphi(G_{\Delta_{\bar v}})$ for the graph product structure of $H$ over $\overline{\Lambda}$.

Indeed, letting $R=hH_{\Upsilon_R}h^{-1}$ for some $h\in H$,
\begin{itemize}
\item we have $\varphi(G_{\bar v})\subseteq h H_{\pi_{\Lambda}(\Upsilon_R)}h^{-1}$;
\item on the other hand, if $R'\subsetneq h H_{\pi_{\Lambda}(\Upsilon_R)}h^{-1}$ is a proper parabolic subgroup for the graph product structure of $H$ over $\overline{\Lambda}$, then its type $\Upsilon'$ is a proper subgraph of $\pi_{\Lambda}(\Upsilon_R)$. Using $\pi_\Lambda^{-1}$, we deduce that $R'$ is a parabolic subgroup for the graph product structure on $\Lambda$ whose type is a proper subgraph of $\Upsilon_R$. Since $R$ is the support of $\varphi(G_{\bar v})$ for the graph product structure on $\Lambda$, it follows that $\varphi(G_{\bar v})\not\subseteq R'$.
\end{itemize}

\medskip

Combining the above two observations shows that the type of the parabolic support of $\varphi(G_{\bar v})$ for the parabolic structure on $\overline{\Lambda}$ is $\pi_\Lambda(\Upsilon_R)$, and it does not split as a join. This proves that $\varphi$ is $(\overline{\Gamma},\overline{\Lambda})$-twistless.
\end{proof}

Our next lemma is the main result of the present section. We refer to page~\pageref{de:saturated} for the notion of a saturated parabolic subgroup. Recall from the introduction (Definition~\ref{de:basic}) that $\varphi\in \Aut(G)$ is called a \deffont{twist} if there exist distinct vertices $v,v'\in V\Gamma$ such that $\st_\Gamma(v)\subseteq\st_\Gamma(v')$, and a homomorphism $\twistor:G_v\to Z(G_{v'})$, such that $\varphi(g)=g\twistor(g)$ for every $g\in G_v$, and $\varphi(h)=h$ for every $h\in G_{\Gamma\setminus\{v\}}$. In this case we refer to $G_v$ as the \deffont{twisted subgroup} and to $\twistor(g)$ as the \deffont{twistor} of $g$. \label{def:twistor}

\begin{lemma}\label{lemma:direct-product}
Let $\Gamma,\Lambda$ be finite simple graphs which are not reduced to one edge. Let $G$ be a graph product over $\Gamma$ with vertex groups $(G_v)_{v\in V\Gamma}$, and let $H$ be a graph product over $\Lambda$ with vertex groups $(H_w)_{w\in V\Lambda}$. Let $\varphi:G\to H$ be an isomorphism. Let $P,Q\subseteq G$ be two standard
parabolic subgroups such that $\varphi(P)=Q$. Assume that $P$ and $Q$ split as direct products of at least two non-trivial parabolic subgroups, and let $P=P_1\times\dots\times P_k$ and $Q=Q_1\times\dots\times Q_m$ be their respective full parabolic direct product decompositions. 

Assume that one of the following holds: 
\begin{itemize}
\item $\Gamma$ and $\Lambda$ are strongly reduced;
     \item $\Gamma$ and $\Lambda$ are almost strongly reduced, no vertex group $G_v,H_w$ splits non-trivially as a free product, and $P,Q$ are saturated.
\end{itemize}

Then $k=m$, and there exist a bijection $\sigma:\{1,\dots,k\}\to\{1,\dots,m\}$ and an automorphism $\tau\in\Aut(G)$ which is a product of twists, such that 
\begin{itemize}
\item for every $i\in\{1,\dots,k\}$, one has $\varphi\circ\tau(P_i)=Q_{\sigma(i)}$;
\item for every parabolic subgroup $P'\subseteq G$, if $\varphi(P')$ is a parabolic subgroup $Q'$, then $\varphi\circ\tau(P')$ is conjugate to $Q'$;
\item if $\varphi$ is $(\Gamma,\Lambda)$-twistless, then $\tau=\mathrm{id}$.
\end{itemize}
\end{lemma}

\begin{rk}
Only the case where $\Gamma$ and $\Lambda$ are strongly reduced is needed for our solution of the isomorphism problem. The case where they are almost strongly reduced is only needed for Part~\ref{part:automorphisms} of this work on automorphisms.
\end{rk}

\begin{proof} We can assume that $\Gamma,\Lambda$ are not reduced to one vertex, otherwise there is no subgroup $P$ (or $Q$) as in the statement.

If $\Gamma$ is disconnected, then $\Lambda$ must be disconnected as well, as this is equivalent to $G,H$ splitting non-trivially as a free product. If this happens, as $\Gamma$ and $\Lambda$ are almost strongly reduced, they must be edgeless. In this case there are no subgroups $P,Q$ as in the statement of the lemma, so the lemma is vacuously true.

We therefore assume that $\Gamma,\Lambda$ are connected. Since they are almost strongly reduced and not reduced to one edge, they cannot split non-trivially as a join. Let $\overline{\Gamma},\overline{\Lambda}$ be their respective strong reductions (which are unique, see Remark~\ref{Rq:strong-reduction-as-at-least-3-vertices}). Note that in the case where $\Gamma,\Lambda$ are assumed strongly reduced, then $\overline{\Gamma}=\Gamma$ and $\overline{\Lambda}=\Lambda$.

Up to permuting the factors $P_i,Q_i$, we can assume that there exist $i_0,j_0,i_1,j_1$ such that $C_P=P_1\times \dots \times P_{i_0}$ and $C_Q=Q_1 \times \dots \times Q_{j_0}$ are the respective clique factors of $P$ and $Q$, and $C^e_P=P_1\times \dots \times P_{i_1}$ and $C^e_Q=Q_1 \times \dots \times Q_{j_1}$ are their respective extended clique factors, as defined on page~\pageref{de:extended-clique-factor}. 
  
   Since $P,Q$ are saturated, they remain parabolic subgroups for the new graph product structures on $\overline{\Gamma}$ and $\overline{\Lambda}$ respectively. They inherit new full parabolic direct product decompositions, that we now describe. To this end, notice that a subgraph $\Delta$ of $\Gamma$ in the type of $P$ can be collapsed to a point only if it is totally disconnected, which forces it to be contained in the type $\Gamma_i$ of some $P_i$. The image of $\Delta$ under the collapse map is then reduced to one point if and only if $P_i$ was contained in the extended clique factor of $P$. It follows that the clique factor of $P$ for the graph product structure on $\overline{\Gamma}$ is $C_P^e$, and its full parabolic direct product decomposition becomes $P=C_P^e\times P_{i_1+1}\times\dots\times P_k$. Likewise, the clique factor of $Q$ for the graph product structure on $\overline{\Lambda}$ is $C_Q^e$, and its full parabolic direct product decomposition is $Q=C_Q^e\times Q_{j_1+1}\times\dots\times Q_m$.
\medskip

\underline{\textbf{Step 1.}} We show that  $\varphi(C^e_P) = C^e_Q$.
\smallskip

  Since $\overline{\Gamma}$ and $\overline{\Lambda}$ are strongly reduced, this follows from \cite[Lemmas~A.20 and~A.22]{EH}. Indeed, \cite[Lemma~A.22]{EH} characterizes the triviality of the clique factor by an algebraic property called $(\mathrm{Q}_{\mathrm{tc}})$ which is invariant under isomorphisms, so $C_P^e=\{e\}$ if and only if $C_Q^e=\{e\}$. Assuming now that $C_P^e\neq\{e\}$, \cite[Lemma~A.20]{EH} characterizes the fact that the $\overline{\Gamma}$-type of $P$ is a clique by an algebraic property which is invariant under isomorphisms: in the language of \cite{EH}, this property is the absence of a subgroup $F\subseteq P$ such that $(G,P,F)$ satisfies $(\mathrm{Q}_{\mathrm{fact}})$. So the $\overline{\Gamma}$-type of $P$ is a clique if and only if the $\overline{\Lambda}$-type of $Q$ is a clique; in this case $P=C_P^e$ and $Q=C_Q^e$ and the desired property holds. Finally, if the $\overline{\Gamma}$-type of $P$ is not a clique, then \cite[Lemma~A.20]{EH} chracterizes the clique factor of $P$ (for the graph product structure over $\overline{\Gamma}$) by an algebraic property which is invariant under isomorphisms: in the language of \cite[Lemma~A.20]{EH}, $C_P^e$ is the only subgroup of $P$ such that $(G,P,C_P^e)$ satisfies Property~$(\mathrm{Q}_{\mathrm{fact}})$. It again follows that $\varphi(C_P^e)=C_Q^e$.

\medskip

\underline{\textbf{Step 2.}} By induction on $i\le i_1$, we will construct automorphisms $\tau_i\in\Aut(G)$ which are products of $\Gamma$-twists, such that  
\begin{enumerate}
    \item for every $j\le i$, there exists $\bar{w_j}\in V\overline{\Lambda}$ such that $\varphi\circ\tau_i(P_j)=H_{\bar{w}_j}$;
    \item $\varphi\circ\tau_i(C_P^e)=C_Q^e$;
    \item for every $\Gamma$-parabolic subgroup $P'\subseteq G$, if $\varphi(P')$ is a $\Lambda$-parabolic subgroup, then $\varphi\circ\tau_i(P')$ is conjugate to $\varphi(P')$;
    \item if $\varphi$ is $(\Gamma,\Lambda)$-twistless, then $\tau_i=\mathrm{id}$.
\end{enumerate}

\smallskip

Let $\tau_0:=\mathrm{id}$. Then let $i\in\{1,\dots,i_1\}$, and assume that $\tau_{i-1}$ has been constructed. Let $\varphi_i:=\varphi\circ\tau_{i-1}$. We will construct $\tau_i$ as a product of twists with twisted subgroup $P_i$, and twistor in a subgroup $R_i$ that we now define. We refer to page~\pageref{def:twistor} for the introduction of twistor and twisted subgroup.

\textbf{$\drsh$ Definition of $S_i$ and $R_i$.} Let $S_i:=P_i^+\cap C_P^e$, where the notation $P_i^+$ is understood for the graph structure on $\overline{\Gamma}$. Equivalently $S_i$ is the parabolic subgroup generated by all vertex groups $D\subseteq C_P^e$ (for the graph product structure on $\overline{\Gamma}$) that satisfy $N_G(P_i)\subseteq N_G(D)$. Note that $P_i\subseteq S_i$ and that $N_G(S_i)=N_G(P_i)$. Let $R_i:=P_i^\perp\cap S_i$, so that $S_i\simeq P_i\times R_i$, where here and from now on in this proof when we write $\perp$, this is understood for the graph product structure on $\overline{\Gamma}$. 

\begin{claim}\label{claim:step-1}
The images $\varphi_i(S_i)$ and $\varphi_i(R_i)$ are $\overline{\Lambda}$-parabolic, whence $\Lambda$-parabolic.
\end{claim}


\begin{proof}[Proof of the claim]~

\textbf{Introducing tight subgroups.} The key point of this proof is the following notion which is invariant under isomorphisms: say that a subgroup $A\subseteq C_P^e$ is \deffont{tight} if for every subgroup $A'\subseteq C_P^e$ that properly contains $A$, one has $\langle C_P^e,N_G(A')\rangle\subsetneq \langle C_P^e,N_G(A)\rangle$. We define tight subgroups of $C_Q^e$ in the same way. Notice that:
\begin{enumerate}[label={(\alph*)}]
\item \label{item:tight1} If $A\subseteq C_P^e$ is a subgroup, and $B$ is the $\overline{\Gamma}$-parabolic support of $A$, Remark~\ref{Fact:Normalizer-support} gives $N_G(A)=N_B(A)\times B^\perp$, while $N_G(B)=B\times B^\perp$. As $B\subseteq C_P^e$, it follows that \begin{equation*}
    \left\langle C^e_P,N_G(A)\right\rangle
    =\left\langle C_P^e, B^\perp\right\rangle 
    \quad \text{and} \quad
    \left\langle C^e_P,N_G(B)\right\rangle
    =\left\langle C_P^e, B^\perp\right\rangle.
\end{equation*}
\item \label{item:tight2} Notice additionally that as the $\overline{\Gamma}$-type of $C_P^e$ is a clique, we have $C_P^e\subseteq B\times B^\perp$, so in fact $\langle C_P^e, B^\perp\rangle=B\times B^\perp=N_G(B)$.
\item \label{item:tight-parab} By point \ref{item:tight1}, tight subgroups of $C^e_P$ are necessarily $\overline{\Gamma}$-parabolic.
\item \label{item:tight-CNS} By points \ref{item:tight1} and~\ref{item:tight2}, a parabolic subgroup $A$ is tight if and only if for every \emph{$\mathit{\overline{\Gamma}}$-parabolic} subgroup $A'\subseteq C^e_P$ that properly contains $A$, one has $N_G(A')\subsetneq N_G(A)$.
\end{enumerate} 
\medskip

\textbf{Showing that {$S_i$ is tight and} $\varphi_i(S_i)$ is $\overline{\Lambda}$-parabolic.} We now prove that $S_i$ is tight, using~\ref{item:tight-CNS}. Since tightness is preserved by isomorphism, this also implies that $\varphi_i(S_i)$ is tight, whence $\overline{\Lambda}$-parabolic by \ref{item:tight-parab}.

So let $U\subseteq C_P^e$ be a parabolic subgroup that properly contains $S_i$. By definition of $S_i$ and since $S_i \subsetneq U$, there exists a vertex group $D\subseteq U$ such that $N_G(P_i)\not\subseteq N_G(D)$. Since $N_G(P_i)=N_G(S_i)$, this rewrites as $N_G(S_i)\not\subseteq N_G(D)$, or in other words 
\begin{equation}\label{eq:1}
    N_G(S_i)\cap N_G(D)\subsetneq N_G(S_i).
\end{equation}
Let $\Delta^*_U$ (resp.\ $\Delta^*_{S_i}$, resp.\ $\Delta^*_D$) be the full subgraph of $\overline{\Gamma}$ whose vertex set consists of all vertices that are in the type of $U^\perp$ (resp.\ $S_i^\perp$, resp.\ $D^\perp$) but not in the type of $C_P^e$. Then,
\begin{equation*}
    N_G(U)=\left\langle{C_P^e, G_{\Delta_U^*}} \right\rangle, \quad 
    N_G(S_i)=\left\langle{C_P^e, G_{\Delta_{S_i}^*}} \right\rangle, \quad 
    N_G(D)=\left\langle{C_P^e, G_{\Delta_D^*}} \right\rangle.
\end{equation*} 
As $D\subseteq U$ and $S_i\subseteq U$, we have $\Delta_U^*\subseteq \Delta_D^*$ and $\Delta_U^*\subseteq \Delta_{S_i}^*$. Therefore 
\begin{equation}\label{eq:2}
    N_G(U)\subseteq N_G(S_i)\cap N_G(D).
\end{equation}
Combining Equations~\eqref{eq:1} and~\eqref{eq:2} shows that $N_G(U)\subsetneq N_G(S_i)$, so by \ref{item:tight-CNS} $S_i$ is tight.

\medskip

\textbf{Showing that $\varphi_i(R_i)$ is $\overline{\Lambda}$-parabolic.}
Let us first prove that every vertex group $B\subseteq R_i$ is contained in a tight parabolic subgroup $T_B\subseteq R_i$.

So let $T_B:=B^+\cap S_i$, in other words $T_B$ is the maximal parabolic subgroup of $S_i$ with $N_G(B)=N_G(T_B)$. Remark that since $B\subseteq S_i$ and $N_G(S_i)=N_G(P_i)$, we have 
\begin{equation}
    \label{eq:normalisateur-B-NGPi}
    N_G(P_i)\subseteq N_G(B).
\end{equation}
\begin{itemize}
    \item As $\overline{\Gamma}$ is strongly reduced, the vertices of $\overline{\Gamma}$ representing the types of $P_i$ and $B$ are not equivalent. Equation~\eqref{eq:normalisateur-B-NGPi} thus implies that $N_G(P_i)\subsetneq N_G(B)$.
    Therefore $T_B$ does not contain $P_i$, in particular $T_B$ is a proper $\overline{\Gamma}$-parabolic subgroup of $S_i$.
    Recalling that $S_i=P_i\times R_i$ shows that $T_B\subseteq R_i$. 
    \item We claim that $T_B$ is tight. If not, then by the above point \ref{item:tight-CNS}, $T_B$ is properly contained in a parabolic subgroup $R\subseteq C_P^e$ such that $N_G(R)=N_G(T_B)$. By maximality of $T_B$ inside $S_i$, necessarily $R$ is not contained in $S_i$. Let $\Delta_R^*$ (resp.\ $\Delta_{S_i}^*$) be the full subgraph of $\overline{\Gamma}$ whose vertex set consists of all vertices that are in the type of $R^\perp$ (resp.\ $S_i^\perp$) but not in the type of $C_P^e$. Then 
\begin{equation}\label{eq:3}
    N_G(S_i)=\left\langle { C_P^e, G_{\Delta_{S_i}^*} }\right\rangle, 
    \quad N_G\left( \left\langle{S_i,R}\right\rangle \right)=\left\langle {C_P^e, G_{\Delta_R^*\cap \Delta_{S_i}^*} } \right\rangle.
\end{equation}
On the other hand, Equation~\eqref{eq:normalisateur-B-NGPi} rewrites as $N_G(S_i)\subseteq N_G(R)$, which shows that $\Delta_{S_i}^*\subseteq \Delta_R^*$. Equation~\eqref{eq:3} therefore gives $N_G(S_i)=N_G(\langle S_i,R\rangle)$, contradicting the fact that $S_i$ is tight. So $T_B$ is tight.
\end{itemize}
Hence, every vertex group $B\subseteq R_i$ is contained in a tight parabolic subgroup $T_B\subseteq R_i$. 

This implies that $R_i$ is generated by tight subgroups of $C_P^e$, and therefore $\varphi_i(R_i)$ is generated by tight subgroups of $C_Q^e$. This and \ref{item:tight-parab}, together with the fact that the type of $C_Q^e$ is a clique, shows that $\varphi_i(R_i)$ is $\overline{\Lambda}$-parabolic.
\end{proof}

\textbf{$\drsh$ Definition of $\tau_i$.}
Remark that:
\begin{itemize}
    \item $S_i$ normalizes $R_i$ so $\varphi_i(S_i)$ normalizes $\varphi_i(R_i)$;
    \item and $R_i\subset S_i$, so $\varphi_i(R_i)\subset \varphi_i(S_i)$.
\end{itemize}
 Since moreover $\varphi_i(S_i)$ and $\varphi_i(R_i)$ are $\overline{\Lambda}$-parabolic subgroups (Claim~\ref{claim:step-1}), we can thus let $P'_i:=(\varphi_i(R_i))^\perp\cap\varphi_i(S_i)$, and write $\varphi_i(S_i)\simeq\varphi_i(R_i)\times P'_i$, a direct product of $\overline{\Lambda}$-parabolic subgroups. Recalling that $S_i\simeq R_i\times P_i$, Lemma~\ref{Lmm:detwister-un-facteur-direct} ensures that there exists a homomorphism $\twistor:P_i\to Z(R_i)$ such that, letting $\tau'_i\in\Aut(S_i)$ be the automorphism which is the identity on $R_i$ and sends every $g\in P_i$ to $g\twistor(g)$, we have $\varphi_i\circ\tau'_i(P_i)=P'_i$. 
 
 Note that if $\varphi$ is $(\Gamma,\Lambda)$-twistless, then by induction $\varphi_{i}$ is $(\Gamma,\Lambda)$-twistless, and by Lemma~\ref{lemma:twistless-collapse} it is also $(\overline{\Gamma},\overline{\Lambda})$-twistless. Therefore the type in $\overline{\Lambda}$ of the parabolic support of $\varphi_i(P_i)$ is not a join, so $\varphi_i(P_i)=P'_i$ and we can take $\tau'_i=\mathrm{id}$. 

Note also that $\tau'_i$ extends to an automorphism of $G$. Indeed, denote by $\bar v_i$ the vertex of $\overline{\Gamma}$ giving the type of $P_i$, and by $\bar w$ a vertex of $\overline{\Gamma}$ in the type of $S_i$. Then, by definition of $S_i$, we have $\lk_{\overline{\Gamma}}(\bar v_i) \subseteq \st_{\overline{\Gamma}}(\bar w)$. Hence we can extend $\tau'_i$ to an automorphism of $G$ by letting it be the identity on all $\overline{\Gamma}$-vertex groups different from $P_i$. This extension is, by definition, a product of $\overline{\Gamma}$-twists, and by Lemma~\ref{Lmm:Aut-and-graph-collapse} it is also a product of $\Gamma$-twists. 

Since $\overline{\Gamma}$ and $\overline{\Lambda}$ are strongly reduced, $P'_i=\varphi_i\circ\tau'_i(P_i)$ is necessarily a vertex group for the graph product structure over $\overline{\Lambda}$ (Lemma~\ref{Lmm:vertex-group-iso-strong-redution}).

We now let $\tau_i:=\tau_{i-1}\circ\tau'_i$. Then $\tau_i$ satisfies the first two and last requirements by construction. The third one is a consequence of Lemma~\ref{lemma:redress} and Remark~\ref{Rk:redress} applied to the isomorphism $\varphi_{i}:G\rightarrow H$, the automorphism $\tau'_i$ and $A=R_i$, $G_v=P_i$, $B=\varphi_{i}(R_i)$ and $K=P'_i$.

This completes Step~2.

\medskip

\underline{\textbf{Notational interlude.}} We now let  $\tilde\tau:=\tau_{i_1}$. The previous step shows that $i_1\le j_1$. Using the isomorphism $\varphi^{-1}$ gives the reverse inequality, so $i_1=j_1$. Step~2 also gives us a bijection $\sigma:\{1,\dots,i_1\}\to\{1,\dots,i_1\}$ such that $\varphi\circ \tilde\tau(P_i)=Q_{\sigma(i)}$ for every $i\le i_1$. 

Let now $\psi:P/C^e_P\rightarrow Q/C^e_Q$ be the isomorphism induced from $\varphi\circ\tilde\tau$. Note that $P/C^e_P$ is isomorphic to a direct product $\bar{P}_{i_1+1}\times\dots\times \bar{P}_k$, where each factor $\bar{P}_i$ is isomorphic to $P_i$ and hence has trivial center (by \cite[Theorem~3.34]{Gre}) and does not split non-trivially as a direct product \cite[Proposition~2.8]{Gen}. Likewise $Q/C^e_Q$ is isomorphic to $\bar Q_{i_1+1}\times\dots\times \bar Q_m$. It follows from e.g.\ \cite[Lemma~3.5]{Fio} that there exists a bijection $\sigma:\{i_1+1,\dots,k\}\to\{i_1+1,\dots,m\}$ (which we think of as extending the bijection $\sigma:\{1,\dots,i_1\}\to\{1,\dots,i_1\}$) such that $\psi(\bar P_i)=\bar Q_{\sigma(i)}$ for every $i\ge i_1+1$.

For every $i\ge {i_1}+1$, denote by $\beta_i:P_i\to Q_{\sigma(i)}$ the isomorphism obtained from $\psi$ after identifying $P_i\simeq\bar P_i$ and $Q_{\sigma(i)}\simeq \bar Q_{\sigma(i)}$ (through the quotient maps $P\to P/C^e_P$ and $Q\to Q/C^e_Q$). Then there exists a homomorphism $\alpha_i:P_i\to Z(C^e_P)$ such that for every $g\in P_i$, one has $\varphi\circ\tilde\tau(g)=\beta_i(g)\varphi\circ\tilde\tau(\alpha_i(g))$.

Let us prove that if $\varphi$ is $(\Gamma,\Lambda)$-twistless then $\alpha_i(g)=e$ for every $g\in P_i$. Indeed, if $\varphi$ is $(\Gamma,\Lambda)$-twistless then $\tilde\tau=\mathrm{id}$. By Lemma~\ref{lemma:twistless-collapse}, $\varphi$ is also $(\overline{\Gamma},\overline{\Lambda})$-twistless. Now assume that $\alpha_i$ is non-trivial. Then there exists a vertex group $G_v\subseteq P_i$ such that $\alpha_i(G_v)\neq\{e\}$. The parabolic support of $\varphi(G_v)$ is then a non-trivial direct product, intersecting both factors of $C_Q^e\times Q_{\sigma(i)}$ non-trivially. This contradicts the fact that $\varphi$ is $(\overline{\Gamma},\overline{\Lambda})$-twistless. 

\medskip

\underline{\textbf{Step 3.}} Let $i>i_1$, and let $v\in V\Gamma$ be such that $G_v\subseteq P_i$. We now aim to show that the map sending $g$ to $g\alpha_i(g)^{-1}$ if $g\in G_v$, and sending $g$ to itself when $g\in G_{\Gamma\backslash \{v\}}$ determines an automorphism $\kappa_v$ of $G$ which is a product of $\Gamma$-twists.

\smallskip

Let $g\in G_v$. Write $\alpha_i(g)=\twistor_1(g)\cdots\twistor_{i_1}(g)$ with $\twistor_j(g)\in Z(P_j)$ for all $j\in \{1,\dots,i_1\}$. We will prove that $\alpha_i(g)\in G_v^+$, which will imply that $\kappa_v:g\mapsto g\alpha_i(g)^{-1}$ is a product of twists of $G$. For this, let $w\in \lk_\Gamma(v)$ and let $j\leq i_1$ be such that $G_w$ does not normalize $P_j$; we aim to prove that $\twistor_{j}(g)=e$.

Assume towards a contradiction that $\twistor_j(g)\neq e$. Then the parabolic support $R\subseteq Q$ of $\varphi\circ \tilde\tau(g)$ contains $Q_{\sigma(j)}$. Let $g'\in G_w$. Since $w\in \lk_{\Gamma}(v)$ then $g'$ commutes with $g$. Hence its image $\varphi\circ \tilde\tau(g')$ commutes with $\varphi\circ \tilde\tau(g)$, so $\varphi\circ \tilde\tau(g')$ normalizes $R$ (Remark~\ref{Fact:Normalizer-support}). Hence $\varphi\circ \tilde\tau(g')\in R\times R^{\perp}$. Write $\varphi\circ \tilde\tau(g')=h_1h_2$ with $h_1\in R$ and $h_2\in R^\perp$. But $\twistor_j(g)$ belongs to $Z(P_j)$, which is contained in $Z(P)$, so $\varphi\circ \tilde{\tau}(\twistor_j(g))\in Z(Q)$ and thus $h_1$ commutes with $\varphi\circ \tilde\tau(\twistor_j(g))$. Since $\varphi\circ \tilde{\tau}(\twistor_j(g))\in Q_{\sigma(j)}$ and $R^\perp\subseteq Q_{\sigma(j)}^\perp$, it follows that $\varphi\circ\tilde\tau(g')$ commutes with $\varphi\circ\tilde\tau(\twistor_j(g))$. Therefore $g'$ commutes with $\twistor_j(g)$, contradicting that $w$ is not adjacent to the vertex corresponding to the type of $P_j$. This contradiction shows that $\twistor_j(g)=e$.
\medskip

\underline{\textbf{Conclusion.}} Let $v_1,\dots,v_m$ be the vertices in the type of $P_{i_1+1}\times\dots\times P_k$. Note that for all $\ell \neq j$ the automorphisms $\kappa_{v_\ell}$ and $\kappa_{v_j}$ commute. Let $\tau:=\tilde\tau\circ\kappa_{v_1}\circ \cdots \circ \kappa_{v_m}$, then $\varphi\circ\tau$ sends $P_i$ to $Q_{\sigma(i)}$, for every $i\in\{1,\dots,k\}$. 

Let $P'\subseteq G$ be a parabolic subgroup such that $\varphi(P')$ is a
parabolic subgroup $Q'$. We now prove by induction that for every
$j\in\{0,\dots,m\}$, the image $\varphi\circ \tilde{\tau}\circ \kappa_{v_1}\circ
\dots \circ \kappa_{v_j}(P')$ is conjugate to $Q'$. With $j=m$, this will yield
the second conclusion of the lemma. The case where $j=0$ follows from the
properties of $\tilde\tau=\tau_{i_1}$ established in Step~2. For the inductive
step, let $j\in\{1,\dots,m\}$ and assume that the conclusion was already
established for $j-1$. Let $i\in\{i_1+1,\dots,k\}$ be such that $v_j$ belongs to
the type of $P_i$. The desired conclusion follows from Lemma~\ref{lemma:redress}
and Remark~\ref{Rk:redress}, applied to to the isomorphism $\varphi\circ
\tilde{\tau}\circ \kappa_{v_1}\circ \dots \circ \kappa_{v_{j-1}}$ (in place of
the automorphism $\varphi$ from Lemma~\ref{lemma:redress}), with $v:=v_j$,
$A:=C_P^e\cap G_v^+$, $B:=\varphi\circ \tilde{\tau}\circ \kappa_{v_1}\circ \dots
\circ \kappa_{v_{j-1}}(A)$, $K:=Q_{\sigma(i)}$ and $\xi(g):=\alpha_i(g)^{-1}$.
This completes the proof of the second conclusion of the lemma.

Finally, the third conclusion is satisfied by construction: if $\varphi$ is $(\Gamma,\Lambda)$-twistless, then $\tilde{\tau}$ and all the $\kappa_i$ are equal to $\mathrm{id}$. 
\end{proof}

\section{The isomorphism problem for graph products}\label{sec:endgame-isomorphism}

In this section we complete our proof of Theorem~\ref{theo:isomorphism}, solving the isomorphism problem for graph products, under the assumption that the graphs $\Gamma$ and $\Lambda$ are strongly reduced. We refer to Definition~\ref{de:clean} for the definition of a graphical and twistless isomorphism.

\begin{lemma}\label{lemma:endgame}
Let $\Gamma,\Lambda$ be strongly reduced finite simple graphs with at least three vertices. Let $G$ be a graph product over $\Gamma$ with vertex groups $(G_v)_{v\in V\Gamma}$, and let $H$ be a graph product over $\Lambda$ with vertex groups $(H_w)_{w\in V\Lambda}$. Let $\varphi:G\to H$ be an isomorphism. 

Then there exist an automorphism $\tau\in\Aut(G)$ which is a product of twists and an isomorphism $\theta:\Gamma\to\Lambda$ such that for every $v\in V\Gamma$, the following two hold: 
\begin{enumerate}
    \item\label{cond:theta1} $(v,\theta(v))$ is exclusively $(\varphi\circ\tau)$-paired;
    \item\label{cond:theta2} the vertex $\theta(v)$ is isolated in the type of the support of $\varphi\circ\tau(G_v)$.
\end{enumerate}
In particular, $\varphi\circ\tau$ is graphical and twistless.
\end{lemma}

\begin{proof}
Enumerate $V\Gamma=\{v_1,\dots,v_k\}$. For every $i\in\{1,\dots,k\}$, consider a chain of standard parabolic subgroups \[G_{v_i}\subseteq F_{n_i}^i\subseteq P_{n_i}^i\subseteq A_{n_i}^i\subseteq\dots\subseteq F_1^i\subseteq P_1^i\subseteq A_1^i\subseteq F^{i}_0=G\] where for every $\ell\in\{1,\dots,n_i\}$,
    \begin{itemize}
        \item $A_\ell^i$ is a thick minimal parabolic free factor of $F_{\ell-1}^i$;
        \item $P_\ell^i$ is a maximal parabolic subgroup of $A_\ell^i$ that splits non-trivially as a direct product of parabolic subgroups (this exists because the type of $A_\ell^i$ is connected and not reduced to one vertex);
        \item $F_\ell^i$ is a minimal parabolic direct factor of $P_\ell^i$;
        \item $G_{v_i}$ is a thin minimal parabolic free factor of $F_{n_i}^i$.
        \end{itemize}
Pairs $(i,\ell)$ with $i\in\{1,\dots,k\}$ and $\ell\in\{1,\dots,n_i\}$, or $(i,\ell)=(0,0)$, are ordered lexicographically. Namely $(j,\ell')\le (i,\ell)$ if and only if either $j<i$, or else $j=i$ and $\ell'\le\ell$.

Given $i\in\{1,\dots,k\}$ and $\ell\in\{1,\dots,n_i\}$, we let $\mathcal{P}_\ell^i$ be the set of all parabolic subgroups of the form $A_{\ell'}^j,P_{\ell'}^j,F_{\ell'}^j$ with $(j,\ell')\le (i,\ell)$. By convention we let $\mathcal{P}^0_0=\emptyset$.

We now prove by induction on the pair $(i,\ell)$ (either $(0,0)$ or with $i\in\{1,\dots,k\}$ and $\ell\in\{1,\dots,n_i\}$) that there exist $\tau_{i,\ell}\in\Aut(G)$ which is a product of twists, such that the isomorphism $\varphi\circ\tau_{i,\ell}$ sends every parabolic subgroup in $\mathcal{P}_\ell^i$ to a parabolic subgroup.

The base case where $(i,\ell)=(0,0)$ holds with $\tau_{i,\ell}=\mathrm{id}$ since the condition is void.

Now let $(i,\ell)\neq (0,0)$, and let $(j,\ell')$ be its immediate predecessor for our ordering. Assume that we have constructed $\tau_{j,\ell'}$. 
\begin{itemize}
    \item Since $A^i_\ell$ is a thick free factor of $F^i_{\ell-1}$, Corollary~\ref{cor:recognizing-free-factors} ensures that $\varphi\circ\tau_{j,\ell'}(A^i_\ell)$ is a parabolic subgroup whose type is connected and which is not conjugate to a vertex group.  
    \item Since $P_\ell^i$ is a maximal direct product parabolic subgroup of $A_\ell^i$, and the type of $A_\ell^i$ is connected and not reduced to one vertex, it follows from Lemma~\ref{lemma:maximal-direct-product} that $P_\ell^i$ is a maximal subgroup of $A_\ell^i$ that splits non-trivially as a direct product. So the same holds for $\varphi\circ\tau_{j,\ell'}(P^i_\ell)$ inside $\varphi\circ\tau_{j,\ell'}(A^i_\ell)$. Applying Lemma~\ref{lemma:maximal-direct-product} again ensures that $\varphi\circ\tau_{j,\ell'}(P_\ell^i)$ is a parabolic subgroup.
    \item Since $F_\ell^i$ is a minimal parabolic direct factor of $P_\ell^i$, we can apply Lemma~\ref{lemma:direct-product} to the isomorphism $\varphi \circ \tau_{j,\ell'}$ and to $P=P^i_\ell$, and deduce that we can precompose $\varphi\circ\tau_{j,\ell'}$ by a product of twists to obtain the desired $\tau_{i,\ell}$. Note that the second conclusion of Lemma~\ref{lemma:direct-product} ensures that precomposition does not affect the fact that the images of $A^i_\ell$, $P^i_\ell$ and of the parabolic subgroups in $\mathcal{P}^{i'}_{\ell'}$ with $(i',\ell')<(i,\ell)$, are still parabolic subgroups. 
\end{itemize}
This completes the construction of the automorphisms $\tau_{i,\ell}$. 

Let now $\tau:=\tau_{k,n_k}$. We will finally prove that $\varphi\circ\tau$ satisfies the conclusion of the lemma, and start by defining $\theta$. Since for every $i\in\{1,\dots,k\}$, the vertex group $G_{v_i}$ is a thin free factor of $F_{n_i}^i$, and since $\varphi\circ\tau$ sends $F_{n_i}^i$ to a parabolic subgroup, we can apply Lemma~\ref{lemma:free-product-1} with $P=F_{n_i}^i$ and with $Q=\varphi\circ\tau(F_{n_i}^i)$. This provides a vertex $w_i\in V\Lambda$, which is isolated in the type $\Delta_i$ of $\varphi\circ\tau(F_{n_i}^i)$, such that $(v_i,w_i)$ is exclusively $(\varphi\circ\tau)$-paired. 
In particular, the type of $\varphi\circ\tau(G_{v_i})$ contains $w_i$ and is contained in $\Delta_i$, so $w_i$ is a connected component of the type of $\varphi\circ\tau(G_{v_i})$. So defining $\theta(v_i)=w_i$, we obtain a map satisfying both conditions~\ref{cond:theta1} and \ref{cond:theta2}.
This moreover proves that $\varphi\circ\tau$ is graphical and twistless. 

Similarly, there exists $\tau'\in\Aut(H)$ such that $\varphi^{-1}\circ\tau'$ is graphical and twistless. By Lemma~\ref{lemma:blueprint}, applied to the isomorphisms $\varphi\circ\tau$ and $\varphi^{-1}\circ\tau'$, the map $\theta:V\Gamma\to V\Lambda$ sending each $v_i$ to $w_i$ is a graph isomorphism. 
\end{proof}

We are now in position to solve the isomorphism problem in the case of strongly reduced graphs. We state a slightly stronger version than the one given in %the introduction
Theorem~\ref{theo:isomorphism}.

\begin{theo}\label{theo:iso-strongly-reduced}
Let $\Gamma,\Lambda$ be finite simple graphs. Assume that either 
\begin{enumerate}
\item\label{item:assumption1} every join component of $\Gamma$ and of $\Lambda$ is strongly reduced with at least three vertices;
\item\label{item:assumption2} or every connected component of $\Gamma$ and of $\Lambda$ is strongly reduced with at least three vertices.
\end{enumerate}
Let $G$ be a graph product over $\Gamma$ with vertex groups $(G_v)_{v\in V\Gamma}$, and let $H$ be a graph product over $\Lambda$ with vertex groups $(H_w)_{w\in V\Lambda}$. 

Then $G$ and $H$ are isomorphic if and only if there is a graph isomorphism $\theta:\Gamma\to\Lambda$ such that for every $v\in V\Gamma$, the groups $G_v$ and $H_{\theta(v)}$ are isomorphic.
\end{theo}

\begin{proof}
The “if” direction is clear, so we focus on the “only if” direction. Let $\varphi:G\to H$ be an isomorphism.

\medskip

\textbf{Case 1.} We first treat the special case of the theorem where $\Gamma$ and $\Lambda$ are strongly reduced and have at least three vertices. 

 In this case, Lemma~\ref{lemma:endgame}, applied to $\varphi$, gives an automorphism $\tau\in \Aut(G)$ and a graph isomorphism $\theta:\Gamma\rightarrow \Lambda$ such that $(v,\theta(v))$ is exclusively $(\varphi\circ\tau)$-paired for all $v\in V\Gamma$. In particular $G_v$ and $H_{\theta(v)}$ are isomorphic for all $v\in V\Gamma$.

\medskip 

\textbf{Case 2.} We work under Assumption~\ref{item:assumption1}:
assume that $\Gamma=\Gamma_1\circ\dots\circ\Gamma_k$ and $\Lambda=\Lambda_1\circ\dots\circ\Lambda_l$, where the factors $\Gamma_i$ and $\Lambda_j$ are strongly reduced with at least three vertices.

Then $G$ splits as a direct product $G=G_{\Gamma_1}\times\dots\times G_{\Gamma_k}$ whose factors have trivial center \cite[Theorem~3.34]{Gre} and do not further split as a direct product \cite[Proposition~2.8]{Gen} (and likewise $H=H_{\Lambda_1}\times\dots\times H_{\Lambda_l}$). Since $G$ and $H$ are isomorphic, by e.g.\ \cite[Lemma~3.5]{Fio}, it follows that $k=l$, and $\varphi$ sends every $G_{\Gamma_i}$ to one of the $H_{\Lambda_j}$. The conclusion thus follows from the strongly reduced case treated above. 

\medskip 

\textbf{Case 3.} We finally work under Assumption~\ref{item:assumption2}: every connected component of $\Gamma$ and of $\Lambda$ is strongly reduced and has at least three vertices.

Let $\Gamma_1,\dots,\Gamma_k$ and $\Lambda_1,\dots,\Lambda_l$ be the connected components of $\Gamma$ and $\Lambda$. Then $G=G_{\Gamma_1}\ast\dots\ast G_{\Gamma_k}$ and $H=H_{\Lambda_1}\ast\dots\ast H_{\Lambda_l}$, and these are Grushko decompositions. If $G$ and $H$ are isomorphic, it follows that $k=l$ and that $\varphi$ sends every $G_{\Gamma_i}$ to a conjugate of some $H_{\Lambda_j}$. Again the conclusion follows from Case 1.
\end{proof}



% ======================================
% Part 2 - Automorphisms
% ======================================
\part{Generating the automorphism group} \label{part:automorphisms}
The goal of this second part is to provide an explicit generating set for the automorphism group of a graph product (Theorem~\ref{theo:main-aut}). As in the previous part, we will analyze direct products and free products separately. The results obtained in Section~\ref{sec:direct-products} are sufficient to deal with direct products, but we have some more work to do on free products to obtain the generating set of our automorphism group (see Section~\ref{sec:free-products-2} below).

\section{Permutationless automorphisms}\label{sec:permutationless}

We will need the following definition.

\begin{de}[Permutationless automorphism]
Let $\Gamma$ be a finite simple graph, and let $G$ be a graph product over $\Gamma$. An automorphism $\varphi\in\Aut(G)$ is \deffont{$\Gamma$-permutationless} if for every $v\in V\Gamma$, the type of the parabolic support of $\varphi(G_v)$ contains $v$.
\end{de}

\begin{rk}\label{rk:permutationless-collapse}
If $\pi:\Gamma\to\overline{\Gamma}$ is a graph collapse, and if $\varphi$ is $\Gamma$-permutationless, then $\varphi$ is also $\overline\Gamma$-permutationless.

Indeed, let $\bar v\in V\overline{\Gamma}$, and let $v\in V\Gamma$ be a preimage of $\bar v$ under $\pi$. Since $\varphi$ is $\Gamma$-permutationless, $\varphi(G_v)$ is not contained in any conjugate of $G_{\Gamma\setminus\{v\}}$. But $\varphi(G_v)\subseteq\varphi(G_{\bar v})$, and $G_{\overline{\Gamma}\setminus\{\bar v\}}\subseteq G_{\Gamma\setminus\{v\}}$. So $\varphi(G_{\bar v})$ is not contained in any conjugate to $G_{\overline{\Gamma}\setminus\{\bar v\}}$. Hence $\bar v$ belongs to the type of the parabolic support of $\varphi(G_{\bar v})$.
\end{rk}

Our solution to the isomorphism problem for graph products ensures the existence of permutationless automorphisms.

\begin{lemma}\label{lemma:permutationless}
 Let $\Gamma$ be an almost strongly reduced finite simple graph which is connected and not an edge. Let $G$ be a graph product over $\Gamma$ with vertex groups $(G_v)_{v\in V\Gamma}$. Assume that all vertex groups are either isomorphic to $\mathbb{Z}$ or freely indecomposable. 

Then for every automorphism $\varphi\in\Aut(G)$, there exists an automorphism $\alpha\in\Aut(G)$ which is a product of transvections, partial conjugations and local automorphisms, such that $\varphi\circ\alpha$ is $\Gamma$-permutationless and $\Gamma$-twistless.
\end{lemma}

\begin{proof}
The conclusion is clear if $\Gamma$ is reduced to one vertex, so we assume otherwise. Since $\Gamma$ is almost strongly reduced and not an edge, it does not split non-trivially as a join. Since in addition $\Gamma$ is connected and not reduced to a vertex, its strong reduction $\overline\Gamma$ has at least three vertices. We denote by $\pi:\Gamma\to\overline{\Gamma}$ the associated graph collapse.  

\medskip

\underline{\textbf{Step 1.}} Lemma~\ref{lemma:endgame} gives us an automorphism $\alpha_1\in\Aut(G)$ which is a product of $\overline{\Gamma}$-twists (and hence, of $\Gamma$-twists by Lemma~\ref{Lmm:Aut-and-graph-collapse}), and an automorphism $\bar\theta\in\Aut(\overline{\Gamma})$, such that for every $\bar v\in V\overline{\Gamma}$, the following two hold:
\begin{enumerate}
    \item $(\bar v,\bar\theta(\bar v))$ is exclusively $(\varphi\circ\alpha_1)$-paired and in particular $G_{\bar{v}}\simeq G_{\bar\theta(\bar v)}$;
    \item $\bar\theta(\bar v)$ is isolated in the type (in $\overline{\Gamma}$) of the parabolic support of $\varphi\circ\alpha_1(G_{\bar v})$.
\end{enumerate} 

\medskip

\underline{\textbf{Step 2.}} We now aim at lifting $\bar\theta$ to an automorphism $\theta$ of $\Gamma$ such that for every $v\in V\Gamma$, the vertex groups $G_v$ and $G_{\theta(v)}$ are isomorphic.
Let $\bar v\in V\overline\Gamma$. Since $\Gamma$ is almost strongly reduced, the preimage $\pi^{-1}(\bar v)$ is an edgeless subgraph (whose vertices we denote by $w(\bar v,1),\dots,w(\bar v,k_{\bar v})$), such that every vertex group is either freely indecomposable or isomorphic to $\mathbb{Z}$. If $\bar v,\bar v'\in V\overline{\Gamma}$ carry isomorphic vertex groups, then the uniqueness of the Grushko decomposition (in the form of Remark~\ref{rk:grushko}) ensures that $k_{\bar v}=k_{\bar v'}$. We can thus choose the numbering of the vertices $w(\bar v,i)$ in such a way that if $\bar v,\bar v'\in V\overline\Gamma$ are two vertices that carry isomorphic vertex groups, then $G_{w(\bar v,i)}$ is isomorphic to $G_{w(\bar v',i)}$ for every $i\in\{1,\dots,k_{\bar v}\}$. Since the automorphism $\bar\theta$ sends every vertex $\bar v$ to a vertex $\bar\theta(\bar v)$ that carries an isomorphic vertex group, it follows that $\bar\theta$ lifts to an automorphism $\theta$ of $\Gamma$ with the desired property, by sending $w(\bar v,i)$ to $w(\bar\theta(\bar v),i)$.

\medskip
 
\underline{\textbf{Step 3.}}~Let $\beta\in\Aut(G)$ be a $\Gamma$-local automorphism such that $\beta(G_v)=G_{\theta(v)}$ for every $v\in V\Gamma$. Note that in particular $\beta(G_{\bar v})=G_{\bar\theta(\bar v)}$ for every $\bar v\in V\overline{\Gamma}$. Recall from Definition~\ref{de:paired} that $(\bar v,\bar\theta(\bar v))$ being $(\varphi\circ\alpha_1)$-paired means in particular that $r_{\bar\theta(\bar v)}\circ\varphi\circ\alpha_1$ restricts to an isomorphism from $G_{\bar v}$ to $G_{\bar\theta(\bar v)}$. For every vertex $\bar v\in V\overline{\Gamma}$, we then let $\alpha_{\bar v}:=(r_{\bar\theta(\bar v)}\circ\varphi\circ\alpha_1)^{-1}\circ\beta_{|G_{\bar v}}$. This is an automorphism of the free product $G_{\bar v}$ such that for every $i\in\{1,\dots,k_{\bar v}\}$, the image  $r_{\bar\theta(\bar v)}\circ\varphi\circ\alpha_1\circ\alpha_{\bar v}(G_{w(\bar v,i)})$ is equal to $G_{w(\bar\theta(\bar v),i)}$. By Theorem~\ref{theo:FR} and Remark~\ref{rk:FR}, the automorphism $\alpha_{\bar v}$ is a product of folds, partial conjugations and local automorphisms, for the graph product structure of $G_{\bar v}$ on the edgeless graph $\pi^{-1}(\bar v)$. Since $\pi^{-1}(\bar v)$ is collapsible, Lemma~\ref{Lmm:Extending-aut-defined-on-collapsible} ensures that these folds, partial conjugations and local automorphisms extend to automorphisms of $G$, which are then $\Gamma$-folds, $\Gamma$-partial conjugations and $\Gamma$-local automorphisms. We denote by $\tilde{\alpha}_{\bar v}$ the extension of $\alpha_{\bar v}$ to an automorphism of $G$ which is the identity on all vertex groups not in $\pi^{-1}(\bar v)$. 

We let $\alpha_2$ be the product of all $\tilde\alpha_{\bar v}$ with $\bar v$ varying in $V\overline\Gamma$ (these commute because the subgraphs $\pi^{-1}(\bar v)$ are pairwise disjoint). 

\medskip

\underline{\textbf{Step 4.}}~ 
Let now $\alpha:=\alpha_1\circ\alpha_2\circ\beta^{-1}$. We claim that $\varphi\circ\alpha$ is both $\Gamma$-permutationless and $\Gamma$-twistless.

\textbf{$\drsh$ $\Gamma$-permutationless.} By construction, for every $v\in V\Gamma$, we have $\beta^{-1}(G_v)=G_{\theta^{-1}(v)}$, and $\theta(v)$ belongs to the type of the parabolic support of $\varphi\circ\alpha_1\circ\alpha_2(G_v)$. It follows that $\varphi\circ\alpha$ is $\Gamma$-permutationless.

\textbf{$\drsh$ $\Gamma$-twistless.} Since $\beta$ sends vertex groups to vertex groups, it is enough to prove that $\varphi\circ\alpha_1\circ\alpha_2$ is $\Gamma$-twistless. Let $v\in V\Gamma$, and let $\bar v:=\pi(v)$. Then, by definition of $\alpha_2$ we have $\varphi\circ\alpha_1\circ\alpha_2(G_v) \subseteq\varphi\circ\alpha_1(G_{\bar v})$. Since $\bar{\theta}(\bar v)$ is isolated in the $\overline\Gamma$-type of the parabolic support of $\varphi\circ\alpha_1(G_{\bar v})$, and $\pi^{-1}(\bar\theta(\bar v))$ is edgeless, it follows that $\theta(v)$ is isolated in the $\Gamma$-type of the parabolic support of $\varphi\circ\alpha_1\circ\alpha_2(G_v)$. In particular, this type does not split non-trivially as a join. 
\end{proof}

\section{More on automorphisms of free products}\label{sec:free-products-2}


The goal of this section is to prove Lemma~\ref{lemma:free-product} below, which builds upon Lemma~\ref{lemma:free-product-1}. Our first lemma revisits the results from Section~\ref{sec:free-products} in the presence of permutationless automorphisms.

\begin{lemma}\label{lemma:behaviour}
Let $\Gamma$ be a finite simple graph with at least three vertices, and let $G$ be a graph product over $\Gamma$ with vertex groups $(G_v)_{v\in V\Gamma}$. Assume that
\begin{itemize}
\item $\Gamma$ is almost strongly reduced;
\item all vertex groups $G_v$ are either free or freely indecomposable;
\item no two vertices of $\Gamma$ carrying free vertex groups are equivalent.
\end{itemize}
Let $P$ be a standard parabolic subgroup, and let $P=P_1\ast\dots\ast P_k$ be a full parabolic free product decomposition, with all $P_i$ standard.

Then there exists a bijection $\sigma$ of $\{1,\dots,k\}$ such that for every $\Gamma$-permutationless automorphism $\psi\in\Aut(G)$ satisfying $\psi(P)=P$, and every $i\in\{1,\dots,k\}$, one of the following holds:
\begin{itemize}
\item $\psi(P_{\sigma(i)})$ is conjugate to $P_{\sigma(i)}$, or
\item $P_{\sigma(i)}$ is free, $\psi(P_{\sigma(i)}^+\cap P)$ is conjugate to $P_{\sigma(i)}^+\cap P$, and $P_{\sigma(i)}^+\cap P\subseteq \langle P_{\sigma(1)},\dots,P_{\sigma(i)}\rangle$.
\end{itemize}
\end{lemma}

\begin{proof}
\textbf{1 -- Definition of $\sigma$.} Let $\sigma$ be a bijection of $\{1,\dots,k\}$ such that there exist $i_0\le i_1$ with the following properties:
\begin{itemize}
\item $P_{\sigma(1)},\dots,P_{\sigma(i_0)}$ are thick; 
\item $P_{\sigma(i_0+1)},\dots,P_{\sigma(i_1)}$ are thin and are not free groups; 
\item $P_{\sigma(i_1+1)},\dots,P_{\sigma(k)}$ are thin and are free groups.
\end{itemize}
By assumption, the vertices of $\Gamma$ corresponding to the types $v_{\sigma(i)}$ of the subgroups $P_{\sigma(i)}$ with $i>i_1$ are pairwise inequivalent. We can (and shall) thus modify the definition of $\sigma$ on $\{i_1+1,\dots,k\}$ to ensure that 
\begin{equation}\label{ordering2}\tag{$\star\star\star$}
\text{for every $j>i>i_1$, we have $v_{\sigma(i)}\not\leq v_{\sigma(j)}$.}
\end{equation}
 This ensures that for $i>i_1$, we have $P_{\sigma(i)}^+\cap P\subseteq\langle P_{\sigma(1)},\dots,P_{\sigma(i)}\rangle$.

\medskip

\textbf{2 -- Proof of the desired properties.} For notational simplicity, from now on we will assume, up to renumbering the factors, that $\sigma=\mathrm{id}$. Let $\psi\in\Aut(G)$ be a $\Gamma$-permutationless automorphism such that $\psi(P)=P$.

\begin{claim}\label{claim:free1} 
For every $i\le i_1$ the subgroup $\psi(P_i)$ is conjugate to $P_i$.
\end{claim}

\begin{proof}[Proof of the claim]
Corollary~\ref{cor:recognizing-free-factors} ensures that for every $i\in\{1,\dots,i_0\}$, the image $\psi(P_i)$ is a minimal parabolic thick free factor and is thus conjugate to $P_{j}$ for some $j\leq i_0$. But $\psi$ is $\Gamma$-permutationless, hence $\psi(P_i)$ is conjugate to $P_i$. Finally, Lemma~\ref{lemma:free} ensures that the same holds for free factors $P_i$ with $i\in\{i_0+1,\dots,i_1\}$.
\end{proof}

Lemma~\ref{lemma:free-product-1}, applied with $\Lambda=\Gamma$ and $Q=P$, gives a permutation $\sigma'$ of $\{1,\dots,k\}$ such that for every $i\in\{1,\dots,k\}$, denoting by $v_i,v'_i$ the respective types of $P_i,P'_{\sigma'(i)}$, the pair $(v_i,v'_i)$ is exclusively $\psi$-paired, and one of the following two holds: 
\begin{enumerate}
\item $\psi(P_i)$ is conjugate to $P_{\sigma'(i)}$, or 
\item $\psi(P_i^+\cap P)$ is conjugate to $P_{\sigma'(i)}^+\cap P$ and $ P_{\sigma'(i)}^+\cap P\subseteq \langle  P_{\sigma'(1)},\dots, P_{\sigma'(i)}\rangle$. 
\end{enumerate}
We aim to prove that $\sigma'=\mathrm{id}$.

Let $X\subseteq\{1,\dots,k\}$ be the subset made of indices such that $\psi(P_i)$ is conjugate to $P_{\sigma'(i)}$. Since $\psi$ is $\Gamma$-permutationless, we have $\sigma'(i)=i$ for every $i\in X$. For $i\le i_1$, Claim~\ref{claim:free1} ensures that $\psi(P_i)$ is conjugate to $P_i$. Since $(v_i,v'_i)$ is exclusively $\psi$-paired, necessarily $v'_i=v_i$ and hence $P_{\sigma'(i)}=P_i$. It follows that $X$ contains $\{1,\dots,i_1\}$.

There remains to see that $\sigma'(i)=i$ for every $i\in\{1,\dots,k\}\setminus X$. So let $n$ be the order of~$\sigma'$. For every $i\in \{1,\dots,k\}\setminus X$, since $\psi$ is $\Gamma$-permutationless and $\psi(P^+_i\cap P)$ is conjugate to $P^+_{\sigma'(i)}\cap P$, the vertex $v_i$ belongs to the type of $P^+_{\sigma'(i)}$; so $v_i\geq  v_{\sigma'(i)}$, which by \eqref{ordering2} implies that $\sigma'(i)\le i$. We thus have 
\[i\geq\sigma'(i)\geq {\sigma'}^2(i)\geq\dots\geq{\sigma'}^n(i)=i,\] 
therefore $\sigma'(i)=i$, as desired.     
\end{proof}

We will also need the following observation extracted from the work of Laurence, whose proof we recall for the convenience of the reader. It will be used in the sequel to extend automorphisms of a free product to automorphisms of an ambient graph product.

\begin{lemma}[{Laurence \cite[Proposition~6.5]{Lau}}]\label{lemma:Laurence} 
Let $\Gamma$ be a finite simple graph. Let $\Delta\subseteq\Gamma$ be a full subgraph, let $v\in V\Delta$ be an isolated vertex of $\Delta$. 
Let $C$ be a connected component of $\Gamma\setminus\st_\Gamma(v)$, and let $u$ be a vertex in $\Delta\setminus C$. 
 
 Then $C$ is contained in a connected component of $\Gamma\setminus\st_\Gamma(u)$.
\end{lemma}

\begin{proof}
The conclusion is clear if $u=v$, so we assume otherwise.

Let $w_1,w_2$ be two vertices in $C$. As $C$ is contained in a connected component of $\Gamma\setminus\st_\Gamma(v)$, we can find a path $\gamma$ in $\Gamma\setminus\st_\Gamma(v)$ joining $w_1$ to $w_2$. 

We claim that $\gamma$ is contained in a connected component of $\Gamma\setminus\st_\Gamma(u)$, so assume for the sake of contradiction that it contains a vertex $x\in\st_\Gamma(u)$. By choice of $\gamma$, we have $x\notin\st_\Gamma(v)$. We also have $u\notin\st_\Gamma(v)$ because $u\in V\Delta\setminus\{v\}$ and $v$ is isolated in $\Delta$. Then the union of the edge $[u,x]$ and of the subpath of $\gamma$ from $x$ to $w_2$ is a path in $\Gamma\setminus\st_\Gamma(v)$ joining $u$ to $w_2$, contradicting that $u\notin C$. 
\end{proof} 

The following lemma generalizes a theorem of Servatius \cite[Theorem~3]{Ser} from right-angled Artin groups to general graph products.

\begin{lemma}\label{lemma:cg}
Let $\Gamma$ be a finite simple graph, and $G$ be a graph product over $\Gamma$. Let $z\in V\Gamma$ and let $x,y\in V\Gamma\backslash \st_{\Gamma}(z)$ such that $x$ and $y$ belong to the same connected component of $\Gamma\setminus\st_\Gamma(z)$. Let $\varphi\in\Aut(G)$ be a $\Gamma$-permutationless automorphism.

If $\varphi(G_x)\subseteq G_{\Gamma\setminus\{z\}}$, then $\varphi(G_y)\subseteq G_{\Gamma\setminus\{z\}}$.
\end{lemma}

\begin{rk}
The assumption that $\varphi$ is $\Gamma$-permutationless cannot be dropped. As an example, assume that $\Gamma$ is a pentagon, and that all vertex groups are isomorphic. Let $x,y$ be adjacent vertices, and $z$ be the unique vertex that is adjacent to neither $x$ nor $y$. In particular $x$ and $y$ belong to the same connected component of $\Gamma\setminus\st_\Gamma(z)$. Let $\theta:\Gamma\to\Gamma$ be a graph automorphism such that $\theta(y)=z$ (and thus $\theta(x)\neq z$), and let $\varphi:G\to G$ be a local automorphism such that $\varphi(G_v)=G_{\theta(v)}$ for every $v\in V\Gamma$. Then $\varphi(G_x)\subseteq G_{\Gamma\setminus\{z\}}$, while $\varphi(G_y)=G_z$, so the conclusion of the lemma does not hold for the automorphism $\varphi$ (which is not $\Gamma$-permutationless). 
\end{rk}

\begin{proof}
It is enough to prove the lemma when $x$ and $y$ are adjacent. Indeed, in general, one can take a finite path contained in $\Gamma\setminus\st_{\Gamma}(z)$ and joining $x$ to $y$, and propagate the conclusion of the lemma along this path.
    
So suppose that $\varphi(G_x)\subseteq G_{\Gamma\setminus\{z\}}$ and assume towards a contradiction that there exists $g_y\in G_y$ such that $\varphi(g_y)\notin G_{\Gamma\setminus\{z\}}$. Then any normal form for $\varphi(g_y)$ (see page~\pageref{Def:Normal-forms}) must contain a syllable in $G_z$. Since $\varphi$ is $\Gamma$-permutationless, $x$ belongs to the type of the parabolic support of $\varphi(G_x)$. Lemma~\ref{lemma:cu} therefore enables us to choose $g_x\in G_x$ such that the type of the parabolic support of $\varphi(g_x)$ contains $x$. Since $\varphi(g_x)\in G_{\Gamma\setminus\{z\}}$, a normal form of $\varphi(g_x)$ does not contain any syllable in $G_z$. By \cite[Lemma~3.16]{Gre} and its proof, we can write $\varphi(g_x)=h h_x h^{-1}$ with $h_x$ graphically cyclically reduced (i.e.\ its parabolic support is standard), in such a way that neither $h$ nor $h_x$ contains a syllable in $G_z$. The parabolic support of $h_x$ is thus equal to $G_\Upsilon$ for some full subgraph $\Upsilon$ of $\Gamma$ that does not contain $z$. Consider the join decomposition \[\Upsilon=\Upsilon_1\circ\dots\circ\Upsilon_{j_0}\circ\Upsilon_{j_0+1}\circ\dots\circ\Upsilon_k,\] where the subgraphs $\Upsilon_j$ do not split non-trivially as a join, and the decomposition is written so that $\Upsilon_1,\dots,\Upsilon_{j_0}$ are exactly the join components reduced to one vertex. Write $h_x=h_1\dots h_k$ with $h_j\in G_{\Upsilon_j}$. By Barkauskas' description of centralizers in graph products (recalled in Theorem~\ref{theo:barkauskas}), we have 
    \begin{equation*}
        C_G\left(\varphi\left(g_x\right)\right)=h\left(
        C_{G_{\Upsilon_1}}(h_1)\times\dots\times C_{G_{\Upsilon_{j_0}}}(h_{j_0})\times Z_{j_0+1} \times\dots\times Z_k\times G_{\Upsilon^\perp} 
    \right)h^{-1},
    \end{equation*}
    where for every $j\in\{j_0+1,\dots,k\}$, the group $Z_j$ is the maximal cyclic subgroup of $G_{\Upsilon_j}$ that contains $h_j$. Since $\varphi(g_y)\in C_G(\varphi(g_x))$, and since we assume that every normal form for $\varphi(g_y)$ has a syllable in $G_z$, necessarily $z$ belongs to $\Upsilon^\perp$. Therefore $h_x$ commutes with $G_z$, so $h_x\in G_{\st_{\Gamma}(z)}$. Hence, the type of the parabolic support of $\varphi(g_x)$ is contained in $\st_{\Gamma}(z)$. But by our choice of $g_x$, the vertex $x$ belongs to this type, so $x\in\st_{\Gamma}(z)$, a contradiction.
\end{proof}

We are now in position to prove the main lemma of the section. Recall that the definition of a saturated parabolic subgroup was introduced on page~\pageref{de:saturated}.
  
\begin{lemma}\label{lemma:free-product}
Let $\Gamma$ be a finite simple graph with at least three vertices, and let $G$ be a graph product over $\Gamma$ with vertex groups $(G_v)_{v\in V\Gamma}$. Assume that $\Gamma$ is almost strongly reduced, and that every vertex group $G_v$ is either isomorphic to $\mathbb{Z}$, or freely indecomposable. Let $\varphi\in\Aut(G)$ be a $\Gamma$-permutationless automorphism, and let $P\subseteq G$ be a saturated standard parabolic subgroup such that $\varphi(P)=P$. Let $P=P_1\ast\dots\ast P_k$ be a full parabolic free product decomposition of $P$, with all $P_i$ standard.

Then there exists an automorphism $\tau\in\Aut(G)$ which is a product of folds, partial conjugations and local automorphisms, such that:
\begin{itemize}
\item $\varphi\circ\tau$ is $\Gamma$-permutationless;
\item if $\varphi$ is twistless, then $\varphi\circ\tau$ is twistless;
\item for every $i\in\{1,\dots,k\}$, the subgroup $\varphi\circ\tau(P_i)$ is conjugate to $P_i$;
\item for every $w\in V\Gamma$, if $\varphi(G_{w})$ is conjugate to $G_w$, then $\varphi\circ\tau(G_{w})$ is conjugate to $G_w$.
\end{itemize}
\end{lemma}

\begin{proof}
We first assume that $\Gamma$ is disconnected. As $\Gamma$ is almost strongly reduced, it follows that $\Gamma$ is edgeless. Since all vertex groups are either isomorphic to $\mathbb{Z}$, or freely indecomposable, the conclusion follows from Theorem~\ref{theo:FR} and Remark~\ref{rk:FR} in this case.

We now assume that $\Gamma$ is connected. 

\medskip

\underline{\textbf{Set-up: Collapsing free groups.}}~Let $\Delta_P$ be the type of $P$. Let $\overline\Gamma$ be the graph obtained from $\Gamma$ by collapsing all maximal collapsible edgeless subgraphs of $\Gamma$ whose vertex groups are all isomorphic to $\mathbb{Z}$. Since $P$ is saturated, it inherits a new full parabolic free product decomposition $P=\bar P_1\ast\dots\ast \bar P_{k'}$, where each $\bar P_i$ is either equal to $P_i$, or is a finitely generated free group that decomposes as a free product $\bar P_i=G_{v_{i,1}}\ast\dots\ast G_{v_{i,k_i}}$ of vertex groups of $G$ (for its graph product structure on $\Gamma$). Remark~\ref{rk:permutationless-collapse} ensures that $\varphi$ remains $\overline\Gamma$-permutationless. 
\medskip

\underline{\textbf{Case $|V\overline{\Gamma}|\leq 2$.}}~Assume first that $\overline{\Gamma}$ has at most two vertices. Since $\Gamma$ is connected, non-empty and not reduced to one vertex, $\overline{\Gamma}$ is an edge. Since $\Gamma$ is almost strongly reduced, $\Gamma$ is a join of two edgeless graphs $\Gamma=\Gamma_1\circ\Gamma_2$. If $P$ does not split non-trivially as a free product, then the conclusion of the lemma holds with $\tau=\mathrm{id}$. Otherwise, as $P$ is saturated, we have $P=G_{\Gamma_1}$ or $P=G_{\Gamma_2}$, a graph product over an edgeless graph –~note that in fact $P$ is a finitely generated free group. In particular, the type $\Delta_P$ of $P$ is edgeless. By Theorem~\ref{theo:FR} and Remark~\ref{rk:FR} there exists $\tau_P\in \Aut(P)$ which is a product of $\Delta_P$-folds, $\Delta_P$-partial conjugations and $\Delta_P$-local automorphisms, such that $\varphi_{|P}\circ \tau_P=\mathrm{id}_P$. In particular $\varphi_{|P}\circ \tau_P$ is $\Delta_P$-permutationless and twistless. By Lemma~\ref{Lmm:Extending-aut-defined-on-collapsible}, extending $\tau_P$ by the identity on $G_{\Gamma\backslash \Delta_P}$ induces an automorphism $\tau$ of $G$ that is a product of $\Gamma$-folds, $\Gamma$-partial conjugations and $\Gamma$-local automorphisms and such that $\varphi\circ\tau$ satisfies the four wanted points.

\medskip

\underline{\textbf{Case $|V\overline{\Gamma}|\geq 3$.}}~From now on we assume that $\overline{\Gamma}$ has at least three vertices. Note that $\overline{\Gamma}$ is again almost strongly reduced, and that two vertices of $\overline{\Gamma}$ whose associated vertex groups are free are never equivalent. We can therefore apply Lemma~\ref{lemma:behaviour} to the graph $\overline{\Gamma}$: this tells us that up to reordering the factors $\bar P_i$, we can assume that for every $\overline{\Gamma}$-permutationless automorphism $\psi\in\Aut(G)$ satisfying $\psi(P)=P$, and every $i\in\{1,\dots,k'\}$, 
\begin{itemize}
\item either $\psi(\bar P_i)$ is conjugate to $\bar P_i$;
\item or $\bar P_i$ is free and $\psi(\bar P_i^+\cap P)$ is conjugate to $\bar P_i^+\cap P$ and $\bar P_i^+\cap P\subseteq \langle \bar P_1,\dots,\bar P_i\rangle$.
\end{itemize}
Here the notation $\bar P_i^+$ is understood for the graph product structure over $\overline{\Gamma}$.

\smallskip

\textbf{Unknitting the automorphism.} We will now prove by induction on $i\in\{0,\dots,k'\}$ that there exists an automorphism $\tau_i\in\Aut(G)$ which is a product of $\Gamma$-folds, $\Gamma$-partial conjugations, and $\Gamma$-local automorphisms, such that the following five conditions hold:
\begin{itemize}
    \item one has $\varphi\circ \tau_i(P)=P$ and for every $j\in\{1,\dots,i\}$, the subgroup $\varphi\circ\tau_i(\bar P_j)$ is conjugate to $\bar P_j$;
    \item for every $j\in\{1,\dots,i\}$, if $\bar P_j=G_{v_{j,1}}\ast\dots \ast G_{v_{j,k_j}}$, then $\varphi\circ\tau_i(G_{v_{j,l}})$ is conjugate to $G_{v_{j,l}}$ for every $l\in \{1,\dots,k_j\}$;
    \item for every $w\in V\Gamma$, if $\varphi(G_w)$ is conjugate to $G_w$, then $\varphi\circ\tau_i(G_w)$ is conjugate to $G_w$;
    \item the automorphism $\varphi\circ\tau_i$ is $\Gamma$-permutationless;
    \item if $\varphi$ is $\Gamma$-twistless, then $\varphi\circ\tau_i$ is $\Gamma$-twistless.
\end{itemize}

Once this is proved, setting $\tau:=\tau_{k'}$ will complete our proof of the lemma.

\medskip

The case $i=0$ holds with $\tau_i=\mathrm{id}$. 

Let now $i\in\{1,\dots,k'\}$, and assume that $\tau_{i-1}$ is constructed. Let $\varphi_i:=\varphi\circ\tau_{i-1}$. By induction $\varphi_i$ is $\Gamma$-permutationless, and therefore by Remark~\ref{rk:permutationless-collapse} it is also $\overline{\Gamma}$-permutationless. So either $\varphi_i(\bar P_i)$ is conjugate to $\bar P_i$, or else $\bar P_i$ is free and $\varphi_i(\bar P_i^+\cap P)$ is conjugate to $\bar P_i^+\cap P$, with $\bar P_i^+\cap P\subseteq\langle \bar P_1,\dots,\bar P_i\rangle$.

\medskip

\textbf{$\drsh$ Case 1:} $\varphi_i(\bar{P}_i)$ is conjugate to $\bar P_i$. 

If $\bar P_i$ is not a free group, we let $\tau_i:=\tau_{i-1}$ and all conclusions hold by induction.

So let us assume that $\bar P_i=G_{v_{i,1}}\ast\dots\ast G_{v_{i,k_i}}$ is a free group, and denote by $\Upsilon_i$ the type of $\bar P_i$ for the graph product structure on $\Gamma$. Since $\varphi(\bar{P}_i)$ is conjugate to $\bar{P}_i$, using \cite[Proposition~3.4]{AM}, there exists $\alpha\in \mathrm{Inn}(G)$ where the conjugating element belongs to $P$, such that $(\varphi\circ \alpha)_{|\bar{P}_i}$ is an automorphism $\beta_i$ of the free product $\bar{P}_i$. By Theorem~\ref{theo:FR}, $\beta_i$ is a product of local automorphisms, partial conjugations and folds of $G_{\Upsilon_i}=\bar{P}_i$. By Lemma~\ref{Lmm:Extending-aut-defined-on-collapsible}, these automorphisms extend to $\Gamma$-local automorphisms, $\Gamma$-partial conjugations and $\Gamma$-folds of $G$ that are the identity on $G_{\Gamma\setminus\Upsilon_i}$. Hence $\beta_i$ also extends to $G$ and this extension is a product of $\Gamma$-local automorphisms, $\Gamma$-partial conjugations and $\Gamma$-folds of $G$. Letting  $\tau_i:=\alpha\circ \beta_i^{-1}$ gives the desired automorphism. 

\medskip 

\textbf{$\drsh$ Case 2:} $\bar P_i$ is free and $\varphi_i(\bar P_i^+\cap P)$ is conjugate to $\bar P_i^+\cap P$, with $\bar P_i^+\cap P\subseteq \langle \bar P_1,\dots,\bar P_i\rangle$. 

Let $\beta\in\Inn(G)$ with conjugating element in $P$ with $\varphi_i\circ\beta(\bar P_i^+\cap P)=\bar P_i^+\cap P$ (which exists by \cite[Lemma~3.7]{AM}).

Let $\bar{v}_i$ be the vertex of $\overline\Gamma$ giving the type of $\bar P_i$. Let $\tilde{\Delta}_1,\dots,\tilde{\Delta}_{p+1}$ be the connected components of $\overline{\Gamma}\setminus\lk_{\overline\Gamma}(\bar{v}_i)$ that have non-empty intersection with the type of $\bar P_i^+\cap P$. For every $j\in\{1,\dots,p+1\}$, let $\Delta_j$ be the intersection of $\tilde\Delta_j$ with the type of $\bar P_i^+\cap P$. Note that one of the subgraphs $\Delta_j$ is reduced to the vertex $\bar{v}_i$. Without loss of generality, we assume that $\Delta_{p+1}=\{\bar{v}_i\}$. Letting $A_j= G_{\Delta_j}$ for every $j\in\{1,\dots,p\}$, we then have $\bar P_i^+\cap P=A_1\ast\dots\ast A_p\ast \bar P_i$.

\begin{claim}\label{claim:complement-star}
Let $\ell\in\{1,\dots,p\}$. The image $\varphi_i\circ\beta(A_\ell)$ is conjugate to $A_\ell$ in $\bar P_i^+\cap P$.
\end{claim}

\begin{proof}[Proof of the claim] 
Recall that $\beta\in\Inn(G)$ satisfies $\varphi_i\circ\beta(\bar P_i^+\cap P)=\bar P_i^+\cap P$.
\medskip
 
    \textbf{Subclaim 1.} 
    Let $\ell\in\{1,\dots,p\}$ and let $x$ be a vertex in the $\overline{\Gamma}$-type of $A_\ell$. There exists $\alpha_{\ell,x}\in \mathrm{Inn}(G)$ with conjugating element in $\bar P_i^+\cap P$ and $\varphi_i\circ \beta\circ \alpha_{\ell,x}(G_x)\subseteq A_\ell$.
    
    \smallskip

    Indeed let $j\in\{1,\dots,k'\}$ be such that $G_x\subseteq \bar P_j$. Since $G_x\subseteq \bar{P}_i^+$ and $x\neq \bar{v}_i$, and since $\bar P_i^+\cap P\subseteq \langle \bar P_1,\dots,\bar P_i\rangle$, we have $j<i$. Moreover, by our induction hypothesis $\varphi_i(\bar P_j)$ is conjugate to $\bar P_j$, so the image $\varphi_i(G_x)$ is contained in a conjugate of $\bar P_j$. By our standing assumption of Case~2, we also have $\varphi_i\circ\beta(G_x)\subseteq \bar P_i^+\cap P$. So by \cite[Proposition~3.4]{AM}, there exists an inner automorphism $\alpha_{\ell,x}\in\mathrm{Inn}(G)$, where the conjugating element belongs to $\bar P_i^+\cap P$, such that $\varphi_i\circ\beta\circ\alpha_{\ell,x}(G_x)\subseteq \bar P_j\cap \bar P_i^+$.  
    
    We will finally prove that $\bar P_j\cap\bar P_i^+\subseteq A_\ell$, which will imply that $\varphi_i\circ\beta\circ\alpha_{\ell,x}(G_x)\subseteq A_\ell$, as desired. Let $\Delta'$ be the $\overline{\Gamma}$-type of $\bar{P}_j$, a connected subgraph of $\overline{\Gamma}$. Since $\bar{P}_j$ and $\bar{P}_i$ are two free factors of $P$, we have $\Delta'\cap \lk_{\overline{\Gamma}}(\bar{v}_i)=\emptyset$. Hence $\Delta'$ is contained in a connected component of $\overline{\Gamma}\backslash \lk_{\overline{\Gamma}}(\bar{v}_i)$, namely $\tilde{\Delta}_\ell$, which contains $x$. Noting that since $\bar{P}_j\subset P$ we have $\bar{P}_j\cap \bar P^+_i \subset \bar{P}^+_i\cap P$, we finally get that the type of $\bar P_j\cap\bar P_i^+$ is a subgraph of $\Delta_\ell$.
    So $\bar P_j\cap \bar P_i^+\subseteq A_\ell$. \medskip
    
    \textbf{Subclaim~2.} Let $\ell\in\{1,\dots,p\}$ and $x\in V\Delta_\ell$. Let $\alpha_{\ell,x}\in\Inn(G)$ be as given by Subclaim~1. Then $\varphi_i\circ\beta\circ\alpha_{\ell,x}(A_\ell)\subseteq A_\ell$.
    \smallskip
    
    Indeed, let $y\in V\Delta_\ell$ and $z\in V\overline\Gamma\setminus V\Delta_\ell$. We will prove that $\varphi_i\circ\beta\circ\alpha_{\ell,x}(G_y)\subseteq G_{\overline\Gamma\setminus\{z\}}$. 
    \begin{itemize}
    \item We first assume that $z$ does not belong to the type $\Delta$ of $\bar P_i^+\cap P$. It follows from the definitions of $\beta$, and of $\alpha_{\ell,x}$ in Subclaim~1, that $\varphi_i\circ\beta\circ\alpha_{\ell,x}(\bar P_i^+\cap P)=\bar P_i^+\cap P$. Therefore $\varphi_i\circ\beta\circ\alpha_{\ell,x}(G_y)\subseteq \bar P_i^+\cap P\subseteq G_{\overline\Gamma\setminus\{z\}}$. 
    \item We now assume that $z\in V\Delta$. Lemma~\ref{lemma:Laurence}, applied to $\overline{\Gamma}$ with $\Delta$ the subgraph spanned by $\Delta_1\cup\dots\cup\Delta_{p+1}$,
    with $v=\bar{v}_i$, $C=\tilde\Delta_\ell$ and $u=z$, ensures that $x,y$ are in the same connected component of $\overline\Gamma\backslash \st_{\overline{\Gamma}}(z)$.
    By Lemma~\ref{lemma:cg} applied to the $\overline{\Gamma}$-permutationless automorphism $\varphi_i\circ\beta\circ\alpha_{\ell,x}$, we have $\varphi_i\circ\beta\circ\alpha_{\ell,x}(G_y)\subseteq G_{\overline\Gamma\setminus\{z\}}$.
    \end{itemize}
    As $\varphi_i\circ\beta\circ\alpha_{\ell,x}(G_y)\subseteq G_{\overline\Gamma\setminus\{z\}}$ for every $z\notin V\Delta_\ell$, it follows that $\varphi_i\circ\beta\circ\alpha_{\ell,x}(G_y)\subseteq A_\ell$. As this is true for every $y\in V\Delta_\ell$, it follows that $\varphi_i\circ\beta\circ\alpha_{\ell,x}(A_\ell)\subseteq A_\ell$, as desired.
    
    \medskip
    
    \textbf{Conclusion of the claim.} Subclaim~2 shows that for every $\ell\in\{1,\dots,p\}$, the image $\varphi_i\circ\beta(A_\ell)$ is contained in a conjugate of $A_\ell$ with conjugating element in $\bar P_i^+\cap P$. Since $\bar P_i^+\cap P$ splits as a free product of finitely many freely indecomposable factors and of a finitely generated free group, no automorphism of $\bar P_i^+\cap P$ sends a free factor of $\bar P_i^+\cap P$ strictly inside itself. Hence $\varphi_i\circ\beta(A_\ell)$ is conjugate to~$A_\ell$. 
\end{proof}

Recall that $\bar P_i$ is a free group. Write $\bar P_i = G_{v_{i,1}}\ast\dots\ast G_{v_{i,k_i}}$, where each $v_{i,j}$ is a vertex of $\Gamma$. For every $j\in\{1,\dots,k_i\}$, let $x_{i,j}$ be a generator of $G_{v_{i,j}}\simeq\mathbb{Z}$. Let $\calf_i$ be the set of $(\bar P_i^+\cap P)$-conjugacy classes of the subgroups $A_1,\dots,A_p$ and recall that $\bar{P}_i^+\cap P = A_1*\dots *A_p*\bar{P}_i$. Claim~\ref{claim:complement-star} ensures that $(\varphi_i\circ\beta)_{|\bar P_i^+\cap P}\in\Aut(\bar P_i^+\cap P,\calf_i)$. By Lemma~\ref{lemma:fouxe-rabinovitch}, there exists an automorphism $\alpha_i\in\Aut(\bar P_i^+\cap P,\calf_i^{(t)})$ such that for every $j\in\{1,\dots,k_i\}$, one has $\varphi_i\circ\beta\circ\alpha_i(x_{i,j})=x_{i,j}$. By Theorem~\ref{theo:FR} and Remark~\ref{rk:FR}, the automorphism $\alpha_i$ is a product of: 
\begin{itemize}
\item transvections, of the form $x_{i,j}\mapsto x_{i,j}g$, where $g$ belongs to some $A_j$ or is equal to (a power of) some $x_{i,l}$ with $l\neq j$;
\item partial conjugations, i.e.\ automorphisms conjugating one of the subgroups $A_j$ by an element in some $A_l$ with $l\neq j$ or by some power of $x_{i,l}$ with $l\in\{1,\dots,k_i\}$;  
\item inversions, i.e.\ automorphisms that send one $x_{i,j}$ to $x_{i,j}^{-1}$, are the identity on every $A_l$, and fix every $x_{i,l}$ with $l\neq j$;
    \item permutations, i.e.\ automorphisms that are the identity on every $A_l$, and send every $x_{i,j}$ to $x_{i,\sigma(j)}$ for some permutation $\sigma$ of $\{1,\dots,k_i\}$.
\end{itemize}
First note that all these automorphisms of $\bar P_i^+\cap P$ extend to automorphisms of $G$ which are products of $\Gamma$-folds, $\Gamma$-partial conjugations and $\Gamma$-local automorphisms. Indeed: 
\begin{itemize}
\item By the definition of $\bar P_i^+$, every transvection $x_{i,j}\mapsto x_{i,j}g$ as above extends to a product of $\Gamma$-folds.
\item Let $\Delta$ be the subgraph spanned by $\Delta_1\cup\dots\cup\Delta_{p+1}$. Let $\ell\in\{1,\dots,p\}$ and $\ell'\in\{1,\dots,p+1\}$ with $\ell\neq \ell'$, and let $z\in V\Delta\setminus V\Delta_\ell$. By Lemma~\ref{lemma:Laurence} applied with $v=\bar v_i$ and $C=\tilde{\Delta}_\ell$, the subgraph $\Delta_\ell$ is contained in some connected component $\Upsilon$ of $\overline{\Gamma}\setminus\st_{\overline{\Gamma}}(z)$. We observe that $\Upsilon\cap\Delta_{\ell'}=\emptyset$: indeed, as $z\ge \bar v_i$, a path joining $\Delta_\ell$ to $\Delta_{\ell'}$ in  $\overline{\Gamma}\setminus\st_{\overline{\Gamma}}(z)$ would be contained in $\overline{\Gamma}\setminus\st_{\overline{\Gamma}}(\bar v_i)$, contradicting that $\ell\neq\ell'$. This proves that for $g\in G_z$, the automorphism of $\bar P_i^+\cap P$ given by the partial conjugation of $A_\ell$ by $g$ extends to an automorphism of $G$ which is the conjugation by $g$ on $G_{\Upsilon}$ and the identity on $G_{\overline{\Gamma}\setminus\Upsilon}$. 
\item All inversions and permutations extend to $\Gamma$-local automorphisms.
\end{itemize}
Therefore $\alpha_i$ extends to an automorphism $\tilde{\alpha}_i$ of $G$. Let $\tau'_i:=\beta\circ\tilde\alpha_i$ and $\tau_i:=\tau_{i-1}\circ\tau'_i$, and observe that by construction $\tau'_i(P)=P$. 

Let us prove that $\tau_i$ satisfies the five desired properties stated at the beginning of the paragraph “Unknitting the automorphism”. 

We first prove that $\varphi\circ\tau_i$, which is equal to $\varphi_i\circ\tau'_i$, is $\Gamma$-permutationless, and that if $\varphi$ is $\Gamma$-twistless, then so is $\varphi\circ\tau_i$. Note that for every $u\in V\Gamma$, the following holds. 
\begin{itemize}
\item If $u$ projects to $\bar v_i$ in $\overline\Gamma$, then $G_u$ is a cyclic group generated by one of the elements $x_{i,j}$, and by construction $\varphi_i\circ\tau'_i(G_u)=G_u$. 
\item If $u$ does not project to $\bar v_i$ in $\overline\Gamma$, then $\tau'_i(G_u)$ is conjugate to $G_u$, so $\varphi_i\circ\tau'_i(G_u)$ is conjugate to $\varphi_i(G_u)$. 
\end{itemize}
It follows in both cases that $u$ belongs to the type of $\varphi_i\circ\tau'_i(G_u)$, using in the second case that $\varphi_i$ is $\Gamma$-permutationless by induction. So $\varphi\circ\tau_i$ is $\Gamma$-permutationless.

If $\varphi$ is $\Gamma$-twistless, it also follows in both cases that the type of the parabolic support of $\varphi_i\circ\tau'_i(G_u)$ does not split non-trivially as a join, using in the second case that $\varphi_i$ is $\Gamma$-twistless by induction. So $\varphi\circ\tau_i$ is $\Gamma$-twistless.

Since $\tau'_i(P)=P$, we obtain by induction $\varphi_i\circ\tau'_i(P)=P$. 
Let us now prove that $\varphi_i\circ \tau'_i(\bar{P}_j)$ is conjugate to $\bar{P}_j$ for all $j\in\{1,\dots,i-1\}$. If $\bar{P}_j$ is thick, this comes from Corollary~\ref{cor:recognizing-free-factors} together with the fact that $\varphi_i\circ\tau'_i$ is $\Gamma$-permutationless. If $\bar{P}_j$ is thin, this comes from noting that $\bar{P}_j$ is either contained in one $A_\ell$ or has trivial intersection with $\bar{P}^+_i\cap P$ and that $\tilde\alpha_i$ (and hence $\tau'_i$) is a product of automorphisms that restrict to a conjugation on each $A_\ell$ with $\ell\in\{1,\dots,p\}$. Finally, by construction $\varphi_i\circ \tau'_i(\bar{P}_i)=\bar{P}_i$. The second and third points follow from the above observations.
\end{proof}

\section{Generating the automorphism group of a graph product}\label{sec:endgame-automorphisms}

We now complete our proof of Theorem~\ref{theo:main-aut}, giving a generating set for the automorphism group of a graph product, in the case where the graph $\Gamma$ is almost strongly reduced.

\begin{theo}[Generators for $\Aut(G)$, almost strongly reduced case]\label{theo:conclusion}
Let $\Gamma$ be a finite simple graph which is not an edge. Let $G$ be a graph product over $\Gamma$ with vertex groups $(G_v)_{v\in V\Gamma}$. Assume that $\Gamma$ is almost strongly reduced and that all vertex groups $G_v$ are freely indecomposable or isomorphic to $\bZ$. 

Then every $\varphi\in\Aut(G)$ is a product of transvections, partial conjugations and local automorphisms.
\end{theo}

\begin{proof}
If $\Gamma$ is reduced to one vertex, all automorphisms are local. If $\Gamma$ is disconnected, then being almost strongly reduced, it must be edgeless. Then $G$ is the free product of the vertex groups $G_v$ and the conclusion follows from Theorem~\ref{theo:FR} and Remark~\ref{rk:FR}. 

From now on we assume that $\Gamma$ is connected and has at least three vertices. In view of Lemma~\ref{lemma:permutationless}, we can assume without loss of generality that $\varphi$ is $\Gamma$-permutationless and $\Gamma$-twistless. We moreover enumerate $V\Gamma=\{v_1,\dots,v_k\}$.
\medskip

\underline{\textbf{Step 1.}} We prove by induction on $i\in\{0,\dots,k\}$ that there exists $\tau_i\in\Aut(G)$ which is a product of transvections, partial conjugations and local automorphisms, such that $\varphi\circ\tau_i(G_{v_j})$ is conjugate to $G_{v_j}$ for every $j\in\{1,\dots,i\}$, and such that $\varphi\circ\tau_i$ is $\Gamma$-permutationless and $\Gamma$-twistless. 

\smallskip

The case $i=0$ is immediate with $\tau_0=\mathrm{id}$ (the condition is void). So let $i\in\{1,\dots,k\}$, assume that $\tau_{i-1}$ has been constructed and let us explain how to build $\tau_{i}$. Consider a chain of subgroups
\begin{equation*}
    G_{v_i}\subseteq F_{n}\subseteq P_{n}\subseteq A_{n}\subseteq\dots\subseteq F_1\subseteq P_1\subseteq A_1\subseteq F_0=G
\end{equation*} where for every  $\ell\in \{1,\dots,n\}$,
    \begin{itemize}
        \item $A_\ell$ is a thick minimal parabolic free factor of $F_{\ell-1}$;
        \item $P_\ell$ is a maximal parabolic subgroup of $A_\ell$ that splits non-trivially as a direct product;
        \item $F_\ell$ is a minimal parabolic direct factor of $P_\ell$;
        \item $G_{v_i}$ is a thin minimal parabolic free factor of $F_n$.
        \end{itemize}
Note that all subgroups $A_\ell,P_\ell$ and $F_\ell$ are saturated (where for $F_\ell$ we use the fact that $\Gamma$ is almost strongly reduced, so two equivalent vertices are always non-adjacent). We clearly have $\varphi(F_0)=F_0$. Arguing by induction, we now prove that for every $\ell\in\{1,\dots,n\}$, the subgroups $A_\ell, P_\ell, F_\ell$ are sent by $\varphi\circ\tau_{i-1}$ to conjugates of themselves.  

\begin{itemize}
    \item Since $A_\ell$ is a thick free factor of $F_{\ell-1}$,  Corollary~\ref{cor:recognizing-free-factors} ensures that $\varphi\circ\tau_{i-1}(A_\ell)$ is a thick free factor of $\varphi\circ\tau_{i-1}(F_{\ell-1})$. Since by induction $\varphi\circ\tau_{i-1}(F_{\ell-1})$ is conjugate to $F_{\ell-1}$, it follows that $\varphi\circ\tau_{i-1}(A_\ell)$ is conjugate to a thick free factor of $F_{\ell-1}$. Finally, since $\varphi\circ\tau_{i-1}$ is $\Gamma$-permutationless, we deduce that $\varphi\circ\tau_{i-1}(A_\ell)$ is conjugate to $A_\ell$.  
    \item Since $P_\ell$ is a maximal parabolic subgroup of $A_\ell$ that splits non-trivially as a direct product, automatically $\varphi\circ\tau_{i-1}$ sends $P_\ell$ to a maximal parabolic subgroup of $\varphi\circ\tau_{i-1}(A_\ell)$ that splits non-trivially as a direct product of parabolic subgroups by Lemma~\ref{lemma:maximal-direct-product}. Since $\varphi\circ\tau_{i-1}$ is $\Gamma$-permutationless, $\varphi\circ\tau_{i-1}(P_\ell)$ is conjugate to $P_\ell$.
    \item Since $F_\ell$ is a minimal parabolic direct factor of $P_\ell$ and $\varphi\circ\tau_{i-1}$ is $\Gamma$-twistless, Lemma~\ref{lemma:direct-product} ensures that $\varphi\circ\tau_{i-1}$ sends $F_\ell$ to minimal parabolic direct factor of $\varphi\circ\tau_{i-1}(P_\ell)$. Since $\varphi\circ\tau_{i-1}$ is $\Gamma$-permutationless, it follows that $\varphi\circ\tau_{i-1}(F_\ell)$ is conjugate to $F_\ell$.
\end{itemize}

We precompose $\varphi\circ\tau_{i-1}$ by an inner automorphism $\alpha$ so that $\varphi\circ\tau_{i-1}\circ\alpha(F_n)=F_n$. Finally, Lemma~\ref{lemma:free-product} gives an automorphism $\sigma_i$ which is a product of folds, partial conjugations and local automorphisms, such that for every $j\in\{1,\dots,i\}$, the subgroup $\varphi\circ\tau_{i-1}\circ\alpha\circ\sigma_{i}(G_{v_j})$ is conjugate to $G_{v_j}$, and such that $\varphi\circ\tau_{i-1}\circ\alpha\circ\sigma_{i}$ is $\Gamma$-permutationless and $\Gamma$-twistless. Letting $\tau_i:=\tau_{i-1}\circ \alpha\circ\sigma_i$ thus completes Step~1. 

\medskip

\underline{\textbf{Step 2: Conclusion.}}
The automorphism $\varphi\circ\tau_k$ sends every vertex group to a conjugate of a vertex group; in other words, it belongs to the subgroup $\PConj(G)\subseteq\Aut(G)$. Since by \cite[Theorem~A]{GM} the group $\PConj(G)$ is generated by local automorphisms and partial conjugations, the theorem follows. 
\end{proof}
% ======================================
% Part 3 - T reduced
% ======================================
\part{Beyond the strongly reduced case}\label{part:3}
In this last part, we generalize our solution to the isomorphism problem, and to finding a generating set for the automorphism group,  to graphs that are not necessarily strongly reduced or almost strongly reduced. This will be done by reducing all other cases to the strongly reduced (or almost strongly reduced) one.

\section{The isomorphism problem}\label{sec:3.1}

We start with our extension of the isomorphism problem. {When we worked with strongly reduced graphs, the graphs $\Gamma$ and $\Lambda$ did not contain any proper collapsible subgraphs, but there were no restrictions at all on vertex groups. Another set of assumptions in Theorem~\ref{theo:isomorphism} only requires that $\Gamma,\Lambda$ do not contain any collapsible cliques, but the counterpart is that we only allow graphically irreducible vertex groups. In general, one can impose that only certain types of subgraphs (in a family $\mathbb{T}$) are not allowed to be collapsible in $\Gamma$, and as a counterpart one needs to assume that the only way in which a vertex group can split as graph product is over a graph in the family $\mathbb{T}$. This is made precise in the following statement.}   

\begin{theo}\label{theo:isomorphism-complete}
Let $\Gamma,\Lambda$ be finite simple graphs with at least three vertices. Let $G$ be a graph product over $\Gamma$ with vertex groups $(G_v)_{v\in V\Gamma}$, and let $H$ be a graph product over $\Lambda$ with vertex groups $(H_w)_{w\in V\Lambda}$. 

Let $\mathbb{T}$ be a family of finite simple graphs with at least two vertices, that contains all cliques (with at least two vertices), and satisfies the following stability properties:  
\begin{itemize}
    \item if $\Delta\in\mathbb{T}$, then every connected component of $\Delta$ which is not reduced to one vertex also belongs to $\mathbb{T}$;
    \item if $\Delta\in\mathbb{T}$, then every join component of $\Delta$ which is not reduced to one vertex also belongs to $\mathbb{T}$.
\end{itemize}
Assume that the two following conditions are satisfied:
\begin{enumerate}
    \item\label{hyp:Treduced1} both $\Gamma$ and $\Lambda$ do not contain any proper collapsible subgraph in the family $\mathbb{T}$;
    \item\label{hyp:Treduced2} if a vertex group $G_v$ or $H_w$ splits non-trivially as a graph product over a graph $\Delta$, then $\Delta\in\mathbb{T}$. 
\end{enumerate}
Then $G$ is isomorphic to $H$ if and only if there exists a graph isomorphism $\theta:\Gamma\to\Lambda$ such that for every $v\in V\Gamma$, one has $H_{\theta(v)}\simeq G_v$.
\end{theo}

Theorem~\ref{theo:isomorphism} from the introduction, in the clique-reduced case, corresponds to the particular case where $\mathbb{T}$ is the family consisting of all cliques. The strongly reduced case, already established in Theorem~\ref{theo:iso-strongly-reduced}, corresponds to the case where $\mathbb{T}$ is the family of all finite simple graphs (not reduced to one vertex). The proof of Theorem~\ref{theo:isomorphism-complete} will consist in reducing to the strongly reduced case.

\begin{proof}
The “if” direction is clear; we focus on the “only if” direction. Assume that $G$ and $H$ are isomorphic. We will argue by induction on $|V\Gamma|$ but start with a few useful facts.

\begin{center}
    \textbf{Part 1: Preliminary remarks on collapsible subgraphs}
\end{center}

\begin{Fact}\label{Fact:Stability}
Let $\Delta\subseteq \Gamma$ be a proper collapsible subgraph, then 
\begin{enumerate}
    \item Any collapsible subgraph of $\Delta$ is collapsible as a subgraph of $\Gamma$ as well, so $\Delta$ also satisfies Condition~\ref{hyp:Treduced1};
    \item And viewing $G_\Delta$ as a graph product over $\Delta$ with vertex groups $(G_v)_{v\in V\Delta}$, it naturally satisfies Condition~\ref{hyp:Treduced2}.
\end{enumerate}
\end{Fact}

\begin{Fact}\label{Fact:ecrasable-pas-arete}
    Let $\Delta\subseteq \Gamma$ be a collapsible subgraph, then $\Delta$ is not an edge.
\end{Fact}
Arguing by contraposition, if $\Delta$ is an edge, since $|V\Gamma|\geq 3$, then $\Delta$ is a proper subgraph of $\Gamma$. But $\mathbb{T}$ contains all cliques, hence $\Delta\in\mathbb{T}$, and thus $\Delta$ is not collapsible.

\begin{Fact}\label{Fact:Cas-un-sommet} Let $\Delta\subseteq \Gamma$ and $\Upsilon\subseteq \Lambda$ be collapsible subgraphs such that $G_\Delta \simeq H_\Upsilon$. Then $\Delta$ is reduced to one vertex if and only if $\Upsilon$ is. 
\end{Fact}
Indeed, assume for the sake of contradiction that $\Delta=\{u\}$ for some $u\in V\Gamma$ and that $\Upsilon$ is not reduced to a vertex. Since $\Upsilon$ is collapsible, we have $\Upsilon \notin \mathbb{T}$. But $H_{\Upsilon}$ is isomorphic to $G_\Delta=G_u$, so $G_u$ splits as a graph product over $\Upsilon$ and hence $\Upsilon\in \mathbb{T}$, a contradiction. This being symmetric in $\Delta$ and $\Upsilon$, we obtain the wanted equivalence.

\begin{Fact}
\label{Fact:Disconnected-with-two-vertices} Let $\Delta\subseteq \Gamma$ and $\Upsilon\subseteq \Lambda$ be collapsible subgraphs. Assume that $G_\Delta \simeq H_\Upsilon$ and that $\Delta$ or $\Upsilon$ is the disconnected graph with exactly two vertices. Then there exists a graph isomorphism
$\theta_\Delta:\Delta\rightarrow \Upsilon$ such that $G_{v}\simeq H_{\theta_\Delta(v)}$ for all $v\in V\Delta$.
\end{Fact}
\begin{proof}[Proof of Fact~\ref{Fact:Disconnected-with-two-vertices}.]
Assume that $\Upsilon$ is the disconnected graph on two vertices, denoted by $w_1$ and $w_2$ (the case where $\Delta$ is the disconnected graph on two vertices being symmetrical). In this case $H_{\Upsilon}=H_{w_1}*H_{w_2}$. 

Since $|V\Lambda|\ge 3$, the subgraph $\Upsilon$ is proper, so Condition~\ref{hyp:Treduced1} implies that $\Upsilon \notin \mathbb{T}$. Thus, Condition~\ref{hyp:Treduced2} ensures that all vertex groups of $G$ and $H$ (and in particular $H_{w_1}$ and $H_{w_2}$) are freely indecomposable or isomorphic to $\bZ$. Therefore, letting $\Delta_1,\dots,\Delta_m$ be the connected components of $\Delta$,
we have $G_\Delta=G_{\Delta_1}*\dots *G_{\Delta_{m}}$, and each $G_{\Delta_i}$ is either freely indecomposable or isomorphic to $\bZ$. Remark that for all $i\in\{1,\dots,m\}$, since $\Delta_i$ is a connected component of the collapsible subgraph $\Delta$, it is either a vertex or collapsible, in which case $\Delta_i\notin \mathbb{T}$ (Condition~\ref{hyp:Treduced1}). 
Since $G_\Delta\simeq H_\Upsilon$, it follows from Remark~\ref{rk:grushko} that $m=2$ and that up to renumbering the factors $H_{w_1}\simeq G_{\Delta_1}$ and $H_{w_2}\simeq G_{\Delta_2}$. Since $\Delta_i\notin \mathbb{T}$, the group $H_{w_i}$ cannot split as a graph product over $\Delta_i$ unless $\Delta_i$ is reduced to one point (Condition~\ref{hyp:Treduced2}). Therefore $\Delta_1$ and $\Delta_2$ are points and the fact follows.
\end{proof}

\begin{Fact}\label{Fact:graph-collapses} Let $\pi_\Gamma:\Gamma \rightarrow \overline{\Gamma}$ and $\pi_\Lambda:\Lambda \rightarrow \overline{\Lambda}$ be two graph collapses of $\Gamma$ and $\Lambda$ respectively. If there exist
\begin{itemize}
    \item a graph isomorphism $\bar{\theta}:\overline{\Gamma} \rightarrow \overline{\Lambda}$,
    \item and for all $\bar{v}\in V\overline{\Gamma}$, a graph isomorphism $\theta_{\bar{v}}:\pi^{-1}_\Gamma(\bar{v}) \rightarrow \pi^{-1}_\Lambda(\bar{\theta}(\bar{v}))$ such that $G_{u}\simeq H_{\theta_{\bar{v}}(u)}$ for all $u\in V\pi^{-1}_\Gamma(\bar{v})$,
\end{itemize}
then there exists a graph isomorphism ${\theta}:{\Gamma} \rightarrow {\Lambda}$ such that for every $v\in V\Gamma$, one has $G_v\simeq H_{\theta(v)}$.
\end{Fact}

\begin{proof}[Proof of Fact~\ref{Fact:graph-collapses}] Let $\bar{v}\in
  V\overline{\Gamma}$. Note that $\pi^{-1}_\Gamma(\bar{v})$ is a collapsible
  subgraph of $\Gamma$ (possibly reduced to one vertex).
  Therefore defining $\theta$ by $\theta(v):=\theta_{\pi_{\Gamma}(v)}(v)$ for all $v\in V\Gamma$, gives the wanted isomorphism.
\end{proof}

\newpage

\begin{center}
    \textbf{Part 2: Induction}    
\end{center}
Let $\varphi:G\to H$ be an isomorphism. Let $\overline{\Gamma},\overline{\Lambda}$ be respective strong reductions of $\Gamma,\Lambda$ and $\pi_\Gamma:\Gamma\to\overline{\Gamma}$ and $\pi_{\Lambda}:\Lambda\to\overline{\Lambda}$ be the respective graph collapses. 

We now distinguish different cases, depending on whether $\overline{\Gamma}$ and $\overline{\Lambda}$ have at least three vertices and we can thus apply our result from Part~\ref{part:isomorphism}, or not. Recall (see Remark~\ref{Rq:strong-reduction-as-at-least-3-vertices} on page~\pageref{Rq:strong-reduction-as-at-least-3-vertices}) that if $\overline{\Gamma}$ has two vertices, then either $\Gamma$ is disconnected or $\Gamma$ is connected and splits non-trivially as a join. These possibilities are dealt with in Cases 2 and 3. 
\medskip


\textbf{Case 1.} $\overline{\Gamma}$ and $\overline{\Lambda}$ both have at least three vertices. 
\smallskip

By Theorem~\ref{theo:iso-strongly-reduced}, there exists an isomorphism $\bar\theta:\overline{\Gamma}\to\overline{\Lambda}$ such that $H_{\bar\theta(\bar v)}\simeq G_{\bar v}$ for all $\bar{v}\in V\overline{\Gamma}$. We claim that for all $\bar{v}\in V\overline{\Gamma}$, letting $\Delta:=\pi^{-1}_\Gamma \left(\bar{v}\right)$ and $\Upsilon:=\pi^{-1}_\Lambda \left(\bar\theta(\bar{v})\right)$, there exists a graph isomorphism $\theta_{\bar{v}}$ from $\Delta$ to $\Upsilon$ such that $H_{\theta_{\bar{v}}(u)}\simeq G_u$ for all $u\in V\Delta$. 

By Fact~\ref{Fact:ecrasable-pas-arete}, neither $\Delta$ nor $\Upsilon$ is an edge. We now distinguish three cases depending on the values of $|V\Delta|$ and $|V\Upsilon|$.
\begin{enumerate}
    \item If $|V\Delta|=1$ or $|V\Upsilon|=1$, then both $\Delta$ and $\Upsilon$ are reduced to a vertex by Fact~\ref{Fact:Cas-un-sommet}. In this case we can let $\theta_{\bar{v}}$ be the map that sends the unique vertex of $\Delta$ to the unique vertex of $\Upsilon$.
    \item If $\Delta$ or $\Upsilon$ is the disconnected graph on two vertices, the graph isomorphism given by Fact~\ref{Fact:Disconnected-with-two-vertices} is the desired $\theta_{\bar{v}}$.
    \item If both $\Delta$ and $\Upsilon$ have at least three vertices, then by Fact~\ref{Fact:Stability}, Conditions~\ref{hyp:Treduced1} and~\ref{hyp:Treduced2} are satisfied for the graph product structures of $G_\Delta$ over $\Delta$ and of $H_\Upsilon$ over $\Upsilon$.
    So our induction hypothesis implies the existence of a graph isomorphism $\theta_{\bar{v}}:\Delta \rightarrow \Upsilon$ such that $G_u\simeq H_{\theta_{\bar v}(u)}$ for all $u\in V\Delta$, as desired.
\end{enumerate}
Fact~\ref{Fact:graph-collapses} then gives wanted graph isomorphism $\theta:\Gamma\rightarrow \Lambda$. 
\medskip

\textbf{Case 2.} $\Gamma$ or $\Lambda$ is disconnected.
\smallskip

Remark that if $\Gamma$ is disconnected then $G$ splits non-trivially as a free product, and so does $H$, hence $\Lambda$ is disconnected as well. Thus in this case, both $\Gamma$ and $\Lambda$ are disconnected.
Let $\Gamma_1,\dots,\Gamma_k$ be the connected components of $\Gamma$, and $\Lambda_1,\dots,\Lambda_m$ be the connected components of $\Lambda$.
Note in particular that $\Gamma_i,\Lambda_j$ are collapsible proper subgraphs of respectively $\Gamma$ and $\Lambda$. Therefore by Fact~\ref{Fact:ecrasable-pas-arete}, these subgraphs are either reduced to a vertex or have at least three vertices.

\begin{claim} 
There exists a bijection $\sigma:\{1,\dots,k\}\to\{1,\dots,m\}$ such that $G_{\Gamma_i}\simeq H_{\Lambda_{\sigma(i)}}$ for all $i\in\{1,\dots,k\}$. 
\end{claim}

\begin{proof}[Proof of the claim]
Up to renumbering $\{\Gamma_i\}_i$ and $\{\Lambda_j\}_j$, we let $\ell\in \{1,\dots,k\}$ and $\ell' \in \{1,\dots,m\}$ be such that $\{G_{\Gamma_i}\}_{1\leq i\leq \ell}$ and $\{H_{\Lambda_j}\}_{1\leq j\leq \ell'}$ are the thick free factors of $G$ and $H$ respectively.

We first prove that for every $i\in\{1,\dots,\ell\}$, the image $\varphi(G_{\Gamma_i})$ is conjugate to some thick free factor $H_{\Lambda_j}$, i.e.\ with $j\in \{1,\dots,\ell'\}$.  

So let $i\in\{1,\dots,\ell\}$. Since $G_{\Gamma_i}$ is freely indecomposable, $\varphi(G_{\Gamma_i})$ is either conjugate to some $H_{\Lambda_j}$ with $j\in\{1,\dots,\ell'\}$, or to a free factor of $H_{\Lambda_j}$ with $j>\ell'$ (possibly to $H_{\Lambda_j}$ itself). We claim that the latter case is impossible. Indeed, in this case $H_{\Lambda_j}$ would split as a graph product over a finite simple graph $\Upsilon_j$, one of whose connected components is $\Gamma_i$. Since $\Gamma_i$ is collapsible in $\Gamma$, Condition~\ref{hyp:Treduced1} ensures that $\Gamma_i\notin\mathbb{T}$. Since by assumption the family $\mathbb{T}$ is stable under passing to connected components, we also have $\Upsilon_j\notin\mathbb{T}$. This yields a contradiction to Condition~\ref{hyp:Treduced2}, and proves the wanted assertion.

Arguing with $\varphi^{-1}$ shows that $\ell=\ell'$, and that up to renumbering the factors, $\varphi(G_{\Gamma_i})$ is conjugate to $H_{\Lambda_i}$ for every $i\in\{1,\dots,\ell\}$. In particular
\begin{equation}
\label{eq:facteur-épais-sous-groupe-distingué}
\varphi\left(\left \llangle G_{\Gamma_i} \ : \  i\in \{1,\dots,\ell\}\right\rrangle\right) =
\left\llangle H_{\Lambda_i} \ : \  i\in \{1,\dots,\ell\}\right\rrangle.
\end{equation}
It it thus now sufficient to prove that $k=m$ and that up to reordering $G_{\Gamma_i}\simeq H_{\Lambda_i}$ for all $i\in \{\ell+1,\dots,k\}$.

\textbf{$\drsh$ Case (i).} $G$ or $H$ has no thin free factor. 

 Then Equation~\eqref{eq:facteur-épais-sous-groupe-distingué} and the fact that $\varphi$ is an isomorphism imply that both groups have no thin free factor and the claim follows.

\smallskip

\textbf{$\drsh$ Case (ii).} Both $\Gamma$ and $\Lambda$ have exactly one connected component reduced to one vertex, namely $\Gamma_k=\Gamma_{\ell+1}$ and $\Lambda_k=\Lambda_{\ell'+1}$.

Then $k=m$ and the claim follows noting that we have:
\begin{equation*}
    G_{\Gamma_{k}} \simeq G_\Gamma/ \llangle G_{\Gamma_i} \ : \ i\in \{1,\dots,\ell\}\rrangle
    \simeq H_\Lambda / \llangle H_{\Lambda_i} \ : \ i\in \{1,\dots,\ell\}\rrangle \simeq H_{\Lambda_{k}}.
\end{equation*}

\textbf{$\drsh$ Case (iii).} $\Gamma$ or $\Lambda$ has (at least) two connected components reduced to one vertex. 

Then the graph consisting of two isolated vertices does not belong to $\mathbb{T}$ (Condition~\ref{hyp:Treduced1}), which forces all vertex groups to be freely indecomposable or isomorphic to $\bZ$ (Condition~\ref{hyp:Treduced2}). As all factors of both decompositions $*^k_{i=1}G_{\Gamma_i}$ and $*^m_{i=1} H_{\Lambda_i}$ are either freely indecomposable or isomorphic to $\bZ$, the claim follows from Remark~\ref{rk:grushko}.
\end{proof}

The conclusion in Case~2 follows from the claim, arguing by induction and using Fact~\ref{Fact:graph-collapses}.
\medskip

%-------------------------------------------
\textbf{Case 3.} $\Gamma$ or $\Lambda$ is connected and splits non-trivially as a join.
\smallskip 

Let $\Gamma_1,\dots,\Gamma_k$ be the join components of $\Gamma$, and let $\Lambda_1,\dots,\Lambda_m$ be the join components of $\Lambda$. We make the following claim.

\begin{claim}
There exists a bijection $\sigma:\{1,\dots,k\}\to\{1,\dots,m\}$ such that $G_{\Gamma_i}\simeq H_{\Lambda_{\sigma(i)}}$ for all $i\in\{1,\dots,k\}$.
\end{claim}

\begin{proof}[Proof of the claim] 
We distinguish two cases.
\smallskip 

\textbf{$\drsh$ Case (i).} No $\Gamma_i$ is reduced to one vertex. 

Then $G_{\Gamma_1},\dots,G_{\Gamma_k}$ are exactly the directly indecomposable direct factors of $G_\Gamma$, and they have trivial center \cite[Theorem~3.34]{Gre}. If some $\Lambda_j$ is reduced to one vertex, it follows from \cite[Proposition~2]{CdlH} that $H_{\Lambda_j}$ is isomorphic to a direct product of factors $G_{\Gamma_i}$, contradicting that the union of the corresponding subgraphs $\Gamma_i$ is not in $\mathbb{T}$. So no $\Lambda_j$ is reduced to one vertex, and hence $k=m$ and up to permuting the factors, one has $G_{\Gamma_i}\simeq H_{\Lambda_i}$ for all $i\in \{1,\dots,k\}$.
\medskip

By symmetry, the argument in Case~3.1 also treats the case where no $\Lambda_i$ is reduced to one vertex.
\medskip

\textbf{$\drsh$ Case (ii).} Some $\Gamma_i$ (say $\Gamma_1$) is reduced to one vertex, and some $\Lambda_j$ (say $\Lambda_1$) is reduced to one vertex. 

Since $\mathbb{T}$ contains all cliques,  $\Gamma_1$ and $\Lambda_1$ are the
only join factors reduced to a vertex. In this case no $G_{\Gamma_i}$ with $i\ge
2$ can be isomorphic to a direct factor of $H_{\Lambda_1}$: indeed otherwise,
$H_{\Lambda_1}$ would split as a graph product over a finite simple graph, one
of whose join components is $\Gamma_i$, contradicting that
$\Gamma_i\notin\mathbb{T}$ and that the family $\mathbb{T}$ is stable under
passing to join components.
It thus follows from \cite[Lemma~3.12]{Gen} (using \cite[Lemma~3.5]{Gen} to prove that the third assumption is satisfied) that $G_{\Gamma_1}$ is isomorphic to $H_{\Lambda_1}$, and that $G_{\Gamma_2}\times\dots\times G_{\Gamma_k}$ is isomorphic to $H_{\Lambda_2}\times\dots\times H_{\Lambda_m}$. Since $\Gamma_1$ and $\Lambda_1$ are the only factors reduced to a vertex, we can apply Case~(i) to the latter two direct products, and get the desired conclusion.
\end{proof}
 
Now let $i\in \{1,\dots,k\}$. If $\Gamma_i$ or $\Lambda_i$ is the disconnected graph on two vertices, Fact~\ref{Fact:Disconnected-with-two-vertices} gives a graph isomorphism $\theta_{i}:\Gamma_i \rightarrow \Lambda_i$ such that $G_{u}\simeq H_{\theta_{i}(u)}$ for all $u\in V\Gamma_i$. If $\Gamma_i$ and $\Lambda_i$ both have at least three vertices, our induction hypothesis gives a graph isomorphism $\theta_{i}:\Gamma_i \rightarrow \Lambda_i$ as above. Applying Fact~\ref{Fact:graph-collapses} now completes the proof in Case~3.
\end{proof}



\section{Generators for the automorphism group}\label{sec:3.2}

We finally generalize Theorem~\ref{theo:main-aut} to graphs which are not necessarily almost strongly reduced, as follows.

\begin{theo}\label{theo:aut-complete}
Let $\Gamma$ be a finite simple graph which is not an edge. Let $G$ be a graph product over $\Gamma$ with vertex groups $(G_v)_{v\in V\Gamma}$. Let $\mathbb{T}$ be a family of connected finite simple graphs with at least two vertices, that contains all cliques (with at least two vertices) and such that if $\Delta\in\mathbb{T}$, then every join component of $\Delta$ which is not reduced to one vertex also belongs to $\mathbb{T}$. Assume that: 
\begin{enumerate}
    \item\label{item:hyp1_aut_complete} $\Gamma$ does not contain any proper collapsible subgraph in the family $\mathbb{T}$;
    \item\label{item:hyp2_aut_complete} if a vertex group $G_v$ splits non-trivially as a graph product over a finite simple graph $\Delta$, then $\Delta\in\mathbb{T}$. 
\end{enumerate}
Then every $\varphi\in\Aut(G)$ is a product of transvections, partial conjugations and local automorphisms.
\end{theo}

A difference between the statements of Theorems~\ref{theo:isomorphism-complete} and~\ref{theo:aut-complete} is that we now assume that all subgraphs in $\mathbb{T}$ are connected, which in  combination with Assumption~\ref{item:hyp2_aut_complete} ensures that no vertex group splits non-trivially as a free product. This is closely related to the need to work with almost strongly reduced graphs in Theorem~\ref{theo:main-aut}, as explained in Remark~\ref{rk:aut}.

Again, Theorem~\ref{theo:main-aut} from the introduction, in the clique-reduced case, corresponds to the particular case where $\mathbb{T}$ is the family consisting of all cliques. The almost strongly reduced case, already established in Theorem~\ref{theo:conclusion}, corresponds to the case where $\mathbb{T}$ is the family of all connected finite simple graphs (not reduced to one vertex). The proof of Theorem~\ref{theo:aut-complete} will consist in reducing to the almost strongly reduced case.

\begin{proof} Let $\varphi\in \Aut(G)$. We will argue by induction on $|V\Gamma|$.

\begin{center}
    \textbf{Part 1: Preliminary remarks on collapsible subgraphs}
\end{center}

Since all subgraphs in $\mathbb{T}$ are connected, the two stability properties of Theorem~\ref{theo:isomorphism-complete} are satisfied. This will allow us to freely use the aforementioned theorem in the sequel, as well as the following three assertions from the previous section, which remain true in our setting.
\begin{enumerate}
    \item \label{item:remarks-on-collapsible-aut1} If $\Delta\subset \Gamma$ is a collapsible subgraph, then $\Delta$ is not an edge (Fact~\ref{Fact:ecrasable-pas-arete}).
    \item \label{item:remarks-on-collapsible-aut-un-sommet} Let $\Delta,\Upsilon \subset \Gamma$ be collapsible subgraphs such that $G_\Delta \simeq G_\Upsilon$. Then $\Delta$ is reduced to one vertex if and only if $\Upsilon$ is (Fact~\ref{Fact:Cas-un-sommet} applied to $\Lambda=\Gamma$).
    \item \label{item:remarks-on-collapsible-aut2} Let $\Delta,\Upsilon \subset \Gamma$ be collapsible subgraphs. Assume that $G_\Delta \simeq G_\Upsilon$ and that $\Delta$ or $\Upsilon$ is the disconnected graph on two vertices. Then there exists a graph isomorphism $\theta_\Delta:\Delta\rightarrow \Upsilon$ such that $G_{v}\simeq G_{\theta_\Delta(v)}$ for all $v\in V\Delta$
    (Fact~\ref{Fact:Disconnected-with-two-vertices} applied to $\Lambda=\Gamma$).
\end{enumerate}
In particular, if $\Delta,\Upsilon\subset\Gamma$ are collapsible subgraphs such that $G_\Delta\simeq G_\Upsilon$, and if $\Delta$ has at least three vertices, then $\Upsilon$ has at least three vertices.

\begin{center}
    \textbf{Part 2: Induction}
\end{center}

If $\Gamma$ is not a join, let $\overline{\Gamma}$ be the graph obtained from $\Gamma$ by collapsing all maximal proper collapsible connected subgraphs. Note that $\overline{\Gamma}$ is almost strongly reduced and not reduced to a vertex nor an edge. If $\Gamma$ is a join, we define $\overline{\Gamma}$ to be the graph obtained from $\Gamma$ by collapsing all join components. In both cases let $\pi_\Gamma:\Gamma\to\overline{\Gamma}$ be the associated graph collapse. 

\begin{Claim}\label{Claim:Maline}
The automorphism $\varphi$ is a product of  $\overline{\Gamma}$-transvections, $\overline{\Gamma}$-partial conjugations and $\overline{\Gamma}$-local automorphisms.
\end{Claim}

\begin{proof}[Proof of the claim]
We first assume that $\Gamma$ is not a join. Recall that $\overline{\Gamma}$ is almost strongly reduced and not an edge. Since by assumption all subgraphs in $\mathbb{T}$ are connected, Assumption~\ref{item:hyp2_aut_complete} ensures that no $\Gamma$-vertex group $G_v$ splits non-trivially as a free product; in other words, $G_v$ is either freely indecomposable or isomorphic to $\mathbb{Z}$. Since all subgraphs that are collapsed by $\pi_\Gamma$ are connected, and any graph product over a finite simple graph with at least two vertices is freely indecomposable, it follows that all vertex groups $G_{\bar v}$ with $\bar v\in V\overline{\Gamma}$ are freely indecomposable or isomorphic to $\mathbb{Z}$. The claim thus follows from Theorem~\ref{theo:conclusion}.

We now assume that $\Gamma$ splits non-trivially as a join, and denote by $\Gamma_1,\dots,\Gamma_k$ its join components.
\begin{itemize}
    \item If no $\Gamma_i$ is reduced to a vertex, then $G=G_{\Gamma_1}\times\dots\times G_{\Gamma_k}$ and the subgroups $G_{\Gamma_i}$ have trivial center and do not split non-trivially as direct products. Thus, by e.g.\ \cite[Lemma~3.5]{Fio}, there exists a bijection $\sigma$ of $\{1,\dots,k\}$ such that $\varphi(G_{\Gamma_i})=G_{\Gamma_{\sigma(i)}}$ for every $i\in\{1,\dots,k\}$. This exactly means that $\varphi$ is a $\overline{\Gamma}$-local automorphism.  
    \item Else, since by assumption $\Gamma$ does not contain any collapsible clique on at least two vertices, there is exactly one join component, say $\Gamma_1$, which is reduced to one vertex. Since $\Gamma_i$ is a proper collapsible subgraph of $\Gamma$, it does not belong to $\mathbb{T}$ (Condition~\ref{item:hyp1_aut_complete})
    and hence $G_{\Gamma_1}$ does not split as a graph product over $\Gamma_i$ (Condition~\ref{item:hyp2_aut_complete}). Moreover $\mathbb{T}$ being stable under taking join components, $G_{\Gamma_i}$ is not a direct factor of $G_{\Gamma_1}$. We can thus apply \cite[Lemma~3.12]{Gen} (using \cite[Lemma~3.5]{Gen} to prove that the third assumption is satisfied) to obtain that $\varphi(G_{\Gamma_1})=G_{\Gamma_1}$. Then Lemma~\ref{Lmm:detwister-un-facteur-direct} applied with $R_1=R_2=G_{\Gamma_1}$ and $P_1=P_2=G_{\Gamma_2}\times \cdots \times G_{\Gamma_k}$, gives a homomorphism ${\xi}:G_{\Gamma_2}\times \cdots \times G_{\Gamma_k}\rightarrow Z(G_{\Gamma_1})$ such that letting $\tau \in \Aut(G)$ be defined by $\tau(g)=g$ for all $g\in G_{\Gamma_1}$ and $\tau(g)=g{\xi}(g)$ for all $g\in G_{\Gamma_2}\times \dots \times G_{\Gamma_k} $ we have
    \begin{equation*}
        \varphi\circ\tau\left(G_{\Gamma_2}\times \dots \times G_{\Gamma_k} \right)=G_{\Gamma_2}\times \dots \times G_{\Gamma_k}.
    \end{equation*}
    Notice that $\tau$ is a product of $\overline{\Gamma}$-twists, and $\varphi\circ\tau$ restricts to an automorphism of $G_{\Gamma_2}\times \dots\times G_{\Gamma_k}$. By the previous point, this restriction is also a $\overline{\Gamma}$-local automorphism. The claim follows. \qedhere
\end{itemize}
\end{proof}

By Lemma~\ref{Lmm:Aut-and-graph-collapse} on page~\pageref{Lmm:Aut-and-graph-collapse}, any $\overline{\Gamma}$-fold (resp.\ $\overline{\Gamma}$-twist, resp.\ $\overline{\Gamma}$-partial conjugation) is a product of $\Gamma$-folds (resp.\ $\Gamma$-twists, resp.\ $\Gamma$-partial conjugations). It now remains to deal with the case of $\overline{\Gamma}$-local automorphisms.

So let $\varphi\in \Aut(G)$ be such that there exists $\sigma\in \Aut(\overline{\Gamma})$ such that $\varphi(G_{\bar{v}})=G_{\sigma(\bar v)}$ for all $\bar{v}\in V\overline{\Gamma}$. For all
$\bar{v}\in V\overline{\Gamma}$ let $\Delta_{\bar{v}}:=\pi^{-1}_{\Gamma}(\bar{v})$.  

We will now construct a $\Gamma$-local automorphism $\psi:G\to G$ such that $\psi(G_{\bar v})=G_{\sigma^{-1}(\bar v)}$. To this end, we first define its restrictions $\psi_{\bar v}$ to the subgroups $G_{\sigma(\bar v)}$, and for this we can distinguish three cases:
\begin{enumerate}
    \item If $\Delta_{\bar{v}}$ has at least three vertices, then by Part~1 of this proof, so has $\Delta_{\sigma(\bar v)}$. Theorem~\ref{theo:isomorphism-complete} thus gives a graph isomorphism $\theta_{\bar{v}} : \Delta_{\bar{v}} \rightarrow \Delta_{\sigma(\bar{v})}$ such that $G_{v}\simeq G_{\theta_{\bar{v}} (v)}$ for all $v\in V\Delta_{\bar{v}}$.
   For all $v\in V\Delta$ choose an isomorphism $\psi_v:G_{\theta_{\bar{v}}(v)} \rightarrow G_v$ and let 
   \begin{equation*}
        \psi_{\bar{v}}:
        \begin{cases}
        G_{\sigma(\bar{v})}&\rightarrow G_{\bar{v}},\\
        g&\mapsto   \psi_v(g)\ \text{if}\ g\in G_{\theta_{\bar{v}}(v)} \ \text{for some}\ v\in V\Delta_{\bar{v}}.
        \end{cases}
    \end{equation*}
    \item If $\Delta_{\bar{v}}$ is reduced to a vertex, then by item~\ref{item:remarks-on-collapsible-aut-un-sommet} from Part~1, so is $\Delta_{\sigma(\bar{v})}$. We denote them respectively by $v$ and $u$ and we choose for $\psi_{\bar v}$ an isomorphism from $G_u$ to $G_v$
    \item If $\Delta_{\bar{v}}$ is the disconnected graph on two vertices, then (item~\ref{item:remarks-on-collapsible-aut2} from Part~1) there exists a graph isomorphism $\theta_{\bar{v}}:\Delta_{\bar{v}}\rightarrow \Delta_{\sigma(\bar{v})}$ such that $G_{v}\simeq G_{\theta_{\bar{v}}(v)}$ for all $v\in V\Delta_{\bar{v}}$. We can thus define $\psi_{\bar{v}}:G_{\sigma(\bar{v})}\rightarrow G_{\bar{v}}$ as in the first case. 
\end{enumerate}
Note that the map $\theta:\Gamma\rightarrow \Gamma$ defined by $\theta_{|\Delta_{\bar{v}}}= \theta_{\bar{v}}$ is a graph automorphism of $\Gamma$. Then, letting $\psi:G\rightarrow G$ be defined by $\psi(g)=\psi_{\bar{v}}(g)$ for all $g\in G_{\sigma(\bar{v})}$ gives a 
$\Gamma$-local automorphism of $G$ such that $\psi(G_v)=G_{\theta^{-1}(v)}$ for all $v\in V\Gamma$. Moreover $(\varphi\circ\psi)_{|G_{\bar{v}}}$ is an automorphism of $G_{\Delta_{\bar{v}}}$ for all $\bar{v}\in V\overline{\Gamma}$. By our induction hypothesis $(\varphi\circ \psi)_{|G_{\bar{v}}}$ is a product of $\Delta_{\bar{v}}$-transvections,  $\Delta_{\bar{v}}$-partial conjugations and  $\Delta_{\bar{v}}$-local automorphisms. By Lemma~\ref{Lmm:Extending-aut-defined-on-collapsible}, extending $(\varphi\circ \psi)_{|G_{\bar{v}}}$ by the identity on $G_{\Gamma\backslash \Delta_{\bar v}}$ gives an automorphism $\alpha_{\bar v}$ of $G$ which is a product of $\Gamma$-transvections, $\Gamma$-partial conjugations and $\Gamma$-local automorphisms. Then $\varphi\circ \psi$ is equal to the product of all (pairwise commuting) automorphisms $\alpha_{\bar v}$ with $\bar v$ varying over $V\overline{\Gamma}$. This proves that $\varphi$ is a product of $\Gamma$-transvections, $\Gamma$-partial conjugations and $\Gamma$-local automorphisms. 
\end{proof}

\small
\bibliographystyle{alpha}
\bibliography{autos}
\addcontentsline{toc}{part}{Bibliography}
\vfill

\vspace{0.5cm}
\begin{flushleft}
  \textcolor{DeepSkyBlue4}{Amandine Escalier}\\ 
  Université Lyon 1, EC Lyon, INSA Lyon, Université Jean Monnet, CNRS, ICJ, UMR 5208,\\ Villeurbanne, France\\
  \emph{e-mail:~}\texttt{amandine.escalier@math.cnrs.fr}
  \medskip
\end{flushleft}
  
\noindent\textcolor{black!70}{\small A.~Escalier acknowledges support from Lyon 1 Université through AAP Accueil EC 2025 and AAP Accueil EC 2026.}

\begin{flushleft}
\textcolor{DeepSkyBlue4}{Camille Horbez}\\ 
Universit\'e Paris-Saclay, CNRS,  Laboratoire de math\'ematiques d'Orsay, 91405, Orsay, France \\
\emph{e-mail:~}\texttt{camille.horbez@universite-paris-saclay.fr}
\medskip
\end{flushleft}

\noindent\textcolor{black!70}{\small C.~Horbez acknowledges support from the European Research Council, through Grant~101040507 Artin-Out-ME-OA. Views and opinions expressed are however those of the authors only and do not necessarily reflect those of the European Union or the European Research Council; neither the European Union nor the granting authority can be held responsible for them.}
\end{document}
